\documentclass[11pt]{article}
\usepackage[a4paper,margin=1in]{geometry}
\usepackage[T1]{fontenc}
\usepackage{lmodern}
\usepackage{microtype}
\usepackage{amsmath,amssymb,amsthm,mathtools}
\usepackage{bm}
\usepackage{graphicx}
\usepackage{booktabs}
\usepackage{enumitem}
\usepackage{float}
\usepackage{tabularx}
\usepackage{booktabs}
\usepackage{makecell}
\usepackage{array}
\usepackage{hyperref}
\usepackage[nameinlink,noabbrev]{cleveref}

\hypersetup{
  colorlinks=true,
  linkcolor=blue,
  citecolor=blue,
  urlcolor=blue
}

\title{A Full Conservative Derivation of the Two-Body $3$PN Sector\\from the Unified $4$D Toy Model}
\author{Trevor Norris}
\date{\today}

\begin{document}
\maketitle

\begin{abstract}
We present the full conservative two-body $3$PN assembly for the $4$D toy-model program within the same declared closure hierarchy used at $1$PN and $2$PN. Starting from the exact $4$D reduction backbone, the frozen Newtonian/$1$PN/$2$PN ledger, and the $2.5$PN narrowing to the grouped real $P_{20}\oplus P_{21}\oplus P_{22}$ quadrupole bundle, we close the complete conservative near-zone ledger through order $c^{-6}$ in the fixed ADM chart. The strict Schwarzschild one-body gate fixes a minimal three-parameter repair ledger
\[
(\mu_{\rho,3},d_3,s_{24})=\left(\frac14,-\frac{45}{4},-\frac1{16}\right).
\]
The comparable-mass sector then splits exactly into three lanes: the ordinary-Lagrangian counterimage of the universal free two-body relativistic Hamiltonian, an exact grouped real $P_2$ middle-block closure, and a unique geometry-side static completion. Equivalently,
\[
\Delta \hat L_3^{\rm GR}
=
\Delta l_1 v^8
+
L_{P_2}^{\rm mid}
+
\Delta l_{15}^{(g)}U^4,
\]
with $\Delta H_3=-\Delta L_3$ in the fixed ADM chart. The apparent pure-kinetic remainder collapses to the compiler image of the free $3$PN Hamiltonian, while the remaining static slot is uniquely geometry-side after sigma-collapse. Thus the conservative $3$PN bottleneck is no longer algebraic: within the declared hierarchy, the two-body $3$PN ledger is fully assigned, and the live unresolved problem is the deeper moving-throat derivation of the grouped-$P_2$ and geometry constitutive lanes together with the already-narrow outgoing quadrupole-normalization gap inherited from the $2.5$PN program.
\end{abstract}

% Section content for 4d_3pn.tex

\section{Introduction and scope}
\label{sec:intro}

\subsection{Motivation: from full conservative \texorpdfstring{$2$}{2}PN and the \texorpdfstring{$2.5$}{2.5}PN narrowing to the first conservative grouped-\texorpdfstring{$P_2$}{P 2} gate}
\label{sec:intro-motivation}

The earlier papers in this program already established a strong conservative
backbone.  On the reduction side, the unified $4$D framework supplies the exact
projection identities, the projected continuity law with leakage, the controlled
Poisson hook, and the worldtube/coherent-defect reduction used to define the
brane-facing point-particle sector.  On the post-Newtonian side, the carried
lower-order ledger fixes the Newtonian and conservative $1$PN and $2$PN blocks
within a declared closure hierarchy, with the surviving viability conditions
\[
\kappa_\rho = 1,
\qquad
n = 5,
\qquad
\kappa_{\rm add} = \frac{1}{2},
\qquad
\kappa_{\rm PV} = \frac{3}{2},
\qquad
\beta_{\rm 1PN} = 3.
\]
Just as important, the conservative $2$PN assembly sharpened the ontology of the
particle itself.  The defect can no longer be treated as a pure scalar monopole
with small corrections; the solved conservative cross sector already looks like a
small throat-response theory with a carried odd dipole wake, a genuine
$P_0\oplus P_2$ mouth/support layer, and a separate geometry-closure channel.

The $2.5$PN program then made the higher-order bridge much sharper.  Its main
structural lesson was that the surviving universal point-particle route is not a
vague ``do more of the whole theory'' direction, but the orbital/worldtube STF
quadrupole branch, packaged conservatively as the grouped real
\[
P_{20}\oplus P_{21}\oplus P_{22}
\]
bundle.  That narrowing changes the meaning of the next conservative step.
After $2$PN, the right question is no longer whether one can add more local
coefficients somehow.  The right question is whether the first honest
higher-order conservative load carried by the grouped real $P_2$ sector can be
closed exactly inside the same hierarchy that already worked at $1$PN and $2$PN.

That is the purpose of the present paper.  At $3$PN the near-zone two-body
ledger through order $c^{-6}$ must close all of the following at once:
free-particle octic kinematics $v^8$, one-body/self terms through $U^4$, the
full comparable-mass middle interaction block, the static quartic term at order
$G^4/r^4$, and the generic-frame uniqueness problem in the fixed ADM chart.
This is the first conservative grouped-$P_2$ gate of the program.  It is also
the first place where the $2.5$PN narrowing becomes operational rather than only
organizational, because the grouped real quadrupole bundle now has to carry a
full middle-block theorem burden instead of merely serving as a suggestive
future interface.

A first structural consequence appears already at the strict one-body level.
The exact isotropic Schwarzschild gate through $3$PN is not reproduced by a
single cubic extension of the carried $2$PN denominator-style self sector.  The
minimal one-body repair ledger is instead
\[
\mu_{\rho,3}=\frac14,
\qquad
 d_3=-\frac{45}{4},
\qquad
 s_{24}=-\frac1{16},
\]
so the $3$PN gate already proves that this order is not a one-parameter
continuation of the $2$PN one-body closure.  But the one-body sector is still
only the front door.  The deeper algebraic issue is whether the full
comparable-mass $c^{-6}$ residual can be assigned uniquely once the lower-order
ledger is frozen, the exact COM target is imported, and the generic-frame chart
is fixed.

The main positive result of the paper is that it can.  In the fixed ADM chart,
the conservative $3$PN comparable-mass residual splits exactly as
\[
\begin{aligned}
\Delta \hat L_3^{\rm GR}
=\;&\underbrace{\Delta l_1 v^8}_{\text{compiler image of the free COM Hamiltonian}}
+\underbrace{L_{P_2}^{\rm mid}}_{\text{exact grouped real }P_2\text{ closure}}
\\
&+\underbrace{\Delta l_{15}^{(g)}U^4}_{\text{unique geometry completion}}.
\end{aligned}
\]
The middle interaction block is carried by a richer grouped real $P_2$ family,
the remaining static slot is uniquely geometry-side after sigma-collapse, and
the apparent pure-kinetic leftover is only the ordinary-Lagrangian counterimage
of the universal free relativistic two-body COM Hamiltonian.  A second crucial
structural lesson is that this statement has to be reached Hamiltonian-first.
Naive ordinary-Lagrangian COM reduction of the imported generic-frame target does
not reproduce the exact COM ordinary target, whereas the exact cubic Legendre
lift followed by COM reduction does.

The scope of the paper is therefore deliberately strong but intermediate.  It is
stronger than a notebook-level target match because the theorem chain is already
enough to show that the conservative $3$PN bottleneck is no longer an algebraic
coefficient search.  It is weaker than a fully first-principles moving-throat
PDE theorem because the grouped real $P_2$ constitutive law, the geometry-side
completion, and the passive/outgoing quadrupole normalization on the true branch
still require a deeper dynamical derivation.  The reader should therefore view
this paper in the same way as the earlier full conservative $1$PN and $2$PN
assemblies: a full assembly theorem \emph{within a declared closure hierarchy},
not an assumption-free theorem of the fully solved bulk problem.

A second purpose of the paper is bookkeeping clarity.  At $3$PN the argument
mixes carried exact identities, reduced-regime small-body and COM statements,
protocol-dependent constitutive closure, and final theorem consequences.  The
reader should be able to see, line by line, which ingredients are inherited
exactly from the lower-order stack, which depend on the declared reduction and
chart choices, which are fixed only within the response hierarchy, and which are
final assembly statements.  That separation is part of the result.

\subsection{What this paper claims}
\label{sec:intro-claims}

The main claims of the paper are the following.
\begin{enumerate}
  \item The exact $4$D parent reduction backbone from the earlier papers can be
  carried forward unchanged as the reduction framework of the present $3$PN
  analysis.

  \item The strict isotropic Schwarzschild one-body gate fixes a minimal
  three-parameter $3$PN repair ledger,
  \[
  \mu_{\rho,3}=\frac14,
  \qquad
  d_3=-\frac{45}{4},
  \qquad
  s_{24}=-\frac1{16},
  \]
  so the $3$PN one-body sector is not a one-parameter continuation of the
  carried $2$PN closure.

  \item The exact isotropic COM target can be imported, solved, and reduced to a
  finite $15$-slot comparable-mass residual vector, and the exact cubic
  ordinary/Hamiltonian compiler gives the authoritative bridge between the COM
  ordinary and Hamiltonian charts.

  \item In the fixed ADM chart the exact generic-frame cubic compiler reduces to
  \[
  \Delta H_3=-\Delta L_3,
  \]
  so the generic-frame ordinary representative is unique once the Hamiltonian
  target is imposed.  In particular, the correct generic-frame lift is
  Hamiltonian-first rather than naive ordinary-Lagrangian COM reduction.

  \item The minimal local grouped-$P_2$ front end is too small, but the richer
  grouped real $P_2$ family
  \[
  (T_{20},T_{21},T_{22},S_{20},S_{21},S_{22},V_{20},V_{21},V_{22})
  \]
  closes the full $9$-slot $3$PN middle block exactly.

  \item The remaining static slot is uniquely geometry-side after the exact
  sigma-collapse removes the apparent $P_0$/geometry ambiguity.

  \item The apparent pure-kinetic leftover is not a new throat datum.  It is
  only the ordinary-Lagrangian counterimage of the universal free relativistic
  two-body COM Hamiltonian.

  \item Therefore, within the same declared hierarchy used successfully at $1$PN
  and $2$PN, the conservative $3$PN comparable-mass ledger is fully assigned in
  the fixed ADM chart, with exact split
  \[
  \Delta \hat L_3^{\rm GR}
  =
  \Delta l_1 v^8
  +
  L_{P_2}^{\rm mid}
  +
  \Delta l_{15}^{(g)}U^4.
  \]
  What remains open is no longer algebraic coefficient matching.  What remains
  is the deeper moving-throat derivation of the grouped-$P_2$ constitutive law,
  the geometry completion, and then their insertion into the already narrowed
  passive/outgoing quadrupole bridge.
\end{enumerate}

Stated differently, the point of the paper is not that the model can be tuned
until a $3$PN-looking ledger appears.  The point is that, once the lower-order
stack and the fixed-chart target are handled correctly, the full $c^{-6}$
problem narrows to a short list of named theorem gates.  By the end of the
present paper those gates are passed, and the remaining unresolved questions are
constitutive and dynamical rather than algebraic.

\subsection{What this paper does not claim}
\label{sec:intro-nonclaims}

The negative side of the scope is equally important.

First, this paper does \emph{not} claim an assumption-free theorem that begins
with a fully solved moving-throat PDE and ends automatically with the complete
point-particle $3$PN dynamics.  As at $1$PN and $2$PN, the particle limit used
here is a controlled reduction inside a declared hierarchy.  The main text
therefore presents a reduced-theory theorem chain, not a finished theorem of the
full moving-throat bulk problem.

Second, this paper does \emph{not} claim that every grouped-$P_2$ or geometry
coefficient has already been derived microscopically from the completed throat
PDE.  The theorem proved here is that the conservative ledger is fully assigned
once the declared hierarchy is accepted.  The constitutive origin of the richer
grouped-$P_2$ family and the unique geometry completion remains a deeper
follow-on problem.

Third, this paper does \emph{not} claim any dissipative, radiative, or final
$2.5$PN outgoing-normalization result.  The point of the present paper is to
close the conservative $3$PN front end that the $2.5$PN narrowing had already
identified as the relevant grouped real quadrupole lane.  The passive/outgoing
quadrupole normalization on the actual moving-throat branch remains outside the
scope of the present theorem.

Fourth, the paper is restricted to the nonspinning conservative two-body sector
through order $c^{-6}$.  It does not attempt a treatment of spin couplings,
radiation reaction, strong-field nonperturbative completion, generic many-body
response, or higher-PN sectors beyond the present conservative $3$PN scope.
Nor does it claim to replace the broader script-side moving-throat or grouped
response record by a single closed PDE derivation.  The main text is
intentionally organized around the smallest theorem chain needed for the present
$3$PN statement.

These non-claims are not qualifications to hide.  They are part of the correct
interpretation of the result.  The paper should be read as establishing the
sharpest honest conservative $3$PN statement currently supported by the program:
a full algebraic assembly theorem within the declared hierarchy, with the
remaining open work shifted from coefficient matching to microscopic response
and normalization.

\subsection{Claim taxonomy}
\label{sec:intro-taxonomy}

To keep the argument referee-safe, we retain the same bookkeeping discipline used
in the earlier full conservative papers.
\begin{description}
  \item[Exact.] Identities that follow directly from the declared lower-order
  ledger, the exact one-body target, the exact compiler algebra, symmetry
  bookkeeping, or exact source-map manipulations once the chart and basis
  choices have been frozen.

  \item[Controlled reduction.] Steps that require an explicit reduced
  description, such as the carried $4$D-to-brane and worldtube reductions, the
  reduced isotropic COM basis, the grouped real $P_2$ packaging, or the
  Hamiltonian-first COM reduction rule used to recover the exact fixed-chart
  representative.

  \item[Protocol closure.] Statements that depend on the same declared response
  hierarchy used below $3$PN rather than on a fully solved moving-throat PDE.
  In the present paper the most important examples are the one-body repair
  packaging, the richer local grouped-$P_2$ constitutive family, and the
  interpretation of the remaining static slot as a geometry-side completion.

  \item[Full assembly.] The final theorem consequence after the previous
  ingredients are inserted.  In the present paper this means that, once the
  lower-order ledger, the exact $3$PN targets, the fixed ADM chart, and the
  declared closure hierarchy are frozen, the final equality is exact algebra.
\end{description}

This taxonomy is not cosmetic.  It prevents four logically distinct moves from
being blurred together: carrying forward exact lower-order structure, using the
declared reductions and chart choices, invoking the constitutive closure needed
to close the grouped-$P_2$ and geometry lanes, and stating the final theorem
consequence that survives after those ingredients are frozen.

\subsection{Roadmap}
\label{sec:intro-roadmap}

The next section gives an executive summary of the final theorem ledger.
Section~\ref{sec:dictionary} then records the minimal carry-forward dictionary
from the earlier $4$D, $1$PN, $2$PN, and $2.5$PN papers and states the exact
$3$PN target.  Section~\ref{sec:onebody} begins the proof chain with the strict
one-body Schwarzschild gate, the minimal repair ledger, the natural self/static
seed, and the exact cubic compiler.

Section~\ref{sec:compmass} lifts that one-body input into the full
comparable-mass problem: the reduced COM basis, the imported exact COM target,
the generic-frame scaffold, the COM-null and contact/gauge structure, the
failure of naive ordinary-Lagrangian COM reduction, the Hamiltonian-first lift,
and the fixed-chart uniqueness theorem.  Section~\ref{sec:grouped-p2} then
isolates the first conservative grouped-$P_2$ gate, explains why the minimal
front end fails, and proves that the richer grouped family closes the full
middle block exactly.

Section~\ref{sec:geometry} shows that the remaining static slot is uniquely
geometry-side after sigma-collapse, while Section~\ref{sec:kinetic} proves that
the apparent pure-kinetic leftover is only a compiler image of the free two-body
Hamiltonian rather than a new throat datum.  Section~\ref{sec:final-theorem}
then states the final exact split in COM and generic-frame form and records the
equality to the standard $3$PN ADM target within the declared hierarchy.

The last main-text blocks serve two carry-forward purposes.  Section~\ref{sec:interface-25pn}
explains exactly how the present conservative theorem interfaces with the
already-narrowed $2.5$PN quadrupole program, why $3$PN is the first
conservative grouped-$P_2$ gate, what the isotropy data mean, and what the full
moving-throat PDE still has to supply.  Section~\ref{sec:fixed-open} separates
what is now fixed from what remains symbolic or open, Section~\ref{sec:verification}
records the role and scope of the referee suite, and Section~\ref{sec:conclusions}
closes the paper.  The appendices collect the detailed algebraic ledgers,
compiler derivations, grouped-$P_2$ constructions, geometry completion, and
verification scaffolding that would obscure the theorem chain if left in the
body of the paper.

% Section content for 4d_3pn.tex

\section{Executive summary of results}
\label{sec:executive}

This paper closes the first honest conservative grouped-$P_2$ gate of the
program.  The earlier papers already froze the exact $4$D reduction backbone,
carried the conservative Newtonian+$1$PN+$2$PN ledger through its declared
closure hierarchy, and narrowed the live higher-order universal point-particle
route to the orbital/worldtube STF quadrupole branch packaged as the grouped real
\[
P_{20}\oplus P_{21}\oplus P_{22}.
\]
The present question is therefore not whether one can add more local coefficients
somehow.  It is whether the full nonspinning conservative two-body ledger at
order $c^{-6}$ can be assigned uniquely once the strict one-body gate, the exact
reduced COM target, and the fixed ADM chart are handled correctly.  The answer is
now sharp: within the same declared closure hierarchy used successfully at $1$PN
and $2$PN, the conservative $3$PN two-body ledger is fully assigned in the fixed
ADM chart.

\paragraph{Headline result.}
Within the stated hierarchy, the exact comparable-mass $3$PN residual splits as
\[
\begin{aligned}
\Delta \hat L_3^{\rm GR}
=\;&\underbrace{\Delta l_1 v^8}_{\text{compiler image of the free COM Hamiltonian}}
+\underbrace{L_{P_2}^{\rm mid}}_{\text{exact grouped real }P_2\text{ closure}}
\\
&+\underbrace{\Delta l_{15}^{(g)}U^4}_{\text{unique geometry completion}}.
\end{aligned}
\]
The two scalar coefficients worth keeping visible already at the executive level
are
\[
\Delta l_1=\frac{3\nu(3\nu-1)(4\nu-1)}{16},
\qquad
\Delta l_{15}^{(g)}=\frac{\nu\bigl(408\nu^2+1232\nu-2080+63\pi^2\bigr)}{96}.
\]
The middle interaction block is carried entirely by a richer grouped real $P_2$
family, the remaining static scalar slot is uniquely geometry-side after
sigma-collapse, and the apparent pure-kinetic leftover is only the
ordinary-Lagrangian counterimage of the universal free relativistic two-body COM
Hamiltonian.  What remains open is therefore no longer algebraic coefficient
matching.  What remains open is the deeper moving-throat derivation of the
grouped-$P_2$ constitutive law and geometry completion, followed by their
insertion into the already narrowed passive/outgoing quadrupole bridge.

Table~\ref{tab:executive-ledger-3pn} lists the load-bearing results established
in the present paper and the role each plays in the theorem chain.

\begin{table}[t]
  \centering
  \caption{Load-bearing results of the present $3$PN paper.}
  \label{tab:executive-ledger-3pn}
  \footnotesize
  \setlength{\tabcolsep}{4pt}
  \renewcommand{\arraystretch}{1.12}
  \begin{tabular}{p{0.22\linewidth}p{0.58\linewidth}p{0.14\linewidth}}
    \hline
    Item & Result & Status \\
    \hline
    One-body gate
    & The strict isotropic Schwarzschild target through $c^{-6}$ is fixed exactly,
      and the minimal repair ledger is
      $\mu_{\rho,3}=1/4$, $d_3=-45/4$, $s_{24}=-1/16$.
    & Settled \\

    Fixed-chart compiler
    & The imported exact GR COM target becomes a finite $15$-slot residual, the
      Hamiltonian-first route is the authoritative lift, and in the fixed ADM
      chart the generic-frame compiler reduces to
      $\Delta H_3=-\Delta L_3$.
    & Settled \\

    Grouped-$P_2$ middle block
    & The minimal local grouped-$P_2$ front end fails, but the richer family
      $(T_{20},T_{21},T_{22},S_{20},S_{21},S_{22},V_{20},V_{21},V_{22})$
      closes the full $9$-slot middle block exactly with
      $\det M_{\rm mid}=-4/27$.
    & Exact closure \\

    Static scalar lane
    & The grouped module predicts a definite static companion, and the remaining
      scalar gap collapses algebraically to one unique geometry-side
      representative after the exact sigma identity
      $\nu(p^3+q^3)=(1-3\nu)(p^2q+pq^2)$.
    & Unique completion \\

    Pure-kinetic slot
    & The apparent leftover COM $v^8$ term is not a new throat datum; it is the
      exact ordinary-Lagrangian compensation required to reproduce the universal
      free relativistic two-body COM Hamiltonian.
    & Resolved \\

    Remaining gap
    & The moving-throat PDE must still supply the microscopic origin of the
      grouped-$P_2$ data, the isotropy verdict $a_2=b_2=0$ on the true branch,
      and the final passive/outgoing quadrupole normalization.
    & Open \\
    \hline
  \end{tabular}
\end{table}

\paragraph{How the ledger closes.}
The proof chain has three main positive gates and one decisive bookkeeping
lesson.  First, the strict one-body Schwarzschild target shows that $3$PN is not
a one-parameter continuation of the carried $2$PN self sector; the $U^3v^2$ and
$U^2v^4$ channels together force the three-parameter repair ledger
\((\mu_{\rho,3},d_3,s_{24})\).  Second, the exact reduced COM target and the
exact cubic Legendre compiler turn the comparable-mass problem into a finite
algebraic residual rather than a blind coefficient search.  Third, the grouped
real $P_2$ ontology survives its first full conservative stress test: the minimal
constitutive front end is too rigid, but the richer grouped family closes the
entire nine-slot middle block exactly.  The bookkeeping lesson is that the
correct quotient principle is Hamiltonian-first.  Naive ordinary-Lagrangian COM
reduction of the imported generic-frame target does not reproduce the exact COM
ordinary target, whereas the exact cubic Legendre lift followed by COM reduction
does.  Once this is done, the fixed-chart generic-frame compiler collapses to the
slotwise sign flip
\[
\Delta H_3=-\Delta L_3,
\]
so no true algebraic ambiguity survives after the Hamiltonian target is imposed.

\paragraph{One-body gate and self/static seed.}
The strict one-body target is the exact isotropic Schwarzschild expansion through
$3$PN,
\[
\frac{L_{3\mathrm{PN},1\mathrm{body}}}{m}
=
\frac5{128}v^8
+\frac{11}{16}Uv^6
+\frac{47}{16}U^2v^4
+\frac{13}{8}U^3v^2
-\frac18U^4.
\]
A single cubic extension of the carried $2$PN denominator-style self sector is
not enough to reproduce this answer.  The exact repair ledger is instead
\[
\mu_{\rho,3}=\frac14,
\qquad
 d_3=-\frac{45}{4},
\qquad
 s_{24}=-\frac1{16},
\]
where $\mu_{\rho,3}$ is the cubic static gate, $d_3$ is the cubic denominator
datum, and $s_{24}$ is the genuinely new $U^2v^4$ self slot.  Symmetric
promotion of the exact one-body target then gives the natural $3$PN self/static
seed,
\[
L_{3,\mathrm{seed}}
=
\frac{5}{128}(m_Av_A^8+m_Bv_B^8)
+\frac{11Gm_A m_B}{16r}(v_A^6+v_B^6)
+\cdots
-\frac{G^4m_A m_B}{8r^4}(m_B^3+m_A^3),
\]
which isolates all universal one-body/self/static content before the genuine
comparable-mass lift begins.

\paragraph{Fixed-chart algebra and unique generic-frame representative.}
In the standard isotropic COM basis the exact cubic ordinary/Hamiltonian compiler
is diagonal.  Importing the exact GR $3$PN COM target and subtracting the frozen
self/static seed therefore yields an explicit $15$-slot residual vector.  At the
generic-frame level the interaction scaffold has block sizes
\[
24,\qquad 17,\qquad 8,\qquad 1,
\]
for the $G/r$, $G^2/r^2$, $G^3/r^3$, and $G^4/r^4$ blocks.  COM projection of
that scaffold reveals a small middle-block obstruction if one keeps only the
naive seed placement, and it also shows that COM-blind directions are larger than
the actual ordinary contact/gauge orbit.  None of this is a failure of the
scaffold itself.  The decisive point is that once the Hamiltonian target is
imposed in the fixed ADM chart, the generic-frame residual compiler becomes
\[
\Delta H_3=-\Delta L_3,
\]
so the generic-frame ordinary representative is unique.  At that stage the
bottleneck is no longer algebraic span or chart ambiguity.

\paragraph{The grouped real \texorpdfstring{$P_2$}{P 2} middle block.}
At $3$PN the conservative grouped data are the three $O(\omega^2)$ coefficients
\[
u_2^{(20)},\qquad u_2^{(21)},\qquad u_2^{(22)},
\]
or equivalently the grouped trace/anomaly variables
\[
\bar u_2=\frac{u_2^{(20)}+2u_2^{(21)}+2u_2^{(22)}}{5},
\qquad
 a_2=\frac{2u_2^{(20)}-u_2^{(21)}-u_2^{(22)}}{10},
\qquad
 b_2=\frac{u_2^{(21)}-u_2^{(22)}}{2}.
\]
The first isotropy test for the minimal branch is therefore
\[
a_2=0,
\qquad
b_2=0.
\]
A direct time-local grouped front end built only from the leading three
$O(\omega^2)$ couplings is too rigid: after the unique local demotion it lands in
the correct nine interaction slots but has rank three and obeys six exact
relations violated by the GR target.  The grouped ontology itself is therefore
not the problem; only the minimal constitutive choice is too small.  Enlarging
the local grouped family to
\[
(T_{20},T_{21},T_{22},S_{20},S_{21},S_{22},V_{20},V_{21},V_{22})
\]
gives an exact $9\times 9$ middle-block compiler with constant determinant
\[
\det M_{\rm mid}=-\frac4{27}.
\]
This richer grouped lift reproduces the full GR middle block
\((\Delta l_6,\dots,\Delta l_{14})\) with no residual left.

\paragraph{Geometry completion and pure-kinetic collapse.}
Once the grouped real $P_2$ block is closed, the only remaining conservative
scalar task is the one-dimensional static remainder.  The grouped module itself
predicts a definite static companion,
\[
\Delta l_{15}^{(P_2\,\rm pred)}
=
\frac{\nu\bigl(293-308\nu-102\nu^2\bigr)}{24},
\]
so the actual scalar gap is explicit before any geometry interpretation is chosen.
Allowing both body-local $P_0$-type and pair/geometry representatives initially
creates an apparent two-family ambiguity, but the exact mass identity
\[
\nu(p^3+q^3)=(1-3\nu)(p^2q+pq^2)
\]
collapses that ambiguity algebraically.  The resulting static counterterm is
therefore unique and geometry-side,
\[
\Delta L_{3,\rm ct}^{\rm scalar}
=
\frac{G^4pq}{96r^4}
\bigl(408\nu^2+1232\nu-2080+63\pi^2\bigr)(p^2q+pq^2),
\]
with Hamiltonian image $\Delta H_{3,\rm ct}^{\rm scalar}=-\Delta L_{3,\rm ct}^{\rm scalar}$.
The final apparent pure-kinetic leftover is equally non-mysterious.  At
generic-frame level the free $3$PN ordinary piece is simply
\[
L_{3,\rm free}^{\rm gen}=\frac{5}{128}(m_Av_A^8+m_Bv_B^8),
\]
while the exact free COM Hamiltonian yields
\[
h_1^{\rm free}=\frac{-5+35\nu-70\nu^2+35\nu^3}{128},
\qquad h_2=h_3=h_4=h_5=0.
\]
Applying the exact COM compiler then gives precisely the leftover
\[
\Delta l_1=\frac{3\nu(3\nu-1)(4\nu-1)}{16},
\]
so the pure-kinetic slot is only the ordinary-Lagrangian compensation required to
reproduce the universal free relativistic two-body COM Hamiltonian.

\paragraph{What the result fixes for the \texorpdfstring{$2.5$}{2.5}PN bridge.}
The main positive interface result is that the grouped real $P_2$ ontology
survives its first honest conservative stress test.  The richer local grouped
compiler closes the full nine-slot middle block exactly, so the throat side does
not need a qualitatively new non-$P_2$ conservative carrier for that sector.  The
remaining static slot is uniquely geometry-side, and the pure-kinetic slot is
only a compiler image of the free COM Hamiltonian, so neither competes with the
grouped quadrupole module for the middle-block burden.  The remaining $2.5$PN
gap is therefore much narrower than ``derive $3$PN somehow.''  The moving-throat
PDE still has to supply the actual low-frequency constitutive origin of the
grouped real $P_2$ coefficients, the isotropy verdict $a_2=b_2=0$ on the true
branch, and then the final passive/outgoing quadrupole normalization on that
branch.  The present theorem should therefore be read as removing the
conservative algebraic uncertainty beneath that outgoing-normalization gap rather
than as closing the gap itself.

% Section content for 4d_3pn.tex

\section{Carry-forward dictionary and theorem target}
\label{sec:dictionary}

This section is intentionally downstream of the parent $4$D action paper, the
full conservative $1$PN and $2$PN assembly papers, and the $2.5$PN theorem-audit
package.  We do not repeat those longer derivations here.  Instead, we record
only the imported objects, frozen lower-order values, chart conventions, and
exact target statements that the present $3$PN proof chain actually uses.  The
purpose of the section is fourfold: to freeze notation, to separate exact
imports from controlled reductions, to make clear what is being carried rather
than rederived, and to state the theorem target that organizes the sections that
follow.

The guiding principle is the same one used in the earlier conservative papers.
What is exact in the parent reduction framework should remain visibly exact;
what enters only after a declared near-zone, small-body, COM, or grouped-channel
assumption should remain visibly reduced; and what is fixed only within the
response hierarchy should remain visibly hierarchical rather than disguised as an
assumption-free theorem of the fully solved moving-throat PDE.  The present
paper is therefore built on an imported dictionary rather than on a fresh
rederivation of the full bulk problem.

\subsection{Minimal imported \texorpdfstring{$4$}{4}D reduction dictionary}
\label{sec:dict-4d}

The parent theory lives on a $(4+1)$-dimensional spacetime with coordinates
\[
X^M=(t,x,y,z,w),
\qquad
\mathbf{X}=(x,y,z,w)\in\mathbb R^4,
\qquad
\mathbf{x}=(x,y,z)\in\mathbb R^3.
\]
The fundamental bulk objects carried into the present paper are the matter field
$\psi(\mathbf{X},t)$, its density
\[
\rho(\mathbf{X},t)=|\psi(\mathbf{X},t)|^2,
\]
the optional localized gauge field $A_M(\mathbf{X},t)$ with localization profile
$Z(w)$, and the effective throat geometry variables
\[
a(t),\qquad L(t).
\]
When electric charge is named explicitly in the inherited $4$D notation, it is
always the corrected branch quantity
\[
q_* = \eta_Q e_*,
\qquad
q_{\rm eff}=\frac{q_*}{\sqrt{Z_{\rm int}}},
\qquad
Z_{\rm int}=\int_{-\infty}^{\infty} Z(w)\,dw,
\]
not the gravity-side scalar coefficient $\kappa_\rho$ that appears in the
post-Newtonian ledger below.

A brane observer does not see the full bulk fields directly.  Brane observables
are introduced by a fixed normalized projection kernel $W(w)$,
\[
\int_{-\infty}^{\infty} W(w)\,dw = 1,
\]
through the projection map
\[
\mathcal P_W[Q](t,\mathbf{x})
\equiv
\int_{-\infty}^{\infty} W(w)\,Q(t,\mathbf{x},w)\,dw.
\]
In particular,
\[
\rho_{\rm br}=\mathcal P_W[\rho],
\qquad
\mathbf J_{\rm br}=\mathcal P_W[(j^x,j^y,j^z)],
\qquad
\mathbf v_{\rm br}=\frac{\mathbf J_{\rm br}}{\rho_{\rm br}}
\quad (\rho_{\rm br}>0).
\]
This projection operation is distinct from \emph{reduction}.  Projection defines
what a brane observer measures; reduction refers to the further controlled
elimination of the extra coordinate under an explicit ansatz or scaling regime.

The exact projected continuity law carried forward from the parent theory is
\[
\partial_t\rho_{\rm br}+\nabla_3\!\cdot\mathbf J_{\rm br}=S_{\rm leak},
\]
with leakage source
\[
S_{\rm leak}
=
-\bigl[W(w)j^w\bigr]_{-\infty}^{+\infty}
+
\int_{-\infty}^{\infty} W'(w)j^w\,dw.
\]
After Helmholtz decomposition of the brane velocity,
\[
\mathbf v_{\rm br}=\nabla_3\phi_{\rm br}+\mathbf v_T,
\qquad
\nabla_3\!\cdot\mathbf v_T=0,
\]
one obtains the exact longitudinal identity
\[
\rho_{\rm br}\,\nabla_3^2\phi_{\rm br}
=
S_{\rm leak}
-\partial_t\rho_{\rm br}
-(\nabla_3\rho_{\rm br})\cdot(\nabla_3\phi_{\rm br}+\mathbf v_T).
\]
Under the carried quasi-static near-zone assumptions this reduces to the Poisson
hook
\[
\nabla_3^2\phi_{\rm br}\simeq \frac{S_{\rm eff}}{\rho_0},
\qquad
\Phi_N=\lambda_\Phi\,\phi_{\rm br},
\]
which is the origin of the Newtonian scalar sector used in the lower-order
particle reduction.

The second imported reduction step is the controlled worldtube/coherent-defect
limit.  For a compact defect of size $a\ll r$ in a smooth external field, the
center-of-mass reduction gives the usual point-particle force law
\[
M_A\,\ddot{\mathbf X}_A
=
-M_A\nabla\Phi_{\rm ext}(\mathbf X_A)
+O\!\left(\frac{a^2}{r^2}\,\frac{\Phi_{\rm ext}}{r}\right),
\]
so that the universal point-particle source at leading order is controlled by
the worldline/worldtube multipoles rather than by arbitrary internal details of
the defect.  That worldtube dictionary is what later allows the higher-order
universal bridge to be packaged in terms of the orbital/worldtube STF
quadrupole rather than a generic internal support excitation.

\subsection{Frozen conservative carry-forward ledger}
\label{sec:dict-carry}

The present paper imports the conservative Newtonian+$1$PN+$2$PN ledger as
frozen input, not as open fit parameters.  The surviving lower-order viability
conditions are
\[
\kappa_\rho = 1,
\qquad
n = 5,
\qquad
\kappa_{\rm add}=\frac12,
\qquad
\kappa_{\rm PV}=\frac32,
\qquad
\beta_{\rm 1PN}=3.
\]
The first four values are the scalar mass-dressing, optical/barotropic,
added-mass, and adiabatic breathing-response coefficients frozen by the earlier
program.  The fifth is the resulting carried $1$PN ledger value.  The present
paper is not allowed to repair a $3$PN mismatch by retuning any of these
numbers.

When referenced, the parity-even wake data inherited from the full conservative
$1$PN assembly are also treated as frozen branch data,
\[
\alpha^2=\frac34,
\qquad
a_H=0,
\qquad
K_{\rm vec}=\frac{2}{\pi^2},
\]
not as live higher-order fit parameters.

The conservative $2$PN assembly sharpened the ontology of the defect further.
By the start of the present $3$PN paper, the particle could no longer be treated
as a pure scalar monopole with small corrections.  The carried conservative
structure already looked like a small throat-response theory with three distinct
pieces:
\begin{enumerate}
  \item a carried odd dipole wake sector,
  \item a genuine even $P_0\oplus P_2$ mouth/support layer,
  \item and a separate geometry-closure channel.
\end{enumerate}
This matters conceptually because the present $3$PN work is \emph{not} starting
from a blank EFT.  It is extending a structure that was already forced by the
conservative $2$PN match.  In particular, the grouped real quadrupole bundle is
not being introduced ad hoc at $3$PN; it is the conservative reorganization of a
support/geometry structure that was already visible at lower order.

What remains genuinely open from the lower-order side is not the existence of
these channels but their full dynamical completion.  In particular, the dynamic
pole or compliance scales
\[
\Omega_{1\perp},\ \Omega_{10},\ \Omega_0,
\ \Omega_{20},\ \Omega_{21},\ \Omega_{22},\ \Omega_g
\]
are not claimed as already derived from the fully solved moving-throat PDE.  This
is one reason the present paper continues to speak in terms of a declared closure
hierarchy rather than an assumption-free first-principles theorem.

Two structural consequences of this imported ledger should be stated explicitly.
First, the present paper is not allowed to repair a higher-order mismatch by
changing lower-order conservative constants that were already fixed below $3$PN.
Second, the defect already carries enough internal support structure that the
$3$PN analysis must distinguish carefully between universal point-particle
content, self/static carry-forward content, grouped-$P_2$ middle-block content,
and the separate static geometry lane.

\subsection{Why the \texorpdfstring{$2.5$}{2.5}PN notes narrow the next conservative target to the grouped real \texorpdfstring{$P_2$}{P 2} bundle}
\label{sec:dict-25pn-bridge}

The $2.5$PN theorem audit is what makes the present $3$PN problem sharply
focused rather than diffuse.  By the end of that audit, the live universal
point-particle route had already narrowed to the orbital/worldtube STF
quadrupole branch.  In grouped real-channel language, that means the only
higher-order conservative payload that still matters for the minimal branch is
the grouped real
\[
P_{20}\oplus P_{21}\oplus P_{22}
\]
bundle.

More specifically, the $2.5$PN notes reduce the future conservative task to the
first even coefficients of that grouped real $P_2$ response.  On the minimal
isotropic branch, once the common conservative quadrupole data
\[
(c_0,c_2,c_4)
\qquad\text{or equivalently}\qquad
(\overline K_0,\overline K_2,\overline K_4)
\]
are known, the effective pole $\Omega_Q$ and the odd Burke--Thorne coefficient
$\overline\Gamma_5$ are fixed automatically.  The isotropy gate itself is also
explicit in grouped language: the grouped triples must collapse to one isotropic
trace, with the anisotropy defects vanishing.

That is why the present paper uses the grouped real $P_2$ bundle as its
conservative front end.  The question is no longer ``derive all of higher PN''
or ``add more coefficients somehow.''  The question is whether the conservative
$3$PN calculation lands on the same grouped-$P_2$ branch that the $2.5$PN notes
identified as the only surviving universal quadrupole route.

At $3$PN the grouped conservative data are precisely the three $O(\omega^2)$
coefficients
\[
u_2^{(20)},\qquad u_2^{(21)},\qquad u_2^{(22)},
\]
or, equivalently, the grouped trace/anomaly variables
\[
\bar u_2=\frac{u_2^{(20)}+2u_2^{(21)}+2u_2^{(22)}}{5},
\qquad
 a_2=\frac{2u_2^{(20)}-u_2^{(21)}-u_2^{(22)}}{10},
\qquad
 b_2=\frac{u_2^{(21)}-u_2^{(22)}}{2}.
\]
The first isotropy test for the minimal branch is therefore
\[
a_2=0,
\qquad
b_2=0.
\]
The present paper does not pretend to settle that isotropy verdict from a fully
solved moving-throat PDE.  What it does is close the conservative $3$PN middle
block in exactly the grouped language that the later bridge will need.

\subsection{PN bookkeeping and invariant variables}
\label{sec:dict-pn}

The conservative expansion is organized exactly as in the earlier papers:
\[
L_{\rm cons}=L_0+c^{-2}L_1+c^{-4}L_2+c^{-6}L_3+O(c^{-8}).
\]
So throughout this paper, ``$3$PN'' always means the coefficient of $c^{-6}$.

In the generic frame, the basic two-body invariants remain
\[
\mathbf r=\mathbf x_A-\mathbf x_B,
\qquad
r=|\mathbf r|,
\qquad
\mathbf n=\mathbf r/r,
\]
\[
v_A^2=\mathbf v_A^2,
\qquad
v_B^2=\mathbf v_B^2,
\qquad
v_{AB}=\mathbf v_A\!\cdot\!\mathbf v_B,
\]
\[
v_{An}=\mathbf v_A\!\cdot\!\mathbf n,
\qquad
v_{Bn}=\mathbf v_B\!\cdot\!\mathbf n.
\]
The usual local pair potentials are
\[
U_A=\frac{Gm_B}{r},
\qquad
U_B=\frac{Gm_A}{r},
\]
and in the strict one-body gate we write
\[
U=\frac{GM}{r}.
\]

In the COM sector, the reduced isotropic ordinary-Lagrangian basis uses the
velocity invariants $(v^2,d^2,u)$ with
\[
d\equiv \dot r,
\qquad
u\equiv\frac{GM}{rc^2},
\qquad
\nu\equiv\frac{m_A m_B}{(m_A+m_B)^2}.
\]
The matching COM Hamiltonian basis uses the corresponding reduced momentum
invariants $(p^2,p_r^2,u)$.  In the present paper the ordinary COM coefficients
are denoted $l_i$, the COM Hamiltonian coefficients $h_i$, and the residuals
beyond the frozen self/static seed by $\Delta l_i$ and $\Delta h_i$.
This placement matters because one of the main $3$PN lessons was that COM
bookkeeping is a compiler convenience, not a license to treat naive ordinary COM
substitution as a generic-frame theorem.  For later compact formulas it is also
convenient to write
\[
v_\perp^2\equiv v^2-d^2,
\]
but the invariant generic-frame quantities displayed above remain primary.

\subsection{COM versus generic-frame conventions}
\label{sec:dict-com-generic}

The fixed convention carried forward from the present session is simple:
use the COM basis whenever it simplifies the exact target or the diagonal
compiler map, but treat the \emph{generic-frame Hamiltonian compiler} as the
authoritative uniqueness statement.

The reason for this convention is now settled.  A naive COM reduction of the
imported generic-frame ordinary $3$PN target does \emph{not} reproduce the
correct COM ordinary target.  By contrast, the Hamiltonian-first route does:
\begin{enumerate}
  \item import the generic-frame ordinary target,
  \item apply the exact cubic Legendre lift first,
  \item then reduce the Hamiltonian to COM.
\end{enumerate}
That sequence recovers the standard $3$PN COM Hamiltonian exactly.  The next
sections therefore treat Hamiltonian-first lifting as the authoritative quotient
principle and do not reopen the earlier naive-COM-reduction detour.

A second settled convention follows from the fixed ADM chart result.  Once the
lower-order ledger is frozen, the full generic-frame $3$PN residual compiler is
exact and trivial in that chart,
\[
\Delta H_3=-\Delta L_3.
\]
So COM-blind families or ordinary contact freedoms that appear before the
Hamiltonian chart is fixed are no longer true ambiguities afterward.  In the
fixed chart, the generic-frame Hamiltonian target determines the ordinary
representative uniquely.

\subsection{Notation firewall}
\label{sec:dict-firewall}

The notation firewall from the corrected $4$D / EM / $2$PN stack must be kept
intact.

First, electric-charge notation is completely separate from the conservative
$3$PN coefficient ledger.  When charge appears at all, it uses the corrected
variables
\[
\eta_Q=\pm1,
\qquad
q_* = \eta_Q e_*,
\qquad
q_{\rm eff}=\frac{q_*}{\sqrt{Z_{\rm int}}},
\]
while the historical gravity-side bare $q=1$ remains renamed as
\[
\kappa_\rho=1.
\]
Quantized circulation belongs to the magnetic/vortical sector, not to the
electric-charge dictionary.

Second, several overloaded symbols from the older literature must not be allowed
to drift.  The appendix-side source-weight vector
\[
J=\left(\frac{4}{\sqrt5},\frac54,0,0,0,0\right)
\]
is not an electromagnetic current.  The pure-potential geometry-closure deficit
$\Delta_{\rm geom}$ and the monopole wall add-on $\Delta K_{00}$ are distinct
objects.  And in the present paper the grouped labels $20,21,22$ always denote
grouped real $P_2$ channels, not tensor indices.

Third, the grouped real $P_2$ variables are kept separate from the COM
coefficient ledgers.  The grouped conservative data are written as
\[
u_2^{(20)},\qquad u_2^{(21)},\qquad u_2^{(22)},
\]
or, after regrouping, as
\[
\bar u_2,\qquad a_2,\qquad b_2,
\]
and are never to be confused with the COM coefficients $l_i$ and $h_i$.
Similarly, geometric throat variables $(a,L)$, pole scales such as $\Omega_Q$,
and PN coefficients such as $d_3$ or $\mu_{\rho,3}$ should remain visibly
distinct.  The rest of the paper uses all three layers at once, and most
avoidable confusion at this stage is not physics but notation drift.

Finally, in the static-completion subsection it will be convenient to write
\[
p\equiv m_A,
\qquad
q\equiv m_B,
\]
but only there, and only as mass shorthands.  They are not electric charges.

\subsection{Exact theorem target}
\label{sec:dict-target}

The theorem target of the paper is not an unconstrained coefficient fit.  It has
a strict one-body front door, an exact imported reduced target, and a final
fixed-chart split.

\paragraph{Strict one-body gate.}
The exact isotropic Schwarzschild target through $c^{-6}$ is
\[
\frac{L_{3\mathrm{PN},1\mathrm{body}}}{m}
=
\frac5{128}v^8
+\frac{11}{16}Uv^6
+\frac{47}{16}U^2v^4
+\frac{13}{8}U^3v^2
-\frac18U^4.
\]
Any acceptable $3$PN one-body/self sector must reproduce this answer exactly.
The first task of the paper is therefore to show that the carried denominator-style
package can be repaired minimally by a finite three-parameter ledger rather than
by an uncontrolled extension of the lower-order self sector.

\paragraph{Exact reduced COM target.}
Once the natural symmetric self/static seed is subtracted from the imported exact
GR COM target, the remaining $15$-slot residual vector is
\[
\Delta l_1=\frac{3\nu(3\nu-1)(4\nu-1)}{16},
\qquad
\Delta l_2=\Delta l_3=\Delta l_4=\Delta l_5=0,
\]
\[
\begin{aligned}
\Delta l_6&=\frac{\nu(38-116\nu-57\nu^2)}{16},
& \Delta l_7&=\frac{\nu^2(20-69\nu)}{16},
\\
\Delta l_8&=\frac{3\nu^2(3-11\nu)}{16},
& \Delta l_9&=\frac{5\nu^3}{16},
\end{aligned}
\]
\[
\begin{aligned}
\Delta l_{10}&=\frac{\nu(129-98\nu+52\nu^2)}{16},
& \Delta l_{11}&=\frac{\nu(-3+52\nu+124\nu^2)}{16},
\\
\Delta l_{12}&=\frac{\nu(-5+11\nu+48\nu^2)}{12}.
\end{aligned}
\]
\[
\Delta l_{13}=-\frac{\nu(244+3\pi^2+1272\nu+96\nu^2)}{192},
\qquad
\Delta l_{14}=\frac{\nu(452+3\pi^2-384\nu-224\nu^2)}{64},
\]
\[
\Delta l_{15}=\frac{\nu(-908+63\pi^2)}{96}.
\]
This vector is the exact reduced target data against which every later
compiler, grouped-$P_2$, and geometry statement is checked.

\paragraph{Fixed-chart theorem target.}
Within the same declared closure hierarchy used successfully at $1$PN and $2$PN,
the goal is to show all of the following.
\begin{enumerate}
  \item The exact one-body repair ledger closes the strict Schwarzschild gate.

  \item The fixed ADM chart generic-frame compiler reduces to
  \[
  \Delta H_3=-\Delta L_3,
  \]
  so the correct generic-frame representative is unique once the Hamiltonian
  target is imposed.

  \item The full nine-slot middle interaction block is carried exactly by the
  richer grouped real $P_2$ family
  \[
  (T_{20},T_{21},T_{22},S_{20},S_{21},S_{22},V_{20},V_{21},V_{22}).
  \]

  \item The remaining static scalar slot is uniquely geometry-side after the
  exact sigma-collapse removes the apparent $P_0$/geometry ambiguity.

  \item The apparent pure-kinetic leftover is only the ordinary-Lagrangian
  counterimage of the universal free relativistic two-body COM Hamiltonian,
  rather than a new throat datum.
\end{enumerate}
Equivalently, the exact comparable-mass $3$PN residual to be proved in the rest
of the paper is
\[
\boxed{
\begin{aligned}
\Delta \hat L_3^{\rm GR}
=\;&\underbrace{\Delta l_1 v^8}_{\text{compiler image of the free COM Hamiltonian}}
+\underbrace{L_{P_2}^{\rm mid}}_{\text{exact grouped real }P_2\text{ closure}}
\\
&+\underbrace{\Delta l_{15}^{(g)}U^4}_{\text{unique geometry completion}}.
\end{aligned}
}
\]
This is the most compact statement of the theorem target.  It says that once the
strict one-body gate, the exact reduced COM target, the Hamiltonian-first lift,
the grouped real $P_2$ module, and the static sigma-collapse are all handled
correctly, the conservative $3$PN bottleneck is no longer algebraic.

The correct theorem-status reading is the same as in the earlier full
conservative papers.  The statement above is a \emph{full assembly target within
a declared closure hierarchy}, not an assumption-free theorem of the fully solved
moving-throat PDE.  What remains open after the present paper is therefore not
``more $3$PN coefficients.''  What remains open is the deeper moving-throat
derivation of the grouped-$P_2$ constitutive law and unique geometry lane, and
then the insertion of that data into the already narrowed passive/outgoing
quadrupole bridge.

% Section content for 4d_3pn.tex

\section{One-body \texorpdfstring{$3$}{3}PN gate, self/static seed, and exact cubic compiler}
\label{sec:onebody}

The first genuine theorem gate of the $3$PN program is the strict test-mass
sector.  Before any comparable-mass lift is attempted, one must know exactly
what the carried one-body/self-sector is allowed to reduce to, and one must know
whether the lower-order denominator/static packaging survives unchanged.  At
$3$PN it does not.  The strict one-body audit shows that the order-$c^{-6}$
sector is already richer than a one-parameter continuation of the $2$PN closure:
there is an exact Schwarzschild target, a minimal three-parameter repair ledger,
and then a unique natural symmetric self/static promotion that isolates all
bodywise content before the genuine comparable-mass residual is solved.  The
second ingredient needed before that lift can become finite is the exact
cubic-order perturbative Legendre compiler.  Once the carried $1$PN/$2$PN
ledger is frozen, that compiler makes the $3$PN Hamiltonian content a rigid
algebraic image of the new ordinary-Lagrangian block rather than a blind full
recomputation.

\subsection{Exact Schwarzschild test-mass target through \texorpdfstring{$3$}{3}PN}
\label{sec:onebody-target}

The strict one-body gate is the isotropic Schwarzschild worldline action.  Using
\[
u\equiv \frac{U}{c^2}=\frac{GM}{rc^2},
\]
the exact test-mass Lagrangian in isotropic Schwarzschild coordinates is
\[
\frac{L_{\rm Schw}}{m}
=
-c^2\sqrt{\left(\frac{1-u/2}{1+u/2}\right)^2-(1+u/2)^4\frac{v^2}{c^2}}.
\]
Expanding through total order $c^{-6}$ gives
\[
\frac{L_{\rm Schw}}{m}
={}-c^2+\frac12v^2+U
+\frac{1}{c^2}\left(\frac18v^4+\frac32Uv^2-\frac12U^2\right)
\]
\[
\qquad
+\frac{1}{c^4}\left(\frac1{16}v^6+\frac78Uv^4+2U^2v^2+\frac14U^3\right)
\]
\[
\qquad
+\frac{1}{c^6}\left(
\frac5{128}v^8
+\frac{11}{16}Uv^6
+\frac{47}{16}U^2v^4
+\frac{13}{8}U^3v^2
-\frac18U^4
\right)
+O(c^{-8}).
\]
So the exact $3$PN one-body destination is
\[
\boxed{
\frac{L_{3\mathrm{PN},1\mathrm{body}}}{m}
=
\frac5{128}v^8
+\frac{11}{16}Uv^6
+\frac{47}{16}U^2v^4
+\frac{13}{8}U^3v^2
-\frac18U^4.
}
\]
This plays exactly the same role here that the strict Schwarzschild test-mass
check played at lower order: it is the unique one-body answer that every later
comparable-mass construction must reduce to.

Two features of the target are worth recording immediately.  First, the
universal $v^8$ and $Uv^6$ coefficients continue the lower-order relativistic
pattern exactly.  Second, the remaining three slots
\[
U^2v^4,
\qquad
U^3v^2,
\qquad
U^4,
\]
already show that the full $3$PN self sector cannot be inferred by naive pattern
continuation alone.  The rest of the subsection fixes precisely how much new
one-body information is needed.

\subsection{Minimal one-body repair ledger}
\label{sec:onebody-repair}

The first algebraic test is whether the carried denominator-style self sector can
simply be pushed one order higher.  Write the minimal cubic extension as
\[
L_{\rm red}
=-mc^2(1-u)\sqrt{1-\frac{v^2/c^2}{D(u)}},
\qquad
D(u)=1-4u+8u^2+d_3u^3+O(u^4).
\]
Expanding through $c^{-6}$ yields
\[
\frac{L_{\rm red}}{m}
={}-c^2+\frac12v^2+U
+\frac{1}{c^2}\left(\frac18v^4+\frac32Uv^2\right)
\]
\[
\qquad
+\frac{1}{c^4}\left(\frac1{16}v^6+\frac78Uv^4+2U^2v^2\right)
\]
\[
\qquad
+\frac{1}{c^6}\left(
\frac5{128}v^8
+\frac{11}{16}Uv^6
+3U^2v^4
+\left(-4-\frac{d_3}{2}\right)U^3v^2
\right)
+O(c^{-8}).
\]
So the carried denominator logic already reproduces the universal $v^8$ and
$Uv^6$ coefficients automatically, but it leaves the $U^3v^2$ and $U^2v^4$
slots unresolved.  Matching the exact Schwarzschild coefficient of $U^3v^2$
forces
\[
-4-\frac{d_3}{2}=\frac{13}{8}
\qquad\Longrightarrow\qquad
\boxed{d_3=-\frac{45}{4}.}
\]
It is useful to note that the unextended $2$PN invariant-denominator philosophy
would have predicted instead
\[
d_3^{\rm carry}=\frac{21783}{2432},
\]
which is nowhere near the exact target.  So already at the $U^3v^2$ level,
$3$PN requires a genuinely new cubic invariant correction if one wants to retain
that denominator-style packaging.

The pure static slot gives an independent gate.  Extend the local mass-scaling
sector one order further by a cubic static coefficient $\mu_{\rho,3}$:
\[
\frac{L_{\rm static,1body}}{m}
=
U-\frac12\frac{U^2}{c^2}+\frac14\frac{U^3}{c^4}
-\frac{\mu_{\rho,3}}{2}\frac{U^4}{c^6}+O(c^{-8}).
\]
Matching the exact static Schwarzschild term
\[
-\frac18\frac{U^4}{c^6}
\]
then gives
\[
-\frac{\mu_{\rho,3}}{2}=-\frac18
\qquad\Longrightarrow\qquad
\boxed{\mu_{\rho,3}=\frac14.}
\]
So the denominator datum and the static gate are both fixed exactly and
independently.

At this point the one-body problem is still not closed.  After imposing
\(d_3=-45/4\) and \(\mu_{\rho,3}=1/4\), the denominator-style self sector still
predicts
\[
3U^2v^4=\frac{48}{16}U^2v^4,
\]
whereas the exact Schwarzschild target demands
\[
\frac{47}{16}U^2v^4.
\]
So a residual one-body discrepancy of
\[
-\frac1{16}U^2v^4
\]
remains.  This cannot be absorbed by $d_3$, which controls only the $U^3v^2$
slot, and it cannot be absorbed by $\mu_{\rho,3}$, which controls only the pure
static $U^4$ slot.  The minimal repair is therefore to admit one explicit new
self invariant,
\[
s_{24}\frac{U^2v^4}{c^6},
\qquad
\boxed{s_{24}=-\frac1{16}.}
\]
Collecting everything, the exact one-body repair ledger is
\[
\boxed{
\mu_{\rho,3}=\frac14,
\qquad
 d_3=-\frac{45}{4},
\qquad
 s_{24}=-\frac1{16}.
}
\]
This is the cleanest strict one-body lesson of the $3$PN analysis.  The $3$PN
gate does \emph{not} say ``continue the $2$PN denominator idea by one more
coefficient.''  It says that three distinct data are already needed at the
strict test-mass level: a cubic static gate, a cubic denominator datum, and one
genuinely new self-sector slot.

\subsection{Natural \texorpdfstring{$3$}{3}PN self/static seed}
\label{sec:onebody-seed}

Once the exact one-body target and its minimal repair ledger are frozen, the
natural next step is the same one that worked at $2$PN: promote the strict
one-body structure symmetrically into a two-body self/static seed while keeping
all genuinely comparable-mass interaction data outside the seed.  The rule is to
remain strictly bodywise and symmetric so that the exact test-mass reduction is
automatic.

Start from the exact one-body $3$PN target
\[
\frac{L_{3\mathrm{PN},1\mathrm{body}}}{m}
=
\frac5{128}v^8
+\frac{11}{16}Uv^6
+\frac{47}{16}U^2v^4
+\frac{13}{8}U^3v^2
-\frac18U^4.
\]
For body $A$ in the field of body $B$, use
\[
U_A=\frac{Gm_B}{r},
\qquad
m\to m_A,
\qquad
v\to v_A,
\]
and for body $B$ in the field of body $A$, use
\[
U_B=\frac{Gm_A}{r},
\qquad
m\to m_B,
\qquad
v\to v_B.
\]
Adding the two bodywise copies gives the natural symmetric self/static $3$PN
seed,
\[
\boxed{
L_{3,\mathrm{seed}}
=
\frac{5}{128}(m_Av_A^8+m_Bv_B^8)
+\frac{11Gm_A m_B}{16r}(v_A^6+v_B^6)
}
\]
\[
\boxed{
\qquad
+\frac{47G^2m_A m_B}{16r^2}(m_Bv_A^4+m_Av_B^4)
+\frac{13G^3m_A m_B}{8r^3}(m_B^2v_A^2+m_A^2v_B^2)
-\frac{G^4m_A m_B}{8r^4}(m_B^3+m_A^3),
}
\]
with the understanding that this entire block is the coefficient of $c^{-6}$.

This seed has three properties that make it the correct starting object for the
$3$PN lift.

First, it is manifestly symmetric under $A\leftrightarrow B$.

Second, it preserves the strict one-body gate exactly.  In the test-mass limit,
the light-body sector reduces to the exact Schwarzschild target, while the
heavy-body copy collapses to a self piece that does not contaminate the test-mass
dynamics.

Third, it isolates the universal bodywise content before the genuine
comparable-mass residual is solved.  In the later COM chart this means the seed
is sparse: before lower-order compiler feedback, only the self/static slots are
switched on directly, and the remaining $3$PN structure can be interpreted
cleanly as residual beyond this frozen bodywise core.

That is why this seed is the right starting point for the full $3$PN lift.  It
is neither a guess about the complete comparable-mass answer nor a vague
``start somewhere'' ansatz.  It is the unique symmetric promotion of the exact
one-body $3$PN data into the same self/static lane that already organized the
$2$PN construction successfully.

\subsection{Exact cubic Legendre-transform identity}
\label{sec:onebody-legendre}

The second indispensable ingredient is the exact cubic-order perturbative
Legendre compiler.  The $2$PN assembly already relied on the perturbative
Legendre transform through quadratic order.  At $3$PN the new content enters at
cubic order, and once the carried $1$PN/$2$PN ledger is frozen, the resulting
Hamiltonian feedback is rigid.

Write the conservative ordinary Lagrangian as
\[
L(v)=L_0(v)+\varepsilon L_1(v)+\varepsilon^2L_2(v)+\varepsilon^3L_3(v),
\qquad
\varepsilon=c^{-2},
\]
with quadratic zeroth-order block
\[
L_0(v)=\frac12 v^T M v,
\]
where the mass matrix $M$ is constant at the order of the perturbative
inversion.  Define the Newtonian momentum-space velocity
\[
v_0=M^{-1}p,
\]
and the lower-order derivative data at $v_0$,
\[
A_0\equiv \left.\frac{\partial L_1}{\partial v}\right|_{v_0},
\qquad
B_0\equiv \left.\frac{\partial L_2}{\partial v}\right|_{v_0},
\qquad
C_0\equiv \left.\frac{\partial^2L_1}{\partial v^2}\right|_{v_0}.
\]
Then the perturbative momentum inversion
\[
v=v_0+\varepsilon v_1+\varepsilon^2v_2+O(\varepsilon^3)
\]
solves exactly to
\[
v_1=-M^{-1}A_0,
\qquad
v_2=M^{-1}C_0M^{-1}A_0-M^{-1}B_0.
\]
Substituting into
\[
H(p)=p\cdot v-L(v)
\]
gives the cubic-order Hamiltonian expansion
\[
H=H_0+\varepsilon H_1+\varepsilon^2H_2+\varepsilon^3H_3+O(\varepsilon^4),
\]
with
\[
H_1=-L_1(v_0),
\]
\[
H_2=-L_2(v_0)+\frac12A_0^TM^{-1}A_0,
\]
and the key $3$PN identity
\[
\boxed{
H_3=-L_3(v_0)+A_0^TM^{-1}B_0-
\frac12A_0^TM^{-1}C_0M^{-1}A_0.
}
\]
This is the exact $3$PN analogue of the quadratic compiler identity used at
$2$PN.  Its practical meaning is that once the lower-order $1$PN/$2$PN blocks
are frozen, the genuinely new $3$PN Hamiltonian content enters only through the
explicit term $-L_3(v_0)$, while the feedback pieces are already determined by
the carried lower-order ledger.  That is precisely what turns the later
comparable-mass lift into a finite algebraic residual problem rather than a
blind variational search.

The next section will apply this compiler first in the reduced COM chart and
then, more importantly, in the fixed generic-frame Hamiltonian chart.  The key
lesson to keep in view is already visible here: at $3$PN the right bridge is not
``reduce first and hope.''  It is exact compiler algebra first, followed only
then by the reduced comparison and fixed-chart uniqueness analysis.

% Section content for 4d_3pn.tex

\section[Comparable-mass scaffold, COM target, and fixed-chart uniqueness]{Comparable-mass scaffold, COM target,\\ and fixed-chart uniqueness}
\label{sec:compmass}

The strict one-body audit carried out immediately before this section fixes everything that is
universal and bodywise at $3$PN, but it does not yet solve the genuinely
comparable-mass problem.  The next task is to convert the full nonspinning
order-$c^{-6}$ ledger into a finite algebraic residual, to determine exactly
what the center-of-mass reduction does and does not fix, and then to prove that
the remaining apparent generic-frame ambiguity disappears once the fixed ADM
Hamiltonian chart is imposed.  This is therefore the section in which the $3$PN
program stops being ``more coefficients somewhere'' and becomes a rigid
scaffold--compiler--quotient chain.

The logic is simple but unforgiving.  First, one freezes a reduced $15$-slot COM
basis and the exact cubic ordinary/Hamiltonian compiler in that chart.  Second,
one imports the exact GR COM target and isolates the full residual beyond the
frozen self/static seed.  Third, one lifts that residual to a finite
exchange-symmetric generic-frame scaffold and checks what COM projection still
fails to determine.  Fourth, one diagnoses the middle-block obstruction and
repairs it by a tiny $\nu$-dressed seed sector that vanishes in the strict
one-body limit.  Finally, one proves that the decisive quotient principle is
Hamiltonian-first rather than naive ordinary-Lagrangian COM reduction, and that
in the fixed ADM chart the full generic-frame residual compiler is just the
slotwise sign flip
\[
\Delta H_3=-\Delta L_3.
\]
At that point the comparable-mass $3$PN bottleneck is no longer algebraic.  What
remains is the physical interpretation of the exact residual already isolated in
this section.

\subsection{Reduced COM \texorpdfstring{$15$}{15}-slot basis and diagonal ordinary/Hamiltonian compiler}
\label{sec:compmass-com-basis}

In the standard isotropic center-of-mass chart, and suppressing the overall
reduced-mass normalization already fixed by the lower-order papers, write the
conservative reduced ordinary Lagrangian as
\[
\hat L=\hat L_0+c^{-2}\hat L_1+c^{-4}\hat L_2+c^{-6}\hat L_3+O(c^{-8}),
\]
with $\hat L_0$, $\hat L_1$, and $\hat L_2$ already frozen.  The new $3$PN
coefficient is expanded in the natural $15$-slot ordinary basis
\[
\hat L_3
= l_1 v^8 + l_2 v^6\dot r^2 + l_3 v^4\dot r^4 + l_4 v^2\dot r^6 + l_5\dot r^8
\]
\[
\qquad
+\frac1r\Bigl(l_6 v^6 + l_7 v^4\dot r^2 + l_8 v^2\dot r^4 + l_9\dot r^6\Bigr)
\]
\[
\qquad
+\frac1{r^2}\Bigl(l_{10} v^4 + l_{11} v^2\dot r^2 + l_{12}\dot r^4\Bigr)
+\frac1{r^3}\Bigl(l_{13} v^2 + l_{14}\dot r^2\Bigr)
+\frac{l_{15}}{r^4}.
\]
The matching isotropic COM Hamiltonian basis is
\[
\hat H_3
= h_1 p^8 + h_2 p^6p_r^2 + h_3 p^4p_r^4 + h_4 p^2p_r^6 + h_5 p_r^8
\]
\[
\qquad
+\frac1r\Bigl(h_6 p^6 + h_7 p^4p_r^2 + h_8 p^2p_r^4 + h_9 p_r^6\Bigr)
\]
\[
\qquad
+\frac1{r^2}\Bigl(h_{10} p^4 + h_{11} p^2p_r^2 + h_{12} p_r^4\Bigr)
+\frac1{r^3}\Bigl(h_{13} p^2 + h_{14} p_r^2\Bigr)
+\frac{h_{15}}{r^4},
\]
where $p^2$ and $p_r$ are the reduced isotropic COM momentum invariants.

The importance of this chart is that the exact cubic Legendre transform carried from the preceding one-body section diagonalizes slot by slot.  The full
forward map is
\[
\boxed{h_1=-l_1+\frac{3}{16}\nu-\frac{21}{16}\nu^2+\frac94\nu^3}
\]
\[
\boxed{h_2=-l_2,
\qquad h_3=-l_3,
\qquad h_4=-l_4,
\qquad h_5=-l_5}
\]
\[
\boxed{h_6=-l_6+\frac14+\frac78\nu-\frac{35}{8}\nu^2-\frac{21}{4}\nu^3}
\]
\[
\boxed{h_7=-l_7+\frac{11}{8}\nu^2-\frac92\nu^3,
\qquad
h_8=-l_8+\frac34\nu^2-\frac94\nu^3,
\qquad
h_9=-l_9}
\]
\[
\boxed{h_{10}=-l_{10}+\frac54+\frac{15}{8}\nu+\frac{123}{8}\nu^2+\frac{13}{4}\nu^3}
\]
\[
\boxed{h_{11}=-l_{11}+\frac78\nu+\frac{41}{8}\nu^2+\frac{31}{4}\nu^3}
\]
\[
\boxed{h_{12}=-l_{12}+\frac92\nu^2+4\nu^3}
\]
\[
\boxed{h_{13}=-l_{13}-\frac32-\frac{59}{4}\nu-\frac{25}{4}\nu^2-\frac12\nu^3}
\]
\[
\boxed{h_{14}=-l_{14}+\frac74\nu-\frac{31}{4}\nu^2-\frac72\nu^3,
\qquad
h_{15}=-l_{15}.}
\]
The inverse map is immediate: one simply moves the known lower-order feedback
polynomials to the opposite side.  This is the decisive reduced-chart fact that
makes the $3$PN problem finite.  Once the lower-order ledger is frozen, the
imported GR COM target can be solved slot by slot, and the entire comparable-mass
question reduces to an explicit residual rather than a blind variational search.

\subsection{Imported exact GR COM target and residual extraction}
\label{sec:compmass-com-target}

With the diagonal compiler in hand, import the exact GR $3$PN reduced COM target
in the standard ADM chart and subtract the natural one-body/self-static seed
constructed in the preceding one-body section.  The result is the exact comparable-mass
residual vector
\[
\Delta l_1=\frac{3\nu(3\nu-1)(4\nu-1)}{16},
\qquad
\Delta l_2=\Delta l_3=\Delta l_4=\Delta l_5=0,
\]
\[
\begin{aligned}
\Delta l_6&=\frac{\nu(38-116\nu-57\nu^2)}{16},
& \Delta l_7&=\frac{\nu^2(20-69\nu)}{16},
\\
\Delta l_8&=\frac{3\nu^2(3-11\nu)}{16},
& \Delta l_9&=\frac{5\nu^3}{16},
\end{aligned}
\]
\[
\begin{aligned}
\Delta l_{10}&=\frac{\nu(129-98\nu+52\nu^2)}{16},
& \Delta l_{11}&=\frac{\nu(-3+52\nu+124\nu^2)}{16},
\\
\Delta l_{12}&=\frac{\nu(-5+11\nu+48\nu^2)}{12}.
\end{aligned}
\]
\[
\Delta l_{13}=-\frac{\nu(244+3\pi^2+1272\nu+96\nu^2)}{192},
\qquad
\Delta l_{14}=\frac{\nu(452+3\pi^2-384\nu-224\nu^2)}{64},
\]
\[
\Delta l_{15}=\frac{\nu(-908+63\pi^2)}{96}.
\]
This is the exact reduced $3$PN target data set against which every later
scaffold, quotient, grouped-$P_2$, and geometry statement is checked.

The pattern of the residual is already highly informative.  The first five slots
show that the only free-particle mismatch beyond the self/static seed is the
single $v^8$ coefficient $\Delta l_1$; there is no comparable-mass residual in
any of the other pure-kinetic monomials.  The nine slots
\[
(\Delta l_6,\dots,\Delta l_{14})
\]
constitute the full middle interaction block, while the single static coefficient
$\Delta l_{15}$ is the entire quartic potential remainder.  So before any
generic-frame lifting is attempted, the COM data already exhibit the exact split
that will survive to the end of the paper:
\begin{enumerate}
  \item a one-dimensional free-kinetic correction,
  \item a nine-dimensional middle interaction block,
  \item and a one-dimensional static completion.
\end{enumerate}
The rest of this section explains how these three pieces appear in the generic
frame, why COM data alone do not yet fix the answer uniquely, and why the fixed
Hamiltonian chart eventually removes that ambiguity completely.

\subsection{Generic-frame residual scaffold}
\label{sec:compmass-generic-scaffold}

To lift the COM data back to the full two-body chart, use the exchange-symmetric
generic-frame invariants
\[
 a=v_A^2,
 \qquad
 b=v_B^2,
 \qquad
 c=\mathbf v_A\!\cdot\!\mathbf v_B,
 \qquad
 d=v_{An},
 \qquad
 e=v_{Bn},
\]
with
\[
 p=m_A,
 \qquad
 q=m_B.
\]
After subtracting the frozen one-body/self-static seed, the honest generic-frame
comparable-mass scaffold is
\[
L_{3,\rm full}=L_{3,\rm seed}+L_{3,\rm res},
\]
with
\[
L_{3,\rm res}=
\frac{Gm_A m_B}{r}\sum_{i=1}^{24}q_iQ_i
+\frac{G^2m_A m_B}{r^2}\sum_{j=1}^{17}t_jT_j
+\frac{G^3m_A m_B}{r^3}\sum_{k=1}^{8}s_kS_k
+\frac{G^4m_A m_B}{r^4}u_1U_{01},
\]
where the unique exchange-symmetric static cross polynomial is
\[
U_{01}=p^2q+pq^2.
\]
The important point is not the eventual sparsity of the solved answer but the
size and symmetry of the residual search space before any contact/gauge quotient
is imposed.  The full generic-frame interaction problem is finite from the
start: the block sizes are
\[
24,\qquad 17,\qquad 8,\qquad 1,
\]
for the $G/r$, $G^2/r^2$, $G^3/r^3$, and $G^4/r^4$ blocks, for a total of $50$
interaction directions.

Projecting this scaffold to COM gives the exact blockwise image dimensions
\[
\boxed{12,\qquad 6,\qquad 4,\qquad 1,}
\]
corresponding to the four $G/r$ slots $(l_6,l_7,l_8,l_9)$ each supporting
$\nu$, $\nu^2$, and $\nu^3$, the three $G^2/r^2$ slots
$(l_{10},l_{11},l_{12})$ supporting only $\nu$ and $\nu^2$, the two
$G^3/r^3$ slots $(l_{13},l_{14})$ again supporting only $\nu$ and $\nu^2$, and
the single static slot $l_{15}$ supporting only a coefficient proportional to
$\nu$.  So the blockwise nullities are
\[
\boxed{12,\qquad 11,\qquad 4,\qquad 0,}
\]
and the total COM-blind interaction freedom is
\[
\boxed{27.}
\]
This is the first exact statement of what the solved COM target can and cannot
determine by itself.  The scaffold is large enough; the question is how to place
the seed correctly and how to quotient the COM-blind family once the Hamiltonian
chart is imposed.

\subsection{COM projection, middle-block obstruction, and seed repair}
\label{sec:compmass-obstruction}

The first lift test is whether the natural one-body/self-static seed already
lands in the correct COM image block by block.  The answer is almost yes, but
not quite.  The $G/r$ sextic block and the $G^4/r^4$ static slot are already
compatible with the exact COM target.  The only obstruction sits in the middle
$G^2/r^2$ and $G^3/r^3$ blocks.

The reason is simple.  In the constant-coefficient generic-frame image, the
$G^2/r^2$ and $G^3/r^3$ blocks carry only $\nu$ and $\nu^2$ support under COM
projection.  But after subtracting the natural seed, the exact COM residual
still contains small $\nu^3$ tails in precisely those blocks.  The obstruction
pieces are
\[
\boxed{
\delta l_{10}^{\rm obs}=\frac{13}{4}\nu^3,
\qquad
\delta l_{11}^{\rm obs}=\frac{31}{4}\nu^3,
\qquad
\delta l_{12}^{\rm obs}=4\nu^3,}
\]
for the $G^2/r^2$ block, and
\[
\boxed{
\delta l_{13}^{\rm obs}=-\frac12\nu^3,
\qquad
\delta l_{14}^{\rm obs}=-\frac72\nu^3,}
\]
for the $G^3/r^3$ block.  This is not a failure of the scaffold.  It is a
seed-placement problem: the natural self/static seed is exact in the strict
one-body sector, but it is not yet aligned with the mass-degree structure of the
generic-frame middle blocks.

The repair is remarkably small.  Define the refined seed correction sector
\[
\boxed{
\Delta T_\nu
=
\frac{pq}{4(p+q)}\Bigl(13ab+31cde+16d^2e^2\Bigr),}
\]
\[
\boxed{
\Delta S_\nu
=
\frac{p^2q^2}{2(p+q)^2}\Bigl(c+7de\Bigr).}
\]
Then the refined seed is
\[
L_{3,\rm seed}^{\rm ref}
=
L_{3,\rm seed}
+\frac{G^2pq}{r^2}\,\Delta T_\nu
+\frac{G^3pq}{r^3}\,\Delta S_\nu.
\]
These extra pieces do exactly what is needed:
\begin{enumerate}
  \item they vanish in the strict one-body limit because they are proportional
  to $\nu$,
  \item they leave the $G/r$ and $G^4/r^4$ blocks untouched,
  \item and under COM reduction they reproduce precisely the missing
  $\nu^3$ obstruction tails.
\end{enumerate}
So the middle-block obstruction is bookkeeping rather than ontology.  After this
repair, each interaction block admits an exact COM-compatible generic-frame
representative, and what remains is no longer a compatibility failure but a
quotient problem.

\subsection{The COM-null ideal and the canonical quotient slice}
\label{sec:compmass-com-null}

Once the refined seed is fixed, the remaining generic-frame freedom is exactly
the family invisible in the COM quotient algebra.  In the COM frame the body
velocities satisfy the three linear-momentum identities
\[
\boxed{
\mathcal C_1\equiv pa+qc=0,
\qquad
\mathcal C_2\equiv pc+qb=0,
\qquad
\mathcal C_3\equiv pd+qe=0,}
\]
and the three collinearity identities
\[
\boxed{
\mathcal C_4\equiv ab-c^2=0,
\qquad
\mathcal C_5\equiv ae-cd=0,
\qquad
\mathcal C_6\equiv bd-ce=0.}
\]
These generate the COM-null ideal
\[
\boxed{
\mathfrak I_{\rm COM}
=
\langle \mathcal C_1,\mathcal C_2,\mathcal C_3,\mathcal C_4,\mathcal C_5,\mathcal C_6\rangle.}
\]
The full $27$-dimensional nullspace of the $24/17/8/1$ COM projection matrix is
exactly the generic-frame freedom invisible in the quotient algebra
\[
\mathbb Q[a,b,c,d,e,p,q]/\mathfrak I_{\rm COM}.
\]
A useful refinement is that the $T$- and $S$-block null directions already lie
in the smaller linear-momentum ideal
\[
\langle \mathcal C_1,\mathcal C_2,\mathcal C_3\rangle,
\]
whereas the $Q$-block null directions need the full ideal
$\mathfrak I_{\rm COM}$.

This makes it natural to define a canonical sparse algebraic representative by
quotienting by the COM-null ideal and setting those $27$ null coefficients to
zero.  In that canonical slice one obtains an explicit generic-frame answer that
reproduces the full solved COM target exactly.  The details are lengthy and are
best recorded appendix-side, but the conceptual point belongs in the main text:
this canonical quotient slice is an exact algebraic representative, not a claim
that every COM-blind direction is ordinary gauge.  That distinction becomes
crucial in the next subsection.

\subsection{Contact generator, gauge orbit, and the \texorpdfstring{$11$}{11}-versus-\texorpdfstring{$27$}{27} distinction}
\label{sec:compmass-gauge-orbit}

A $3$PN ordinary contact transformation is generated by total derivatives of odd
scalars.  With the correct exchange rule
\[
A\leftrightarrow B,
\qquad
c\mapsto c,
\qquad
d\mapsto -e,
\qquad
e\mapsto -d,
\qquad
p\leftrightarrow q,
\]
the raw odd family contains
\[
\begin{aligned}
31&\quad (r\times\text{degree-}7),
& 12&\quad (G\times\text{degree-}5),
\\
8&\quad \left(\frac{G^2}{r}\times\text{degree-}3\right),
& 2&\quad \left(\frac{G^3}{r^2}\times\text{degree-}1\right),
\end{aligned}
\]
for a total of
\[
\boxed{53}
\]
candidates.  Imposing the scaffold-preserving condition cuts this to
\[
\boxed{22}
\]
dimensions, and imposing zero COM image cuts it once more.  The actual
scaffold-preserving and COM-preserving ordinary contact/gauge orbit is therefore
\[
\boxed{11\text{-dimensional}.}
\]
So the exact reduction chain is
\[
53\longrightarrow 22\longrightarrow 11.
\]

A sparse basis for this true ordinary contact orbit factors only through the
smaller subideal generated by
\[
\mathcal C_3,\qquad \mathcal C_4,\qquad \mathcal C_5,\qquad \mathcal C_6,
\]
not through all of $\mathfrak I_{\rm COM}$.  In particular, no generator uses
$\mathcal C_1$ or $\mathcal C_2$, and the simple ordinary contact family leaves
the static block untouched.  Its image inside the compact generic-frame scaffold
has blockwise ranks
\[
\boxed{\operatorname{rank}_Q=6,
\qquad
\operatorname{rank}_T=9,
\qquad
\operatorname{rank}_S=4,
\qquad
\operatorname{rank}_U=0.}
\]
So the actual ordinary contact orbit is only a strict subspace of the larger
$27$-dimensional COM-blind family.  This is the exact content of the
$11$-versus-$27$ distinction:
\begin{quote}
COM blindness and ordinary contact/gauge freedom are not identical.
\end{quote}
The canonical quotient slice is therefore the right sparse algebraic reference
chart, but it is not yet the final physical uniqueness theorem.

\subsection{Why naive ordinary COM reduction fails}
\label{sec:compmass-naive-com}

The decisive subtlety appears once the full generic-frame ordinary target is
imported from the standard $3$PN literature, with ambiguity fixed by
\[
\lambda=-\frac{3}{11}\omega_s-\frac{1987}{3080},
\qquad
\omega_s=0,
\qquad
\Rightarrow
\qquad
\lambda=-\frac{1987}{3080}.
\]
After subtraction of the frozen natural one-body/self-static seed, the exact
ordinary generic-frame residual lands entirely inside the compact $24/17/8/1$
scaffold and still vanishes in the strict test-mass branch, so the imported
target itself is perfectly consistent with the one-body gate.

The problem is not the target.  The problem is the order of reduction.
Projecting the imported generic-frame \emph{ordinary} target directly to COM does
not reproduce the exact COM ordinary target.  The mismatches are
\[
\delta l_6=\frac{\nu(80\nu^2+48\nu-17)}{16},
\qquad
\delta l_7=\frac{3\nu(24\nu^2-10\nu+1)}{16},
\qquad
\delta l_8=\frac{3\nu^2(4\nu-1)}{8},
\qquad
\delta l_9=0,
\]
\[
\delta l_{10}=\frac{\nu(-52\nu^2-123\nu+34)}{16},
\qquad
\delta l_{11}=\frac{\nu(-124\nu^2-97\nu+32)}{16},
\qquad
\delta l_{12}=\nu^2(1-4\nu),
\]
\[
\delta l_{13}=\frac{\nu(4\nu^2+27\nu-7)}{8},
\qquad
\delta l_{14}=\frac{\nu(28\nu^2+69\nu-19)}{8},
\qquad
\delta l_{15}=0.
\]
So naive ordinary COM reduction agrees only in the $l_9$ and $l_{15}$ slots.
This is not a contradiction.  It is exactly the same subtlety already known from
standard $3$PN dynamics: the reduced relative ordinary Lagrangian does not
follow from the general-frame ordinary representative by simple substitution.
Naive ordinary COM reduction is therefore not the final quotient principle.

\subsection{Hamiltonian-first lift and exact COM recovery}
\label{sec:compmass-hamiltonian-first}

The cure is to apply the exact cubic Legendre map in the generic frame first,
and only then reduce to COM.  Carrying forward the exact identity from the
preceding one-body section,
\[
H_3=-L_3(v_0)+A_0^TM^{-1}B_0-\frac12A_0^TM^{-1}C_0M^{-1}A_0,
\]
with
\[
v_0=M^{-1}P,
\qquad
A_0=\left.\frac{\partial L_1}{\partial v}\right|_{v_0},
\qquad
B_0=\left.\frac{\partial L_2}{\partial v}\right|_{v_0},
\qquad
C_0=\left.\frac{\partial^2L_1}{\partial v^2}\right|_{v_0},
\]
one computes the generic-frame $3$PN Hamiltonian first and substitutes the COM
relations only afterward.  That procedure gives the exact theorem
\[
\boxed{H_{3,\rm COM}^{\rm(lift)}=H_{3,\rm COM}^{\rm target}.}
\]
Equivalently,
\[
\boxed{
\text{generic-frame Legendre lift first}
\Longrightarrow
H_{3,\rm COM}^{\rm lift}=H_{3,\rm COM}^{\rm GR}.}
\]
So the earlier mismatch is diagnosed completely:
\[
\text{naive COM reduction of ordinary }L_3\neq\text{correct COM ordinary target},
\]
but the imported generic-frame ordinary target itself is fine.  The order of
operations was wrong, not the target.  This is the point where the $3$PN lift
stops being a suspicion about the generic frame and becomes a clean chart issue.

\subsection{Fixed-chart generic-frame uniqueness theorem}
\label{sec:compmass-uniqueness}

Once the Hamiltonian-first route is accepted, the remaining ambiguity question is
no longer ``is the target right?'' but ``which ordinary representative survives
in the fixed Hamiltonian chart?''  The answer is now exact.

Use the same formal invariant symbols as in the ordinary scaffold, but reinterpret
them in the generic-frame Hamiltonian chart as Newtonian-order momentum
invariants:
\[
 a=\frac{P_A^2}{m_A^2},
 \qquad
 b=\frac{P_B^2}{m_B^2},
 \qquad
 c=\frac{P_A\!\cdot\!P_B}{m_A m_B},
 \qquad
 d=\frac{n\!\cdot\!P_A}{m_A},
 \qquad
 e=\frac{n\!\cdot\!P_B}{m_B},
\]
with again
\[
 p=m_A,
 \qquad
 q=m_B.
\]
In this chart the generic-frame Hamiltonian interaction scaffold has exactly the
same $24/17/8/1$ block structure as the ordinary one.  Write
\[
L_3=L_{3,\rm seed}+\Delta L_3,
\]
where $L_{3,\rm seed}$ is the frozen one-body/self-static seed and
$\Delta L_3$ is the pure interaction residual on the $50$-slot scaffold.
Because the lower-order feedback terms in the cubic Legendre identity depend only
on the already frozen $1$PN/$2$PN ledger, they are insensitive to the new
$3$PN residual.  Therefore the full generic-frame residual compiler is exact and
trivial:
\[
\boxed{
\Delta H_3=-\Delta L_3(v_0).
}
\]
In the fixed generic-frame momentum basis, this is literally the slotwise
statement
\[
\boxed{\Delta h=-\Delta l,}
\]
so the full $50\times 50$ compiler matrix is simply
\[
\boxed{C_{\rm gen}=-I_{50}.}
\]
That is the uniqueness theorem.  None of the $27$ COM-blind directions is
Hamiltonian-blind in the fixed ADM chart.  Once the full generic-frame
Hamiltonian target is imposed, the entire COM-null family disappears at once.
Equivalently,
\[
\boxed{\text{COM blindness is not Hamiltonian blindness.}}
\]
The actual $11$-generator ordinary contact family is therefore not invisible in
the fixed Hamiltonian chart; it corresponds to moving to a different canonical
chart.  After the ADM chart is frozen, the Hamiltonian target selects a single
ordinary representative.

At that point the comparable-mass $3$PN lift is finished as an algebraic problem.
What remains is entirely physical: the nine-slot middle interaction block must be
carried by the grouped real $P_2$ conservative module, the last static slot must
be assigned to the scalar/geometry lane, and the isolated $v^8$ remainder must be
recognized as the ordinary-Lagrangian counterimage of the universal free
relativistic two-body COM Hamiltonian.  Those are precisely the tasks of the next
sections.

% Section content for 4d_3pn.tex

\section{The grouped real \texorpdfstring{$P_2$}{P 2} conservative module}
\label{sec:grouped-p2}

The previous section reduced the full comparable-mass $3$PN problem to three COM
pieces: one pure-kinetic slot, a nine-slot middle interaction block, and a
one-dimensional static remainder.  The present section addresses the middle
block.  It is the first place where the $2.5$PN narrowing to the
orbital/worldtube STF quadrupole branch becomes a hard conservative theorem gate
rather than a later radiative aspiration.  The question is not whether a
quadrupole-flavored ontology looks plausible in the abstract; the question is
whether the grouped real
\[
P_{20}\oplus P_{21}\oplus P_{22}
\]
bundle can actually reproduce the exact solved $3$PN middle block in the fixed
ADM chart.

The answer is two-part.  The smallest local grouped-$P_2$ demotion is too rigid:
it produces only a rank-three image and obeys six exact slot relations violated
by the GR target.  But the grouped ontology itself is not too small.  Once the
next natural local grouped families are added, the full nine-slot middle block
is reached exactly, with no middle-block residual left.  This is the first fully
successful conservative grouped-$P_2$ stress test in the series.

\subsection{Why grouped real \texorpdfstring{$P_{20}\oplus P_{21}\oplus P_{22}$}{P 20 + P 21 + P 22} is the right \texorpdfstring{$3$}{3}PN front end}
\label{sec:grouped-why}

The $2.5$PN audit already reduced the surviving universal point-particle route
to the orbital\allowbreak/worldtube STF quadrupole branch.  In conservative language, that
means the next higher-order front end is not an unfocused collection of generic
$3$PN coefficients, but the grouped real quadrupole bundle
\[
P_{20}\oplus P_{21}\oplus P_{22}.
\]
The exact $3$PN target inherited from the comparable-mass analysis just completed makes that statement
quantitative.  After the frozen self/static seed is subtracted, the full COM
residual splits into
\[
\Delta l_1 v^8,
\qquad
(\Delta l_6,\dots,\Delta l_{14}),
\qquad
\Delta l_{15}U^4.
\]
The grouped real $P_2$ lane can only hope to carry the middle block
\((\Delta l_6,\dots,\Delta l_{14})\), so it faces a solved nine-slot target from
the beginning.

At the same time, the grouped conservative front end through order
$O(\omega^2)$ contains only three real data,
\[
u_2^{(20)},\qquad u_2^{(21)},\qquad u_2^{(22)},
\]
or equivalently one grouped trace plus two anisotropy defects.  So even before
any microscopic constitutive law is chosen, the conservative grouped problem is
already overdetermined.  On the isotropic branch it becomes still sharper,
because only one real datum survives.  This is why $3$PN is the right first
conservative grouped-$P_2$ gate: it is the earliest order at which the surviving
quadrupole route is forced to confront a nontrivial exact target, but it is
still rigid enough that failure or success is unambiguous.

\subsection{Grouped conservative data and isotropy variables}
\label{sec:grouped-data}

At $3$PN the grouped conservative payload is carried by exactly three real
$O(\omega^2)$ moments,
\[
u_2^{(20)},\qquad u_2^{(21)},\qquad u_2^{(22)}.
\]
It is usually cleaner to rewrite them in terms of the grouped trace and two
anisotropy defects,
\begin{equation}
\bar u_2=\frac{u_2^{(20)}+2u_2^{(21)}+2u_2^{(22)}}{5},
\qquad
 a_2=\frac{2u_2^{(20)}-u_2^{(21)}-u_2^{(22)}}{10},
\qquad
 b_2=\frac{u_2^{(21)}-u_2^{(22)}}{2},
\label{eq:grouped-trace-defects}
\end{equation}
with inverse map
\begin{equation}
 u_2^{(20)}=\bar u_2+4a_2,
\qquad
 u_2^{(21)}=\bar u_2-a_2+b_2,
\qquad
 u_2^{(22)}=\bar u_2-a_2-b_2.
\label{eq:grouped-inverse-map}
\end{equation}
The exact isotropy gate for the minimal branch is therefore
\begin{equation}
\boxed{a_2=0,
\qquad
b_2=0.}
\label{eq:grouped-isotropy-gate}
\end{equation}
If that gate passes, then all three grouped moments coincide and one writes
\[
u_2^{(20)}=u_2^{(21)}=u_2^{(22)}\equiv u_2.
\]

The exact co-rotating pair-frame source map makes the front end completely
explicit.  Write
\[
d\equiv \dot r,
\qquad
u_\perp^2\equiv v^2-d^2,
\qquad
U\equiv \frac{GM}{r},
\]
so that the pair velocity decomposes as
\(\mathbf v=\mathbf u+d\,\mathbf n\) with
\(\mathbf u\!\cdot\!\mathbf n=0\).
Then the real-STF orbital source coefficients of
\(\ddot I_{ij}^{\rm orb}/(2\mu)\) imply the exact grouped derivative norms
\begin{equation}
\dot C_{20}^2=\frac{2d^2}{3r^2}(3u_\perp^2-U)^2,
\qquad
\dot C_{21c}^2+\dot C_{21s}^2=\frac{2u_\perp^2}{r^2}(u_\perp^2-d^2-U)^2,
\qquad
\dot C_{22c}^2+\dot C_{22s}^2=\frac{2d^2u_\perp^4}{r^2}.
\label{eq:grouped-derivative-norms}
\end{equation}
No constitutive assumption has entered yet.  Equation
\eqref{eq:grouped-derivative-norms} is a kinematic identity of the grouped
source map itself, and it is the unique starting point for any honest local
$3$PN grouped lift.

\subsection{Minimal local grouped-\texorpdfstring{$P_2$}{P 2} front end and why it fails}
\label{sec:grouped-minimal-failure}

The smallest time-local grouped conservative ansatz built from the exact source
norms is
\begin{equation}
L_{P_2,\omega^2}^{\rm front}
=
\frac{c_2^{(20)}}{3r^2}d^2(3u_\perp^2-U)^2
+
\frac{c_2^{(21)}}{r^2}u_\perp^2(u_\perp^2-d^2-U)^2
+
\frac{c_2^{(22)}}{r^2}d^2u_\perp^4.
\label{eq:grouped-minimal-front}
\end{equation}
Under the standard virial bookkeeping
\(U\sim v^2\sim d^2\), every monomial in
\eqref{eq:grouped-minimal-front} has uniform weight five.  So the direct front
end is formally too high to be a $3$PN ordinary-Lagrangian block.  If one insists
on the smallest possible local isotropic demotion, the unique one-step choice is
multiplication by one power of
\(U^{-1}\sim r\).

That test fails exactly.  After the one-step demotion, the grouped image lands
only in the nine interaction slots \((l_6,\dots,l_{14})\), has rank three, and
obeys six exact relations,
\begin{equation}
2l_6+2l_7+l_8=0,
\qquad
l_9=l_6+l_7,
\qquad
l_{10}=-2l_6,
\label{eq:grouped-six-rel-1}
\end{equation}
\begin{equation}
l_{11}+l_{12}=2l_6,
\qquad
l_{13}=l_6,
\qquad
l_{14}=-\frac16 l_{11}.
\label{eq:grouped-six-rel-2}
\end{equation}
Every minimally demoted local grouped-$P_2$ image must satisfy all six
simultaneously.  The exact GR middle block does not.  The referee audit checks
that each obstruction polynomial is generically nonzero once the solved target
\((\Delta l_6,\dots,\Delta l_{14})\) is substituted.

The isotropic specialization is even more rigid.  Setting
\(a_2=b_2=0\) gives
\begin{equation}
\begin{aligned}
l_6&=u_2,
& l_7&=-u_2,
& l_8&=l_9=0,
& l_{10}&=-2u_2,
\\
 l_{11}&=4u_2,
& l_{12}&=-2u_2,
& l_{13}&=u_2,
& l_{14}&=-\frac23u_2.
\end{aligned}
\label{eq:grouped-isotropic-minimal-image}
\end{equation}
So the first moral of the section is negative but precise: the grouped real
$P_2$ ontology is not ruled out by $3$PN, but its smallest local constitutive
lift is too rigid to carry the exact middle block.

\subsection{Richer local grouped-\texorpdfstring{$P_2$}{P 2} family}
\label{sec:grouped-richer-family}

The failure of the minimal lift does not mean that the grouped ontology is too
small.  It means only that the first constitutive guess was too poor.  The next
natural local grouped families are obtained by introducing the source-square
blocks
\begin{equation}
C_{20}^2=\frac16(3d^2-v^2-2U)^2,
\qquad
C_{21}^2=2d^2u_\perp^2,
\qquad
C_{22}^2=\frac12u_\perp^4,
\label{eq:grouped-C-squares}
\end{equation}
and then dressing them in the three simplest local ways compatible with the
$3$PN middle block,
\begin{equation}
T_{20}=\frac13Ud^2(3u_\perp^2-U)^2,
\qquad
T_{21}=Uu_\perp^2(u_\perp^2-d^2-U)^2,
\qquad
T_{22}=Ud^2u_\perp^4,
\label{eq:grouped-T-family}
\end{equation}
\begin{equation}
S_A=U^2C_A^2,
\qquad
V_A=\frac{v^2}{U}S_A=Uv^2C_A^2,
\qquad
A\in\{20,21,22\}.
\label{eq:grouped-SV-family}
\end{equation}
The three minimally demoted dynamic norms \((T_A)\) still span only rank three,
and appending the static-support family \((S_A)\) raises the rank only to six.
But the full nine-element family
\[
(T_{20},T_{21},T_{22},S_{20},S_{21},S_{22},V_{20},V_{21},V_{22})
\]
acts on the nine-slot middle block by an exact $9\times 9$ compiler matrix with
constant determinant
\begin{equation}
\boxed{\det M_{\rm mid}=-\frac4{27}\neq 0.}
\label{eq:grouped-mid-det}
\end{equation}
So the richer grouped family is exactly invertible on the whole middle block.
This is the decisive positive grouped theorem of the $3$PN conservative audit:
the grouped real $P_2$ ontology itself is already large enough to carry the full
nine-slot middle block once the constitutive lift is enlarged in the smallest
natural way.

It is therefore natural to define the richer grouped module by
\begin{equation}
L_{P_2}^{\rm mid}
=
\sum_{A\in\{20,21,22\}}
\bigl(\lambda_A T_A+\sigma_A S_A+\tau_A V_A\bigr).
\label{eq:grouped-mid-ansatz}
\end{equation}
Because of \eqref{eq:grouped-mid-det}, every exact middle-block target has one
and only one coefficient triple
\((\lambda_A,\sigma_A,\tau_A)\).

\subsection{Exact GR grouped coefficients}
\label{sec:grouped-gr-coeffs}

Substituting the solved GR middle block
\((L_6,\dots,L_{14})=(\Delta l_6,\dots,\Delta l_{14})\)
into the exact inverse compiler yields the unique grouped coefficient set
\begin{align}
\lambda_{20}^{\rm GR}
&=\frac{\nu\,(-276\nu^2-3154\nu+3\pi^2+2672)}{32},
\nonumber\\
\lambda_{21}^{\rm GR}
&=\frac{\nu\,(2568\nu^2+2236\nu-3044-3\pi^2)}{192},
\nonumber\\
\lambda_{22}^{\rm GR}
&=\frac{\nu\,(7650\nu^2+32858\nu-30136-33\pi^2)}{96},
\label{eq:grouped-lambda-gr}
\end{align}
\begin{align}
\sigma_{20}^{\rm GR}
&=\frac{\nu\,(-102\nu^2-308\nu+293)}{16},
\nonumber\\
\sigma_{21}^{\rm GR}
&=\frac{\nu\,(-4188\nu^2-23778\nu+21\pi^2+22006)}{192},
\nonumber\\
\sigma_{22}^{\rm GR}
&=\frac{\nu\,(3906\nu^2+2478\nu-2791-3\pi^2)}{48},
\label{eq:grouped-sigma-gr}
\end{align}
\begin{align}
\tau_{20}^{\rm GR}
&=\frac{3\nu\,(-154\nu^2-87\nu+38)}{32},
\nonumber\\
\tau_{21}^{\rm GR}
&=\frac{\nu\,(-6672\nu^2-5824\nu+3\pi^2+4412)}{384},
\nonumber\\
\tau_{22}^{\rm GR}
&=\frac{\nu\,(-2790\nu^2-3367\nu+3\pi^2+3386)}{96}.
\label{eq:grouped-tau-gr}
\end{align}
With exactly these coefficients, the richer grouped lift
\eqref{eq:grouped-mid-ansatz} reproduces the entire solved GR middle block
identically:
\begin{equation}
(l_6,\dots,l_{14})=(\Delta l_6,\dots,\Delta l_{14}).
\label{eq:grouped-exact-middle-fit}
\end{equation}
So the conservative $3$PN answer to the grouped front-end question is now sharp:
\begin{equation}
\boxed{\text{grouped real }P_2\text{ can close the full $3$PN middle block exactly.}}
\label{eq:grouped-main-theorem}
\end{equation}
This is the first place in the program where the conservative quadrupole bundle
has survived a fully overdetermined exact target rather than a qualitative
consistency test.

\subsection{Structural corollaries}
\label{sec:grouped-corollaries}

Two corollaries are immediate and both matter for the rest of the paper.

First, every basis element in \((T_A,S_A,V_A)\) carries at least one power of
\(U\).  Therefore the richer grouped module forces
\begin{equation}
l_1=l_2=l_3=l_4=l_5=0
\label{eq:grouped-no-pure-kinetic}
\end{equation}
identically.  So the grouped real $P_2$ closure does \emph{not} compete with the
pure-kinetic sector.  The free-particle residual remains outside the grouped
module and will later be shown to be only the ordinary-Lagrangian counterimage
of the universal free COM Hamiltonian.

Second, the richer grouped compiler predicts a definite static companion rather
than leaving the quartic potential slot arbitrary.  One finds
\begin{equation}
l_{15}^{(P_2\,\rm pred)}=\frac23\sigma_{20}
=l_{10}+l_{11}+l_{12}+2(l_6+l_7+l_8+l_9).
\label{eq:grouped-static-companion}
\end{equation}
For the exact GR middle block this becomes
\begin{equation}
\Delta l_{15}^{(P_2\,\rm pred)}
=\frac{\nu(293-308\nu-102\nu^2)}{24}.
\label{eq:grouped-static-prediction}
\end{equation}
Comparing with the true static residual isolated in the exact COM target gives the leftover scalar gap
\begin{equation}
\Delta l_{15}^{(0/g)}
\equiv
\Delta l_{15}^{\rm GR}-\Delta l_{15}^{(P_2\,\rm pred)}
=\frac{\nu(408\nu^2+1232\nu-2080+63\pi^2)}{96}.
\label{eq:grouped-static-gap}
\end{equation}
So once the richer grouped real $P_2$ compiler closes the nine middle slots
exactly, the full solved $3$PN COM residual already splits as
\begin{equation}
\Delta \hat L_3^{\rm GR}
=
\Delta l_1 v^8
+
L_{P_2}^{\rm mid}[\lambda_A^{\rm GR},\sigma_A^{\rm GR},\tau_A^{\rm GR}]
+
\Delta l_{15}^{(0/g)}U^4.
\label{eq:grouped-final-split}
\end{equation}
That is the exact point where the old scalar/geometry lane re-enters.  The
richer grouped module has done everything it can do: it has carried the whole
middle block and fixed its unavoidable static shadow.  The only remaining
conservative scalar task is the one-dimensional static remainder, which is the
subject of the next section.

% Section content for 4d_3pn.tex

\section{Static scalar/geometry completion}
\label{sec:geometry}

The preceding section closed the entire nine-slot middle interaction block by the
richer grouped real $P_2$ compiler.  What remains at conservative $3$PN is
therefore not a diffuse family of leftover coefficients but a sharply isolated
one-dimensional static remainder.  The main point of the present section is that
this scalar remainder is more rigid than it first appears.  Inside the old
constant-coefficient generic-frame scaffold, the static block is too small to
carry the exact gap.  If one temporarily enlarges the scalar lane so as to allow
both a body-local $P_0$ carrier and a pair/geometry carrier, the center-of-mass
projection seems to leave a one-parameter family.  But that family is
algebraically redundant: an exact mass identity collapses it to a unique
pair/geometry representative.  In the fixed ADM chart, the Hamiltonian image of
that completion is then obtained by the same exact sign flip that already
governs the rest of the residual.  After this section, the static conservative
$3$PN lane is completely fixed; only its microscopic moving-throat realization
remains open.

\subsection{Residual static slot after grouped closure}
\label{sec:geometry-residual}

Once the richer grouped real $P_2$ compiler has reproduced the exact nine-slot
middle block, the solved COM residual splits as
\begin{equation}
\Delta \hat L_3^{\rm GR}
=
\Delta l_1 v^8
+
L_{P_2}^{\rm mid}[\lambda_A^{\rm GR},\sigma_A^{\rm GR},\tau_A^{\rm GR}]
+
\Delta l_{15}^{(0/g)} U^4,
\qquad
U\equiv \frac{GM}{r}.
\label{eq:geometry-static-split}
\end{equation}
So the grouped module has already done all it can do.  The only remaining scalar
job is to identify the carrier of the one-dimensional static remainder.

A useful way to display that remainder is to compare the static companion forced
by the grouped middle-block fit with the exact GR target.  The grouped compiler
predicts
\begin{equation}
l_{15}^{(P_2\,\mathrm{pred})}
=
\frac{\nu\,(-102\nu^2-308\nu+293)}{24},
\label{eq:geometry-static-pred}
\end{equation}
whereas the exact GR coefficient is
\begin{equation}
l_{15}^{\rm GR}
=
\frac{\nu\,(-908+63\pi^2)}{96}.
\label{eq:geometry-static-target}
\end{equation}
The static gap is therefore
\begin{equation}
\boxed{
\Delta l_{15}^{(0/g)}
=
 l_{15}^{\rm GR}-l_{15}^{(P_2\,\mathrm{pred})}
=
\frac{\nu\,(408\nu^2+1232\nu-2080+63\pi^2)}{96}.
}
\label{eq:geometry-static-gap}
\end{equation}
Equation \eqref{eq:geometry-static-gap} is the full conservative scalar datum
left after the grouped real $P_2$ module succeeds.

\subsection{\texorpdfstring{$U$}{U}-block no-go inside the old constant-coefficient generic scaffold}
\label{sec:geometry-u-block}

The first exact result is negative.  In the old constant-coefficient generic-frame
interaction scaffold, the only static basis polynomial is
\begin{equation}
\mathcal U_1=p^2q+pq^2,
\qquad p\equiv m_A,
\qquad q\equiv m_B,
\label{eq:geometry-U1}
\end{equation}
inside the common prefactor $G^4pq/r^4$.  Writing
\begin{equation}
M=p+q,
\qquad
\mu=\frac{pq}{M},
\qquad
\nu=\frac{pq}{M^2},
\qquad
U=\frac{GM}{r},
\label{eq:geometry-mass-dictionary}
\end{equation}
one finds the exact COM image
\begin{equation}
\frac{1}{\mu}\,\frac{G^4pq}{r^4}\,\mathcal U_1
=
\nu U^4.
\label{eq:geometry-U1-com-image}
\end{equation}
So every constant-coefficient static interaction term in the old scaffold has a
COM coefficient of the form
\begin{equation}
l_{15}=\nu\,u_1,
\label{eq:geometry-linear-in-nu}
\end{equation}
linear in $\nu$.

But the exact remainder \eqref{eq:geometry-static-gap} still contains $\nu^2$ and
$\nu^3$ after the overall factor of $\nu$ is extracted.  Therefore no new
constant coordinate inside the original $24/17/8/1$ generic-frame scaffold can
carry the actual $3$PN static gap.  The residual scalar lane must instead be
represented as a $\nu$-dressed counterterm family.

\subsection{Apparent \texorpdfstring{$P_0$}{P 0} versus geometry ambiguity}
\label{sec:geometry-ambiguity}

Once one allows such $\nu$-dressed scalar carriers, two natural static mass
families appear.  The first is a body-local scalar family,
\begin{equation}
\mathcal U_0=p^3+q^3,
\label{eq:geometry-U0}
\end{equation}
which has the right bookkeeping character for a body-local $P_0$-type wall or
support dressing.  The second is the intrinsic pair family
\begin{equation}
\mathcal U_g=p^2q+pq^2,
\label{eq:geometry-Ug}
\end{equation}
which is the natural candidate for a geometry/compressibility completion.

At COM level these two families span only one scalar observable.  Indeed,
\begin{equation}
\frac{1}{\mu}\,\frac{G^4pq}{r^4}\,\nu\mathcal U_0
=
\nu(1-3\nu)U^4,
\qquad
\frac{1}{\mu}\,\frac{G^4pq}{r^4}\,\mathcal U_g
=
\nu U^4.
\label{eq:geometry-U0-Ug-images}
\end{equation}
So before the exact mass algebra is used, the most natural generic-frame static
counterterm looks like a regular one-parameter $P_0\oplus g$ family,
\begin{equation}
\Delta L_{3,\rm ct}^{(0)}
=
\frac{G^4pq}{r^4}\,\sigma\nu\,\mathcal U_0,
\label{eq:geometry-sigma-P0}
\end{equation}
\begin{equation}
\Delta L_{3,\rm ct}^{(g)}
=
\frac{G^4pq}{r^4}
\left[
\frac{408\nu^2+1232\nu-2080+63\pi^2}{96}
-\sigma(1-3\nu)
\right]\mathcal U_g.
\label{eq:geometry-sigma-g}
\end{equation}
Here $\sigma$ labels how much of the COM scalar lane is assigned to a body-local
$P_0$ carrier rather than to the pair/geometry carrier.  At this stage the split
looks real, because COM fixes only the sum.

\subsection{Sigma-collapse and unique geometry completion}
\label{sec:geometry-sigma-collapse}

The apparent $\sigma$-freedom is illusory.  The exact mass identity
\begin{equation}
\boxed{
\nu\,(p^3+q^3)=(1-3\nu)(p^2q+pq^2)
}
\label{eq:geometry-mass-identity}
\end{equation}
collapses the full family
\eqref{eq:geometry-sigma-P0}--\eqref{eq:geometry-sigma-g} identically.
Equivalently,
\begin{equation}
\frac{d}{d\sigma}
\Bigl(
\Delta L_{3,\rm ct}^{(0)}+\Delta L_{3,\rm ct}^{(g)}
\Bigr)=0.
\label{eq:geometry-dsigma-zero}
\end{equation}
So the full static scalar counterterm reduces to one unique pair/geometry
representative,
\begin{equation}
\boxed{
\Delta L_{3,\rm ct}^{\rm scalar}
=
\frac{G^4pq}{96r^4}
\bigl(408\nu^2+1232\nu-2080+63\pi^2\bigr)
\,(p^2q+pq^2).
}
\label{eq:geometry-unique-counterterm}
\end{equation}
Its COM image is precisely the residual coefficient
\eqref{eq:geometry-static-gap}:
\begin{equation}
\frac{1}{\mu}\,\Delta L_{3,\rm ct}^{\rm scalar}
=
\Delta l_{15}^{(g)}U^4,
\qquad
\Delta l_{15}^{(g)}
=
\frac{\nu\,(408\nu^2+1232\nu-2080+63\pi^2)}{96}.
\label{eq:geometry-unique-com-image}
\end{equation}
It is also convenient to phrase the same result directly at the reduced-coefficient
level:
\begin{equation}
\begin{aligned}
 c_{\rm target}
 &=-\frac{227}{24}+\frac{21\pi^2}{32},
 &
 c_{\rm pred}^{(P_2)}
 &=\frac{293}{24}-\frac{77}{6}\nu-\frac{17}{4}\nu^2,
 \\
 c_{\rm geom}
 &=\frac{408\nu^2+1232\nu-2080+63\pi^2}{96},
\end{aligned}
\label{eq:geometry-c-target-pred-gap}
\end{equation}
with
\begin{equation}
 c_{\rm target}=c_{\rm pred}^{(P_2)}+c_{\rm geom}.
\label{eq:geometry-target-decomposition}
\end{equation}
So once the exact mass algebra is respected, there is no independent body-local
$P_0$ dressing left.  The leftover $3$PN static scalar sector belongs uniquely
to the geometry side.

\subsection{Hamiltonian image of the static completion}
\label{sec:geometry-hamiltonian-image}

The fixed-chart generic-frame compiler established earlier acts on the residual
by exact sign flip,
\begin{equation}
\Delta H_3=-\Delta L_3.
\label{eq:geometry-sign-flip}
\end{equation}
So the Hamiltonian image of the unique static completion follows immediately:
\begin{equation}
\boxed{
\Delta H_{3,\rm ct}^{\rm scalar}
=
-\Delta L_{3,\rm ct}^{\rm scalar}.
}
\label{eq:geometry-hamiltonian-image-eq}
\end{equation}
No further algebraic resolution is needed.  The ordinary and Hamiltonian static
remainders are already the same unique polynomial with opposite sign, and the
only remaining question is microscopic rather than algebraic: what scalar-side
wall or geometry dynamics on the moving-throat branch realizes the coefficient in
\eqref{eq:geometry-unique-counterterm}?

At this stage, therefore, the full conservative $3$PN static lane is fixed.  The
only untreated COM slot is the apparent pure-kinetic $v^8$ remainder, and that
slot is shown in the next section to be nothing more than the ordinary-Lagrangian
counterimage of the universal free two-body relativistic Hamiltonian.

% Section content for 4d_3pn.tex

\section{Pure-kinetic sector and the free two-body Hamiltonian}
\label{sec:kinetic}

The preceding section showed that the last genuine scalar ambiguity collapses to a
unique geometry-side completion.  So after the grouped real $P_2$ middle block
and the static scalar lane are both fixed, only one apparently separate
center-of-mass pure-kinetic slot remains.  The point of the present section is
that this leftover is not a new throat-response datum.  At generic-frame level,
the exact free relativistic two-body Lagrangian contains no mixed comparable-mass
$3$PN pure-kinetic interaction monomial at all.  The apparent center-of-mass
residual appears only because the exact ordinary/Hamiltonian compiler maps the
universal free relativistic two-body Hamiltonian to a shifted ordinary
coefficient.  Once that compiler compensation is taken into account, the entire
pure-kinetic $3$PN sector closes identically.  So the conservative $3$PN split
really does stop with three lanes only: the compiler image of free kinematics,
the exact grouped real $P_2$ middle block, and the unique geometry-side static
completion.

\subsection{Why the apparent pure-kinetic residual is not a new throat datum}
\label{sec:kinetic-why}

After the exact grouped middle-block closure and the static scalar/geometry
completion, the only seemingly independent center-of-mass kinetic remainder is
\begin{equation}
\boxed{
\Delta l_1=
\frac{3\nu(3\nu-1)(4\nu-1)}{16},
\qquad
\Delta l_2=\Delta l_3=\Delta l_4=\Delta l_5=0.
}
\label{eq:kinetic-delta-l1-only}
\end{equation}
In the standard isotropic $15$-slot center-of-mass basis, this is the entire
leftover pure-kinetic block.  Taken in isolation, \eqref{eq:kinetic-delta-l1-only}
looks as though one more independent $3$PN kinetic throat coefficient must still
be supplied.

That interpretation is too quick.  A genuinely new kinetic datum would have to
appear first as a new generic-frame interaction family.  But the exact free-body
audit shows that no such family exists.  The whole residual arises instead from
the exact ordinary/Hamiltonian compiler applied to the already-fixed universal
free relativistic two-body Hamiltonian.  So the right question is not ``what new
kinetic interaction must the throat supply?'' but rather ``what ordinary-side
coefficient is required so that the fixed free Hamiltonian target is reproduced
exactly in center-of-mass variables?''

\subsection{Generic-frame free kinematics}
\label{sec:kinetic-generic}

At generic-frame level, the exact free ordinary relativistic two-body Lagrangian
is simply the sum of the two separate free-body square roots,
\begin{equation}
L_{\rm free}^{\rm gen}
=
-m_Ac^2\sqrt{1-\frac{v_A^2}{c^2}}
-
 m_Bc^2\sqrt{1-\frac{v_B^2}{c^2}}.
\label{eq:kinetic-free-gen-exact}
\end{equation}
Expanding through $3$PN gives
\begin{equation}
\begin{aligned}
L_{\rm free}^{\rm gen}
=
&-(m_A+m_B)c^2
+\frac12\bigl(m_Av_A^2+m_Bv_B^2\bigr)
+\frac{1}{8c^2}\bigl(m_Av_A^4+m_Bv_B^4\bigr)
\\
&+\frac{1}{16c^4}\bigl(m_Av_A^6+m_Bv_B^6\bigr)
+\frac{5}{128c^6}\bigl(m_Av_A^8+m_Bv_B^8\bigr)
+O(c^{-8}).
\end{aligned}
\label{eq:kinetic-free-gen-series}
\end{equation}
So the $3$PN pure-kinetic block is already exactly
\begin{equation}
\boxed{
L_{3,\rm free}^{\rm gen}
=
\frac{5}{128}\bigl(m_Av_A^8+m_Bv_B^8\bigr).
}
\label{eq:kinetic-free-gen-3pn}
\end{equation}
There is no mixed comparable-mass pure-kinetic interaction monomial here.  In
particular, the generic-frame free block contains no independent $m_Am_B$-type
$octic$ interaction that could serve as a new $3$PN constitutive datum.

If one now performs the naive center-of-mass reduction
\begin{equation}
\mathbf v_A=\frac{m_B}{M}\,\mathbf v,
\qquad
\mathbf v_B=-\frac{m_A}{M}\,\mathbf v,
\qquad
M\equiv m_A+m_B,
\label{eq:kinetic-com-velocities}
\end{equation}
with symmetric mass ratio
\begin{equation}
\nu\equiv \frac{m_Am_B}{M^2},
\label{eq:kinetic-nu-def}
\end{equation}
then the carried free ordinary seed coefficient becomes
\begin{equation}
\boxed{
 l_1^{\rm seed}
 =
 \frac{5}{128}-\frac{35}{128}\nu+\frac{35}{64}\nu^2-\frac{35}{128}\nu^3.
}
\label{eq:kinetic-l1-seed}
\end{equation}
This is exactly the coefficient that one obtains by reducing the separate-body
free kinematics before the exact center-of-mass compiler is imposed.

\subsection{Exact free COM Hamiltonian target}
\label{sec:kinetic-com-target}

The correct reference object is instead the exact free relativistic two-body
center-of-mass Hamiltonian,
\begin{equation}
H_{\rm free}^{\rm COM}
=
\sqrt{m_A^2c^4+P^2c^2}
+
\sqrt{m_B^2c^4+P^2c^2},
\label{eq:kinetic-free-com-h-exact}
\end{equation}
with reduced momentum variable
\begin{equation}
P=\mu p,
\qquad
\mu\equiv \frac{m_Am_B}{M}.
\label{eq:kinetic-P-mu-p}
\end{equation}
Subtracting the rest energy and dividing by $\mu$ gives the reduced expansion
\begin{equation}
\frac{H_{\rm free}^{\rm COM}-Mc^2}{\mu}
=
\frac12 p^2
-
\frac{1-3\nu}{8c^2}p^4
+
\frac{1-5\nu+5\nu^2}{16c^4}p^6
+
\frac{h_1^{\rm free}}{c^6}p^8
+O(c^{-8}),
\label{eq:kinetic-free-com-h-series}
\end{equation}
where the exact $3$PN free coefficient is
\begin{equation}
\boxed{
 h_1^{\rm free}
 =
 \frac{-5+35\nu-70\nu^2+35\nu^3}{128},
 \qquad
 h_2=h_3=h_4=h_5=0.
}
\label{eq:kinetic-h1-free}
\end{equation}
This coefficient is not a new observable of the throat.  It is simply the
universal free relativistic two-body Hamiltonian written in center-of-mass
variables.  Accordingly, the remaining kinetic question is entirely a compiler
question: what ordinary coefficient corresponds to the already-fixed Hamiltonian
slot \eqref{eq:kinetic-h1-free}?

\subsection{Collapse of the ordinary pure-kinetic residual}
\label{sec:kinetic-collapse}

The exact center-of-mass compiler established earlier acts on the first kinetic
slot by
\begin{equation}
\boxed{
 l_1=F_1(\nu)-h_1,
 \qquad
 F_1(\nu)=\frac{3}{16}\nu-\frac{21}{16}\nu^2+\frac94\nu^3.
}
\label{eq:kinetic-compiler-l1}
\end{equation}
Applying \eqref{eq:kinetic-compiler-l1} to the naive free ordinary seed
\eqref{eq:kinetic-l1-seed} gives the Hamiltonian image of that seed,
\begin{equation}
\boxed{
 h_1^{\rm seed}
 =
 F_1(\nu)-l_1^{\rm seed}
 =
 -\frac{5}{128}+\frac{59}{128}\nu-\frac{119}{64}\nu^2+\frac{323}{128}\nu^3.
}
\label{eq:kinetic-h1-seed}
\end{equation}
The entire pure-kinetic remainder is therefore fixed by the difference between
this seed Hamiltonian image and the exact free Hamiltonian target:
\begin{equation}
\begin{aligned}
\Delta l_1
&=
 h_1^{\rm seed}-h_1^{\rm free}
\\
&=
\frac{3}{16}\nu-\frac{21}{16}\nu^2+\frac94\nu^3
=
\boxed{\frac{3\nu(3\nu-1)(4\nu-1)}{16}}.
\end{aligned}
\label{eq:kinetic-delta-l1-collapse}
\end{equation}
This is exactly the leftover coefficient displayed in
\eqref{eq:kinetic-delta-l1-only}.  Equivalently, the full ordinary target is
\begin{equation}
 l_1^{\rm GR}
 =
 F_1(\nu)-h_1^{\rm free}
 =
 \frac{5-11\nu-98\nu^2+253\nu^3}{128},
\label{eq:kinetic-l1-target}
\end{equation}
with the shift from \eqref{eq:kinetic-l1-seed} to
\eqref{eq:kinetic-l1-target} given entirely by the compiler compensation
\eqref{eq:kinetic-delta-l1-collapse}.

So the apparent pure-kinetic leftover is not a new constitutive $3$PN throat
coefficient.  It is only the ordinary-Lagrangian counterimage of the universal
free relativistic two-body center-of-mass Hamiltonian.  Since the free target
also has
\begin{equation}
 h_2=h_3=h_4=h_5=0,
\label{eq:kinetic-higher-free-zero}
\end{equation}
no extra generic-frame pure-kinetic interaction module should be added in the
fixed ADM chart.  The grouped real $P_2$ module already closes the whole middle
block, the previous section already fixes the static scalar lane uniquely on the
geometry side, and the present computation shows that the only remaining octic
center-of-mass term is generated automatically by the exact compiler.  After
this collapse, no separate algebraic conservative $3$PN bottleneck remains.

% Section content for 4d_3pn.tex

\section{Final conservative \texorpdfstring{$3$}{3}PN theorem}
\label{sec:final-theorem}

The preceding three sections have removed all remaining algebraic uncertainty from
the conservative $3$PN problem.  The nine-slot middle interaction block is
carried exactly by the richer grouped real $P_2$ module; the apparent static
$P_0/g$ ambiguity collapses to a unique geometry-side completion; and the
remaining center-of-mass pure-kinetic slot is only the ordinary-Lagrangian image
of the universal free relativistic two-body Hamiltonian.  The final theorem is
therefore much shorter than the derivation chain that leads to it.  In the fixed
ADM chart, the full comparable-mass $c^{-6}$ coefficient is completely assigned,
and what remains open is no longer algebraic coefficient matching but the deeper
moving-throat derivation of the grouped real $P_2$ constitutive law, the
geometry completion, and then the passive/outgoing quadrupole normalization that
enters the later $2.5$PN bridge.

\subsection{Final exact split in COM form}
\label{sec:final-com}

Within the declared hierarchy, the exact comparable-mass $3$PN coefficient in the
isotropic COM chart is
\begin{equation}
\boxed{
\Delta \hat L_3^{\rm GR}
=
\underbrace{\Delta l_1 v^8}_{\text{free-COM compiler image}}
+
\underbrace{L_{P_2}^{\rm mid}}_{\text{exact grouped real }P_2\text{ closure}}
+
\underbrace{\Delta l_{15}^{(g)}U^4}_{\text{unique geometry completion}}.
}
\label{eq:final-com-split}
\end{equation}
Here \(U=GM/r\) is the usual COM Newtonian potential.
The two scalar coefficients worth keeping explicit are
\begin{equation}
\boxed{
\Delta l_1=\frac{3\nu(3\nu-1)(4\nu-1)}{16},
\qquad
\Delta l_{15}^{(g)}
=
\frac{\nu(408\nu^2+1232\nu-2080+63\pi^2)}{96}.
}
\label{eq:final-com-scalars}
\end{equation}
The meaning of \eqref{eq:final-com-split} is now completely sharp.  The term
$\Delta l_1v^8$ is not a new throat-response coefficient; it is exactly the
ordinary-side compensation required by the free relativistic two-body COM
Hamiltonian.  The grouped module
\begin{equation}
L_{P_2}^{\rm mid}
=
\sum_{A\in\{20,21,22\}}
\bigl(
\lambda_A^{\rm GR}T_A+\sigma_A^{\rm GR}S_A+\tau_A^{\rm GR}V_A
\bigr)
\label{eq:final-com-grouped-module}
\end{equation}
reproduces the full solved nine-slot middle block
\[
(\Delta l_6,\ldots,\Delta l_{14})
\]
exactly, with the coefficient triples
\((\lambda_A^{\rm GR},\sigma_A^{\rm GR},\tau_A^{\rm GR})\) fixed uniquely by the
GR target.  Finally, the static term
\(\Delta l_{15}^{(g)}U^4\) is not a free $P_0/g$ admixture.  After sigma-collapse
it is the unique geometry-side completion of the remaining scalar slot.

So the COM theorem really does stop with three and only three lanes:

\begin{enumerate}
  \item a compiler image of universal free kinematics;
  \item an exact grouped real $P_2$ middle-block closure;
  \item a unique geometry-side static completion.
\end{enumerate}

There is no fourth independent conservative $3$PN carrier hiding behind the
center-of-mass bookkeeping.  Every solved coefficient belongs to one of the
three structures displayed in \eqref{eq:final-com-split}.

\subsection{Final exact split in generic-frame form}
\label{sec:final-generic}

The generic-frame statement is not a second independent theorem; it is the unique
Hamiltonian-first lift of the same three-lane solution.  Let
\[
p\equiv m_A,
\qquad
q\equiv m_B,
\qquad
\nu\equiv\frac{pq}{(p+q)^2}.
\]
Then the full generic-frame ordinary $3$PN coefficient can be written as
\begin{equation}
\boxed{
L_3^{\rm gen}
=
L_{3,\rm seed}
+
\sum_{A\in\{20,21,22\}}
\bigl(
\lambda_A^{\rm GR}T_A+\sigma_A^{\rm GR}S_A+\tau_A^{\rm GR}V_A
\bigr)
+
\Delta L_{3,\rm ct}^{\rm scalar},
}
\label{eq:final-generic-split}
\end{equation}
where the natural self/static seed is the symmetric promotion of the exact
Schwarzschild one-body target, the grouped real $P_2$ families
\((T_A,S_A,V_A)\) carry the full middle interaction block, and the unique static
counterterm is
\begin{equation}
\boxed{
\Delta L_{3,\rm ct}^{\rm scalar}
=
\frac{G^4pq}{96r^4}
\bigl(408\nu^2+1232\nu-2080+63\pi^2\bigr)(p^2q+pq^2).
}
\label{eq:final-generic-static}
\end{equation}

Two features of \eqref{eq:final-generic-split} should be emphasized.  First, the
free-body part of the seed remains exactly
\[
L_{3,\rm free}^{\rm gen}
=
\frac{5}{128}\bigl(m_Av_A^8+m_Bv_B^8\bigr),
\]
so there is no separate mixed comparable-mass generic-frame pure-kinetic
interaction family to be added by hand.  The COM term \(\Delta l_1v^8\) appears
only after the exact compiler maps the universal free Hamiltonian back to the
ordinary chart.  Second, once the lower-order ledger is frozen, the fixed ADM
chart obeys the exact generic-frame relation
\begin{equation}
\boxed{\Delta H_3=-\Delta L_3.}
\label{eq:final-sign-flip}
\end{equation}
So the earlier COM-blind null directions and the smaller ordinary contact/gauge
orbit are intermediate bookkeeping structures only.  They do not survive as true
freedom after the full generic-frame Hamiltonian target is imposed.

In this sense the generic-frame theorem is stronger than a COM target match.
There is a unique ordinary representative in the fixed chart, and it is exactly
the Hamiltonian-first lift of the three solved pieces already displayed in the
COM decomposition.

\subsection{Equality to the standard \texorpdfstring{$3$}{3}PN ADM target}
\label{sec:final-adm}

Because the generic-frame Hamiltonian compiler has zero kernel in the fixed ADM
chart, the assembled ordinary representative Legendre-transforms to the standard
$3$PN ADM-type Hamiltonian target in exactly the same sense that the full
conservative $1$PN paper closes to EIH and the full conservative $2$PN paper
closes to ADM.  After the one-body repair ledger, the grouped middle-block
closure, the unique geometry completion, and the pure-kinetic compiler image are
all inserted, no residual Hamiltonian discrepancy remains:
\begin{equation}
\boxed{
H_3^{\rm derived}=H_3^{\rm ADM},
\qquad
\Delta H_3
\equiv
H_3^{\rm derived}-H_3^{\rm ADM}=0.
}
\label{eq:final-adm-equality}
\end{equation}

The role of the Hamiltonian-first route should be stated explicitly.  A naive
ordinary-Lagrangian COM reduction of the imported generic-frame target does not
reproduce the exact COM ordinary target.  By contrast, the sequence
\[
\text{generic-frame ordinary target}
\;\longrightarrow\;
\text{exact cubic Legendre lift}
\;\longrightarrow\;
\text{COM reduction}
\]
does reproduce the exact COM Hamiltonian and therefore fixes the correct generic
representative.  Once that route is used, the apparent $27$-dimensional
COM-blind family and the smaller $11$-generator ordinary contact orbit are seen
for what they are: bookkeeping artifacts of an incomplete chart choice rather
than surviving physical or algebraic freedom.

So the conservative $3$PN result should not be read as a target match obtained
after a sequence of special chart adjustments.  It is an exact final equality in
the fixed ADM chart once the correct compiler route is respected.

\subsection{Main theorem statement}
\label{sec:final-main}

\paragraph{Main theorem.}
Fix the carried Newtonian, $1$PN, and $2$PN conservative ledger and work within
the same declared closure hierarchy used in the earlier full conservative
papers.  Let the exact one-body $3$PN Schwarzschild gate be repaired by
\begin{equation}
\mu_{\rho,3}=\frac14,
\qquad
 d_3=-\frac{45}{4},
\qquad
 s_{24}=-\frac1{16},
\label{eq:final-onebody-ledger}
\end{equation}
let the grouped real $P_2$ middle-block compiler be the richer
\((T_A,S_A,V_A)\) family with
\begin{equation}
\det M_{\rm mid}=-\frac4{27},
\label{eq:final-mid-det}
\end{equation}
and let the fixed ADM chart be imposed by the exact generic-frame sign-flip
identity \eqref{eq:final-sign-flip}.  Then the full conservative two-body $3$PN
sector is uniquely fixed.  Equivalently,
\begin{equation}
\boxed{
L_3^{\rm gen}
=
L_{3,\rm seed}
+
L_{P_2}^{\rm mid}[\lambda_A^{\rm GR},\sigma_A^{\rm GR},\tau_A^{\rm GR}]
+
\Delta L_{3,\rm ct}^{\rm scalar},
}
\label{eq:final-main-theorem-generic}
\end{equation}
with COM residual
\begin{equation}
\boxed{
\Delta \hat L_3^{\rm GR}
=
\Delta l_1v^8
+
L_{P_2}^{\rm mid}
+
\Delta l_{15}^{(g)}U^4,
}
\label{eq:final-main-theorem-com}
\end{equation}
and Hamiltonian equality
\begin{equation}
\boxed{
H_3^{\rm derived}=H_3^{\rm ADM}.
}
\label{eq:final-main-theorem-h}
\end{equation}

The point of the theorem is not merely that the model can reproduce the $3$PN
target after enough algebra.  The stronger conclusion is structural: the
conservative $3$PN bottleneck is no longer an algebraic coefficient search.  The
higher-order conservative burden has now been compressed to an exact grouped real
$P_2$ module plus a unique geometry-side completion, while the apparent
pure-kinetic remainder is absorbed completely by the exact free-Hamiltonian
compiler image.

This should also be read with the same status discipline as the earlier full
conservative papers.  The theorem is exact \emph{within the declared hierarchy};
it is not being claimed as an assumption-free theorem of a fully solved
moving-throat PDE.  What remains open is therefore not another conservative $3$PN
algebra problem.  What remains open is the deeper moving-throat derivation of the
grouped real $P_2$ constitutive law and the geometry completion, followed by
their insertion into the already narrowed passive/outgoing quadrupole
normalization bridge.  That interface is the subject of the next section.

% Section content for 4d_3pn.tex

\section{Interface to the \texorpdfstring{$2.5$}{2.5}PN program}
\label{sec:interface-25pn}

The conservative $3$PN theorem does not replace the $2.5$PN program; it tells us
exactly how that program should now be read.  The $2.5$PN audit had already
reduced the universal higher-order route to the orbital/\allowbreak worldtube STF
quadrupole, packaged conservatively as the grouped real
\[
P_{20}\oplus P_{21}\oplus P_{22}
\]
bundle.  What remained unclear at that stage was whether the same grouped
ontology could survive its first honest conservative stress test, or whether the
future outgoing-normalization problem would have to compete with some different
hidden conservative carrier.  The present $3$PN result answers that question.
The full nine-slot middle interaction block is carried exactly by the richer
local grouped-$P_2$ module, the residual static slot is uniquely geometry-side,
and the apparent COM pure-kinetic leftover is only the ordinary-Lagrangian image
of the universal free relativistic two-body Hamiltonian.  So the $2.5$PN bridge
should now be read against a much sharper conservative background: the remaining
gap is no longer a vague higher-order bookkeeping problem, but the passive/
outgoing normalization of one already identified grouped-quadrupole branch.

\subsection{Why \texorpdfstring{$3$}{3}PN is the first conservative grouped-\texorpdfstring{$P_2$}{P 2} gate}
\label{sec:interface-why}

The $2.5$PN notes already showed that the surviving universal point-particle
route is the orbital/\allowbreak worldtube STF quadrupole branch.  In conservative
language, that means the next higher-order front end is not an unfocused list of
new PN coefficients, but the grouped real quadrupole bundle
\[
P_{20}\oplus P_{21}\oplus P_{22}.
\]
The exact $3$PN comparable-mass residual makes that statement quantitative.
After the carried self/static seed is subtracted, the solved COM target splits
into three pieces,
\begin{equation}
\Delta \hat L_3^{\rm GR}
=
\Delta l_1 v^8
+
(\Delta l_6,\ldots,\Delta l_{14})
+
\Delta l_{15}U^4,
\label{eq:interface-com-split-raw}
\end{equation}
so the grouped real $P_2$ lane can only hope to carry the nine-slot middle block
\((\Delta l_6,\ldots,\Delta l_{14})\).  That is exactly what makes $3$PN the
first hard conservative grouped-$P_2$ gate.  It is the earliest order at which
the grouped branch is forced to confront a nontrivial exact target, but it is
still low enough that the pass/fail test is sharp rather than diffuse.

This point should be emphasized because it changes the interpretation of the
later radiative program.  If the grouped real $P_2$ ontology had failed at
$3$PN, then the $2.5$PN narrowing would have been only suggestive packaging.
Because it succeeds, the grouped quadrupole route becomes the unique surviving
conservative carrier of the middle interaction block in the same fixed chart in
which the later passive/outgoing normalization question is posed.

\subsection{The \texorpdfstring{$a_2=b_2=0$}{a 2=b 2=0} isotropy test and grouped low-frequency data}
\label{sec:interface-isotropy}

At the first conservative grouped stage, the low-frequency data are still only
three real $O(\omega^2)$ moments,
\[
u_2^{(20)},\qquad u_2^{(21)},\qquad u_2^{(22)}.
\]
It is cleaner to rewrite them as one grouped trace plus two anisotropy defects,
\begin{equation}
\bar u_2=
\frac{u_2^{(20)}+2u_2^{(21)}+2u_2^{(22)}}{5},
\qquad
 a_2=
\frac{2u_2^{(20)}-u_2^{(21)}-u_2^{(22)}}{10},
\qquad
 b_2=
\frac{u_2^{(21)}-u_2^{(22)}}{2},
\label{eq:interface-trace-defects}
\end{equation}
with inverse map
\begin{equation}
 u_2^{(20)}=\bar u_2+4a_2,
\qquad
 u_2^{(21)}=\bar u_2-a_2+b_2,
\qquad
 u_2^{(22)}=\bar u_2-a_2-b_2.
\label{eq:interface-inverse-map}
\end{equation}
So the first grouped branch test is simply
\begin{equation}
\boxed{a_2=0,
\qquad
b_2=0.}
\label{eq:interface-first-isotropy-gate}
\end{equation}
If \eqref{eq:interface-first-isotropy-gate} fails, then the minimal isotropic
quadrupole branch is already dead before any outgoing-normalization calculation
is attempted.  If it passes, then the three grouped moments collapse to one
trace datum \(\bar u_2\), and the future conservative problem is narrowed again.

That is precisely where the $3$PN and $2.5$PN programs meet.  The present paper
fixes the first grouped data that a future moving-throat calculation must
reproduce; the $2.5$PN package then explains what comes next on the minimal
isotropic branch.  At the next even order one must determine
\(u_4^{(20)},u_4^{(21)},u_4^{(22)}\), test isotropy again, and then check the
single-pole branch identity
\begin{equation}
\boxed{
 a_4=b_4=0,
\qquad
 \bar u_4=4\bar u_2^{\,2},
\qquad
 \Omega_Q^2=\frac{1}{4\bar u_2}.
}
\label{eq:interface-second-isotropy-gate}
\end{equation}
Equivalently, once the isotropic conservative trace data are known through
$O(\omega^4)$, the minimal outgoing quadrupole branch is parametrized by a pole
scale \(\Omega_Q\) and one static normalization.  In the grouped-response
language used by the $2.5$PN package, the target normalization may be written as
\begin{equation}
\boxed{
\overline K_0^{\rm target}
=
\frac{2G}{45c^5\bar u_2^{5/2}}
=
\frac{64G}{45c^5}\,\Omega_Q^5.
}
\label{eq:interface-k0-target}
\end{equation}
So the practical theorem test has become very concrete: $3$PN supplies the first
half of the grouped conservative verdict, and the later outgoing bridge asks
whether the same branch continues through
\eqref{eq:interface-second-isotropy-gate} and
\eqref{eq:interface-k0-target}.

\subsection{What the present \texorpdfstring{$3$}{3}PN result fixes for the minimal outgoing \texorpdfstring{$\ell=2$}{ell=2} branch}
\label{sec:interface-fixes}

The main positive conclusion is that the grouped real $P_2$ ontology survives
its first exact conservative stress test.  The richer local grouped family
closes the entire middle block,
\begin{equation}
L_{P_2}^{\rm mid}
=
\sum_{A\in\{20,21,22\}}
\bigl(
\lambda_A^{\rm GR}T_A+\sigma_A^{\rm GR}S_A+\tau_A^{\rm GR}V_A
\bigr),
\label{eq:interface-grouped-module}
\end{equation}
with no middle-block residual left.  This is the crucial conservative interface
statement.  The throat side does \emph{not} need a qualitatively new non-$P_2$
carrier to reproduce the nine-slot $3$PN interaction burden.  The remaining
static scalar slot is uniquely geometry-side, and the apparent pure-kinetic COM
residual is only the ordinary-chart counterimage of the exact free relativistic
two-body Hamiltonian.  Neither of those lanes competes with the grouped
quadrupole module for the middle-block role.

This matters directly for the outgoing $\ell=2$ branch isolated at $2.5$PN.
Once the grouped real $P_2$ sector is accepted as the correct conservative front
end, the later passive/outgoing problem no longer has to search over unrelated
higher-order carriers.  It only has to determine the low-frequency constitutive
data of the same grouped branch on the true moving-throat solution.  On the
minimal isotropic branch the payoff is large: the conservative and dissipative
coefficients are no longer independent.  If one writes the isotropic conservative
module as
\begin{equation}
\overline K_Q^{\rm cons}(\omega)
=
\overline K_0
\left[
\frac34+\frac14\frac{1}{1-\omega^2/\Omega_Q^2}
\right],
\label{eq:interface-single-pole-module}
\end{equation}
then the higher even moment and the odd Burke--Thorne coefficient follow
automatically,
\begin{equation}
\boxed{
\overline K_2=\frac{\overline K_0}{4\Omega_Q^2},
\qquad
\overline K_4=\frac{\overline K_0}{4\Omega_Q^4},
\qquad
\overline\Gamma_5=\frac{9\overline K_0}{32\Omega_Q^5}.
}
\label{eq:interface-single-pole-ledger}
\end{equation}
So the present $3$PN theorem does not yet determine the outgoing coefficient
numerically, but it does determine something almost as important: where that
coefficient must come from.  It must arise from the grouped real $P_2$ lane,
with the geometry lane supplying only the separate static completion and the
free-kinetic lane supplying only the universal compiler image.

\subsection{What the full moving-throat PDE still has to supply}
\label{sec:interface-open}

The remaining theorem gap is therefore much narrower than ``derive higher PN
somehow.''  What the full moving-throat PDE still has to supply is now specific.
First, it must produce the actual low-frequency constitutive origin of the
conservative grouped data on the true branch, rather than only their reduced
bookkeeping image.  Second, it must decide whether that branch is genuinely
isotropic, first at $O(\omega^2)$ through
\eqref{eq:interface-first-isotropy-gate} and then at $O(\omega^4)$ through
\eqref{eq:interface-second-isotropy-gate}.  Third, it must realize the final
passive/outgoing normalization of the surviving quadrupole route.

The last point is the one surviving universal $2.5$PN theorem condition.  In the
canonically normalized STF language of the theorem audit, it may be written as
\begin{equation}
\boxed{\hat m_0^{\,2}\Gamma_5=\frac{2G}{5c^5}.}
\label{eq:interface-final-normalization}
\end{equation}
On the minimal isotropic outgoing branch,
\eqref{eq:interface-final-normalization} is equivalent to the static scaling law
already displayed in \eqref{eq:interface-k0-target}.  So even here the gap is
no longer diffuse: the missing PDE work must compute one isotropic normalization
product on the actual moving-throat branch, not reopen the whole coefficient
problem.

That is the correct way to read the present $3$PN theorem.  It removes the
conservative algebraic uncertainty beneath the $2.5$PN bridge, but it does not
pretend to close the remaining passive/outgoing normalization problem by itself.
The future moving-throat calculation still has to show that the true branch
realizes the same grouped-$P_2$ constitutive law, the same isotropy verdict, and
the same normalized outgoing quadrupole coefficient that the reduced $2.5$PN
package has already isolated as the final serious theorem gate.

% Section content for 4d_3pn.tex

\section{Fixed versus open quantities}
\label{sec:fixed-open}

The final theorem closes the conservative $3$PN algebraic ledger, but it does not
erase the status boundary that runs through the whole series.  As in the full
$1$PN, $2$PN, and $2.5$PN papers, one must keep separate what is exact in the
carried reduction backbone, what is frozen as lower-order input, what is fixed
by the present $3$PN audit, and what still belongs to the unsolved
moving-throat completion.  This section records that boundary in one place.
Its purpose is not to weaken the theorem, but to state its true strength: the
conservative $3$PN coefficient problem is now closed in the fixed ADM chart,
while the remaining open work has been narrowed to a much smaller dynamical and
constitutive problem.

\subsection{Frozen carry-forward inputs}
\label{sec:fixed-open-inputs}

The present paper does not begin from a blank effective theory.  It imports the
exact $4$D parent reduction backbone already frozen by the earlier papers: the
projection map, projected continuity with leakage, the exact longitudinal
identity, the quasi-static Poisson hook, and the coherent-defect / worldtube
particle reduction.  These formulas remain the structural dictionary that turns
the bulk ontology into the near-zone two-body language used throughout the
post-Newtonian series.

The lower-order conservative ledger is imported in the same way.  The carried
Newtonian and $1$PN constants are
\begin{equation}
\kappa_\rho=1,
\qquad
n=5,
\qquad
\kappa_{\rm add}=\frac12,
\qquad
\kappa_{\rm PV}=\frac32,
\qquad
\beta_{\rm 1PN}=3,
\label{eq:fixed-open-lower-order-ledger}
\end{equation}
and the full conservative Newtonian+$1$PN and $2$PN assemblies are already
frozen in the same declared closure hierarchy.  In particular, the $1$PN EIH
match and the $2$PN ADM match are inherited here as fixed lower-order output,
not reopened as free data.

The $2.5$PN narrowing is also a frozen input.  The surviving universal
higher-order route had already been reduced to the orbital/worldtube STF
quadrupole, packaged conservatively as the grouped real bundle
\begin{equation}
P_{20}\oplus P_{21}\oplus P_{22}.
\label{eq:fixed-open-grouped-bundle}
\end{equation}
That means the present $3$PN paper is not searching over arbitrary new higher-PN
carriers.  It is testing whether the same grouped real $P_2$ ontology survives
its first exact conservative stress test.

\subsection{Quantities fixed by the present \texorpdfstring{$3$}{3}PN audit}
\label{sec:fixed-open-fixed}

The first new fixed data are the strict one-body $3$PN repair parameters.  The
exact Schwarzschild gate forces
\begin{equation}
\mu_{\rho,3}=\frac14,
\qquad
 d_3=-\frac{45}{4},
\qquad
 s_{24}=-\frac1{16}.
\label{eq:fixed-open-onebody-ledger}
\end{equation}
So the carried $2$PN denominator-style self sector does \emph{not} extend to
$3$PN with one new scalar knob.  The one-body gate already fixes a cubic static
datum, a cubic denominator datum, and one genuinely new self-sector slot.

The second new fixed object is the comparable-mass middle-block compiler.  The
richer grouped real $P_2$ family closes the full nine-slot middle interaction
block exactly,
\begin{equation}
L_{P_2}^{\rm mid}
=
\sum_{A\in\{20,21,22\}}
\bigl(
\lambda_A^{\rm GR}T_A+\sigma_A^{\rm GR}S_A+\tau_A^{\rm GR}V_A
\bigr),
\label{eq:fixed-open-grouped-closure}
\end{equation}
with a nondegenerate exact compiler matrix obeying
\begin{equation}
\det M_{\rm mid}=-\frac4{27}.
\label{eq:fixed-open-det-mid}
\end{equation}
This is the precise sense in which the algebraic $3$PN bottleneck is gone: the
full middle block is no longer a symbolic comparable-mass residual waiting for a
future solve.

The third new fixed result is the chart-level uniqueness statement.  Once the
lower-order ledger is frozen and the fixed ADM Hamiltonian target is imposed,
the generic-frame cubic compiler obeys the exact sign-flip rule
\begin{equation}
\Delta H_3=-\Delta L_3.
\label{eq:fixed-open-signflip}
\end{equation}
So the earlier COM-blind null directions and the smaller ordinary contact/gauge
orbit are intermediate bookkeeping structures only.  They do not survive as
true freedom after the full generic-frame Hamiltonian target is imposed.

Finally, the scalar pieces of the exact $3$PN split are fixed.  The unique
geometry-side completion is
\begin{equation}
\Delta l_{15}^{(g)}
=
\frac{\nu\bigl(408\nu^2+1232\nu-2080+63\pi^2\bigr)}{96},
\label{eq:fixed-open-geometry-coeff}
\end{equation}
while the apparent pure-kinetic COM remainder is
\begin{equation}
\Delta l_1
=
\frac{3\nu(3\nu-1)(4\nu-1)}{16}.
\label{eq:fixed-open-kinetic-coeff}
\end{equation}
The meaning of \eqref{eq:fixed-open-kinetic-coeff} is important: it is not a new
throat-response datum.  It is only the ordinary-Lagrangian counterimage of the
universal free relativistic two-body COM Hamiltonian.

Taken together, the fixed $3$PN theorem package may be summarized as
\begin{equation}
\Delta \hat L_3^{\rm GR}
=
\underbrace{\Delta l_1 v^8}_{\text{free-Hamiltonian compiler image}}
+
\underbrace{L_{P_2}^{\rm mid}}_{\text{exact grouped real }P_2\text{ closure}}
+
\underbrace{\Delta l_{15}^{(g)}U^4}_{\text{unique geometry completion}}.
\label{eq:fixed-open-final-split}
\end{equation}
So the quantities fixed by the present audit are not just isolated coefficients;
they are the complete three-lane conservative $3$PN decomposition in the chosen
chart.

\begin{center}
\renewcommand{\arraystretch}{1.15}
\begin{tabular}{@{}>{\raggedright\arraybackslash}p{0.28\linewidth}>{\raggedright\arraybackslash}p{0.16\linewidth}>{\raggedright\arraybackslash}p{0.48\linewidth}@{}}
\textbf{Object} & \textbf{Status} & \textbf{Role at $3$PN} \\
\hline
Parent reduction backbone & frozen & Exact/reduced dictionary inherited from the $4$D parent theory \\
Lower-order Newtonian+$1$PN+$2$PN ledger & frozen & Fixed conservative input and chart conventions \\
$\mu_{\rho,3},d_3,s_{24}$ & fixed here & Exact one-body $3$PN repair ledger \\
$\lambda_A^{\rm GR},\sigma_A^{\rm GR},\tau_A^{\rm GR}$ & fixed here & Exact grouped real $P_2$ closure of the nine-slot middle block \\
$\Delta l_{15}^{(g)}$ & fixed here & Unique geometry-side static completion \\
$\Delta l_1$ & fixed here & Free-COM compiler image, not a new throat datum \\
\end{tabular}
\end{center}

\subsection{Quantities that remain symbolic or open}
\label{sec:fixed-open-open}

What remains open is no longer algebraic coefficient matching.  The open problem
is the deeper moving-throat realization of already identified lanes.  First, the
present paper does not derive the grouped real $P_2$ constitutive map from a
completed moving-throat PDE.  It proves that the grouped family is the correct
conservative carrier of the middle block, but it does not yet compute the true
branch-side low-frequency origin of the coefficients in
\eqref{eq:fixed-open-grouped-closure}.  Likewise, the paper fixes the unique
geometry-side completion coefficient \eqref{eq:fixed-open-geometry-coeff}, but
it does not yet derive that completion from a finished distributed geometry
problem on the moving throat.

Second, the dynamic throat-response observables carried forward from the earlier
papers remain genuinely open.  A representative list is
\begin{equation}
\Omega_{1\perp},\quad
\Omega_{10},\quad
\Omega_0,\quad
\Omega_{20},\quad
\Omega_{21},\quad
\Omega_{22},\quad
\Omega_g.
\label{eq:fixed-open-dynamic-poles}
\end{equation}
These are not missing algebraic residues in the conservative $3$PN theorem.
They are the dynamical pole/compliance scales of the underlying mouth / support /
geometry response problem.  The present paper is allowed to use the already
solved zero-frequency conservative ledger, but it is not allowed to pretend that
\eqref{eq:fixed-open-dynamic-poles} have already been derived from the completed
moving-throat PDE.

Third, the narrowed $2.5$PN quadrupole bridge remains open in exactly the sense
already isolated by the theorem audit.  The future moving-throat calculation
still has to determine the actual low-frequency grouped real $P_2$ data on the
true branch, decide the isotropy verdict
\begin{equation}
a_2=0,
\qquad
b_2=0,
\label{eq:fixed-open-isotropy-open}
\end{equation}
and then realize the final passive/outgoing quadrupole normalization on that
branch.  The $3$PN theorem removes the conservative algebraic uncertainty
\emph{beneath} this bridge; it does not close the outgoing normalization problem
itself.

Finally, several sectors remain outside the honest scope of the present paper:
dissipative leakage and radiation reaction beyond the already narrowed
quadrupole route, spin couplings, strong-field completion, and higher-PN work
beyond the solved conservative $3$PN scope.  So the correct negative statement
is the same one already used in the lower-order series: the present paper is a
full conservative derivation within a declared hierarchy, not an assumption-free
theorem of the fully solved moving-throat bulk dynamics.

\subsection{Interpretation}
\label{sec:fixed-open-interpretation}

The point of this section is therefore simple.  After the present audit, the
$3$PN problem no longer waits on another coefficient search, a hidden
contact-transformation freedom, or an unexplained comparable-mass residual.  All
of that is fixed.  What remains is the dynamical throat explanation of data that
are already algebraically solved: the grouped real $P_2$ constitutive law, the
geometry-side completion, the dynamic pole structure, and the passive/outgoing
quadrupole normalization on the true branch.

That is the sharpest way to state the current roadmap.  The conservative $3$PN
ledger is now part of the fixed carry-forward backbone.  The open work begins
only when one asks the completed moving-throat PDE to realize the same three
lanes microscopically and then feed them into the already narrowed $2.5$PN
quadrupole bridge.

% Section content for 4d_3pn.tex

\section{Verification and reproducibility ledger}
\label{sec:verification}

The derivation in this paper is analytic rather than computational.  The
accompanying symbolic archive is therefore not presented as a black-box
generator of the theorem.  Its role is narrower and more auditable: it serves
as a referee ledger for the exact symbolic identities, compiler maps, residual
extractions, grouped-$P_2$ closure statements, sigma-collapse formulas, and
fixed-chart assembly relations used in the main text.  This is the same stance
already adopted elsewhere in the series.  The archive verifies the paper as
written \emph{within the declared closure hierarchy}; it does not replace any
analytic step, and it does not upgrade the present paper into an
assumption-free theorem of a fully solved moving-throat bulk PDE.

At $3$PN the verification story is naturally split into two layers.  The first
layer is the inherited lower-order archive carried forward from the $4$D/$1$PN
/$2$PN papers, which freezes the reduction backbone, conventions, and already
solved coefficient ledger used unchanged here.  The second layer is the
dedicated $3$PN referee suite, whose job is to check the new theorem chain
specific to the present paper: the strict one-body gate, the cubic
ordinary/Hamiltonian compiler, the exact COM target extraction, the
Hamiltonian-first generic-frame lift, the richer grouped real $P_2$ middle-block
closure, the sigma-collapse and geometry completion, the pure-kinetic collapse,
and the final conservative $3$PN theorem package.  In the published bundle this
second layer is dual-CAS: the stagewise SymPy master runner is paired with a
bundle-local Mathematica manifest whose standalone \texttt{.wl} audits mirror
the same theorem chain and carry embedded \texttt{Output:} transcripts.

\subsection{Role of the \texorpdfstring{$3$}{3}PN referee suite}
\label{sec:verification-role}

The present referee suite should be read as a symbolic reproducibility archive,
not as an independent source of physical assumptions.  Its function is to make
each load-bearing algebraic step re-runnable in a transparent way.  In
particular, the suite is designed to verify the statements that are most likely
to matter to a referee reading the paper as a theorem chain:

\begin{itemize}
  \item whether the exact Schwarzschild one-body gate really forces the repair
  ledger \(\mu_{\rho,3}=1/4\), \(d_3=-45/4\), and \(s_{24}=-1/16\);
  \item whether the cubic Legendre-transform identity used throughout the
  comparable-mass lift is exact;
  \item whether the imported GR COM target and the reduced self/static seed
  really produce the quoted residual vector;
  \item whether the grouped real $P_2$ middle block closes exactly only in the
  richer local family, with the stated determinant and without hidden leftover
  middle-block freedom;
  \item and whether the apparent static and pure-kinetic leftovers really do
  collapse into the unique geometry-side completion and the free-Hamiltonian
  compiler image claimed in the main text.
\end{itemize}

So the archive is not being asked to prove a different theorem from the one
stated in the paper.  It is being asked to verify that the paper's own exact,
reduced, and closure-level statements close where they are claimed to close.

\subsection{Master archives}
\label{sec:verification-master}

The main SymPy verification artifact for the present paper is
\path{3pn_referee_master_sympy_audit.py}.  It is a standalone SymPy
referee script whose explicit purpose is to replay the full symbolic derivation
chain corresponding to the current $3$PN writeup.  The master runner replays a
bundle-local stage archive, executes the stages in isolated namespaces, and
stops immediately on the first failed identity.  This design keeps the
executable theorem chain visible to a referee: there is no hidden notebook
state, no dependence on local project imports, and no ambiguity about which
exact stage source is being executed.

Two implementation details are especially important for reproducibility.  First,
each extracted stage file carries a pinned SHA256 digest inside the master
runner, so the suite is not merely a loose script collection but a self-checking
archive with explicit integrity tags.  Second, the master script uses only the
Python standard library plus SymPy.  That makes the audit easy to rerun on a
clean symbolic environment and aligns with the paper's stated goal: to provide
an auditable backstop for the algebra, not a fragile local build system.

The parallel Mathematica entry point is the manifest file
\path{mathematica/wl_notes.txt}.  It lists the referee-facing
\texttt{3pn\_*\_mathematica\_audit.wl} scripts that mirror the same theorem
chain in grouped tranches: one-body and grouped kickoff, comparable-mass
compiler foundations, imported GR COM target, generic-frame projection/contact,
Hamiltonian-first lift and fixed-chart compiler, grouped-$P_2$ closure, and the
scalar/kinetic/final-theorem closeout.  Each of those scripts runs under
\texttt{math -script} and ends with an embedded \texttt{Output:} block so that
the archived transcript can be diffed directly against live stdout.

\subsection{Stage audits replayed by the master runner}
\label{sec:verification-stages}

The SymPy stages are organized in the same order as the theorem chain in the
paper, and the Mathematica manifest mirrors the same decomposition at a coarser
but still referee-readable tranche level.

\begin{enumerate}
  \item \textbf{Stages 1--4: one-body gate, grouped kickoff, and exact COM
  compiler.}  The first block consists of
  \texttt{3pn\_onebody\_audit.py},\\
  \texttt{3pn\_grouped\_p2\_audit.py},\\
  \texttt{3pn\_comparable\_mass\_audit.py}, and\\
  \texttt{3pn\_com\_linear\_map\_audit.py}.  These stages expand the exact
  Schwarzschild target, solve the minimal one-body repair ledger, freeze the
  natural self/static seed, set up the grouped real $P_2$ bookkeeping, and
  derive the exact diagonal COM ordinary/Hamiltonian compiler.

  \item \textbf{Stages 5--11: imported GR target, generic-frame projection, and
  Hamiltonian-first fixed-chart lift.}  The second block consists of
  \texttt{3pn\_com\_gr\_target\_audit.py},\\
  \texttt{3pn\_generic\_frame\_com\_projection\_audit.py},\\
  \texttt{3pn\_seed\_repair\_and\_com\_null\_ideal\_audit.py},\\
  \texttt{3pn\_contact\_generator\_and\_gauge\_orbit\_audit.py},\\
  \texttt{3pn\_generic\_frame\_target\_import\_audit.py},\\
  \texttt{3pn\_hamiltonian\_level\_lift\_audit.py}, and\\
  \texttt{3pn\_generic\_frame\_hamiltonian\_compiler\_audit.py}.  This block
  checks the exact GR COM target and residual extraction, the COM projection and
  seed-repair story, the COM-null and contact/gauge structures, the
  Hamiltonian-first lift, and the fixed-chart uniqueness theorem.

  \item \textbf{Stages 12--18: grouped real $P_2$ closure, static geometry
  completion, and final theorem ledger.}  The last block consists of
  \texttt{3pn\_grouped\_p2\_target\_pack\_audit.py},\\
  \texttt{3pn\_grouped\_p2\_middle\_block\_test\_audit.py},\\
  \texttt{3pn\_grouped\_p2\_richer\_lift\_audit.py},\\
  \texttt{3pn\_static\_p0\_geometry\_counterterm\_audit.py},\\
  \texttt{3pn\_sigma\_collapse\_and\_unique\_geometry\_completion\_audit.py},\\
  \texttt{3pn\_pure\_kinetic\_collapse\_audit.py}, and\\
  \texttt{3pn\_conservative\_theorem\_audit.py}.  These stages check the
  grouped target pack, the failure of the minimal grouped lift, the success of
  the richer grouped family, the sigma-collapse, the unique geometry-side static
  completion, the pure-kinetic collapse to the free COM Hamiltonian, and the
  final conservative $3$PN split.
\end{enumerate}

This modular staging is intentional.  It keeps the one-body gate, the
comparable-mass compiler, the grouped-middle-block closure, and the final static
and kinetic collapses separate, so that failure in one sector cannot be hidden
by compensating changes in another.

\subsection{What the suite verifies}
\label{sec:verification-verifies}

The suite verifies the load-bearing symbolic statements of the paper sector by
sector.

First, the early stages verify the exact one-body and compiler backbone.  They
check the strict Schwarzschild $3$PN target, confirm that the carried $2$PN
one-body package is insufficient by itself, solve the minimal repair ledger,
and then prove the cubic perturbative Legendre-transform identity used to move
cleanly between ordinary and Hamiltonian descriptions at $3$PN.

Second, the middle stages verify the full comparable-mass reduction and
fixed-chart uniqueness chain.  They check the exact imported COM target and the
resulting residual vector, the generic-frame projection and seed-repair story,
the COM-null ideal and the smaller contact/gauge orbit, and finally the
Hamiltonian-first lift that fixes the generic-frame representative uniquely in
the ADM chart.  In the language of the main text, this is the stage chain that
justifies the exact sign-flip rule
\begin{equation}
\Delta H_3=-\Delta L_3.
\label{eq:verification-signflip}
\end{equation}

Third, the later stages verify the grouped real $P_2$ theorem chain.  They check
that the grouped target data are packaged correctly, that the minimal grouped
local lift fails, that the richer grouped family succeeds, and that the exact
middle-block compiler is nondegenerate with
\begin{equation}
\det M_{\rm mid}=-\frac{4}{27}.
\label{eq:verification-det-mid}
\end{equation}
The suite therefore verifies the strongest positive conservative claim of the
paper: the full nine-slot $3$PN middle block is carried by the grouped real
$P_2$ module and does not require a qualitatively new non-$P_2$ conservative
carrier.

Fourth, the terminal stages verify the two apparent leftovers and the final
three-lane theorem package.  They check the sigma-collapse that removes the
apparent static \(P_0/g\) ambiguity, the unique geometry-side completion, the
pure-kinetic collapse to the universal free relativistic two-body COM
Hamiltonian, and the exact final split
\begin{equation}
\Delta \hat L_3^{\rm GR}
=
\Delta l_1 v^8
+
L_{P_2}^{\rm mid}
+
\Delta l_{15}^{(g)}U^4.
\label{eq:verification-final-split}
\end{equation}
In other words, the suite verifies that the conservative $3$PN bottleneck is no
longer algebraic once the declared hierarchy and the fixed ADM chart are
accepted.

Four reproducibility conventions are worth recording explicitly.

\paragraph{Named stage checks.}
Each stage carries a human-readable title and reduces its main claim to an
explicit symbolic residual, determinant check, solved coefficient relation, or
exact coefficient comparison.  The reader is therefore not asked to trust a
single final ``zero'' or ``PASS'' line without context.

\paragraph{Local assumptions and chart choices.}
The relevant hypotheses are attached to the stages that need them.  COM
reduction, Hamiltonian-first lifting, grouped isotropy bookkeeping, and the
fixed-chart convention are all made explicit where they enter, so the script
structure reflects the same claim taxonomy used in the paper.

\paragraph{Embedded integrity metadata.}
Because each stage source is shipped inside the master script together with its
SHA256 digest, the archive records exactly which symbolic derivation text is
being executed.  This makes accidental stage drift easier to detect than in a
loose notebook workflow.

\paragraph{Explicit bases and modular audit trail.}
Reproducibility comes from exposing the actual bases and invariant ledgers being
used: the $15$-slot COM basis, the generic-frame block decomposition, the
COM-null ideal, the grouped real $P_2$ target pack, and the final three-lane
split.  The theorem is therefore audited in the same modular language in which
it is proved.

\subsection{What the suite does not verify}
\label{sec:verification-does-not}

It is equally important to state the limits of the present archive.

First, the referee suite does not solve the fully dynamical moving-throat bulk
PDE.  The paper remains a full conservative $3$PN derivation only within the
same declared closure hierarchy already used at lower order.  The scripts verify
the algebra of that theorem attempt; they do not replace the missing
microscopic completion.

Second, the suite does not by itself derive the future microscopic origin of the
grouped real $P_2$ constitutive law or of the geometry-side completion.  Those
objects are algebraically fixed in the present paper, but their true
moving-throat realization is still deferred to the later PDE-side program.
Likewise, the suite does not close the passive/outgoing quadrupole normalization
problem on the true branch.  The $3$PN result is a conservative theorem beneath
that bridge, not a replacement for it.

Third, the suite is intentionally limited to the sectors addressed in the
present paper.  It does not claim to verify spin couplings, strong-field
completion, dissipative leakage or radiation reaction, or higher-PN sectors
beyond the present conservative $3$PN gate.

\paragraph{Bottom line.}
The purpose of the archive is not to ask the reader to trust a computation
instead of a derivation.  It is to provide an auditable backstop showing that,
under the stated hierarchy, the load-bearing symbolic statements close exactly:
\begin{equation}
\mu_{\rho,3}=\frac14,
\qquad
 d_3=-\frac{45}{4},
\qquad
 s_{24}=-\frac1{16},
\qquad
\Delta H_3=-\Delta L_3,
\qquad
\det M_{\rm mid}=-\frac{4}{27},
\label{eq:verification-bottom-line-ledger}
\end{equation}
while the final theorem package itself takes the exact form
\begin{equation}
\Delta \hat L_3^{\rm GR}
=
\Delta l_1 v^8
+
L_{P_2}^{\rm mid}
+
\Delta l_{15}^{(g)}U^4.
\label{eq:verification-bottom-line-split}
\end{equation}
What remains open is not another algebraic coefficient search, but the
microscopic moving-throat realization of the grouped real $P_2$ lane, the
geometry-side completion, and the final passive/outgoing quadrupole
normalization on the true branch.

% Section content for 4d_3pn.tex

\section{Conclusions}
\label{sec:conclusions}

This paper closes the conservative two-body $3$PN problem in the same precise
sense that the earlier full papers closed the conservative $1$PN and $2$PN
sectors: not as an assumption-free theorem of a fully solved moving-throat PDE,
but as an exact algebraic assembly within a declared closure hierarchy.  The
gain is larger than one more coefficient list.  By $3$PN the defect is forced to
carry a genuinely nontrivial comparable-mass middle interaction block, and the
present theorem shows that this burden is borne exactly by the grouped real
$P_2$ lane, with the remaining static scalar burden collapsing to a unique
geometry-side completion and the apparent pure-kinetic remainder collapsing to
the ordinary-chart image of the universal free relativistic two-body
Hamiltonian.  In that sharpened sense, the conservative $3$PN bottleneck is no
longer an unresolved coefficient search.

\subsection{What has been established}
\label{sec:conclusions-established}

The proof line established in the paper may be stated compactly.

\begin{enumerate}
  \item The strict one-body Schwarzschild gate through $c^{-6}$ is fixed
  exactly, and its minimal repair ledger is not one-dimensional but
  \begin{equation}
  \mu_{\rho,3}=\frac14,
  \qquad
   d_3=-\frac{45}{4},
  \qquad
   s_{24}=-\frac1{16}.
  \label{eq:conclusions-onebody-ledger}
  \end{equation}
  So the carried $2$PN self-sector package does not extend to $3$PN by a single
  denominator coefficient alone.

  \item The exact cubic ordinary/Hamiltonian compiler is rigid, and once the
  lower-order ledger is frozen the correct comparable-mass lift is
  Hamiltonian-first.  In the fixed ADM chart this becomes the exact sign-flip
  identity
  \begin{equation}
  \Delta H_3=-\Delta L_3,
  \label{eq:conclusions-signflip}
  \end{equation}
  so the generic-frame ordinary representative is unique after the Hamiltonian
  target is imposed.

  \item The richer local grouped real $P_2$ family closes the full nine-slot
  middle interaction block exactly.  Its compiler is nondegenerate,
  \begin{equation}
  \det M_{\rm mid}=-\frac4{27},
  \label{eq:conclusions-det-mid}
  \end{equation}
  and the exact GR coefficient triples
  \((\lambda_A^{\rm GR},\sigma_A^{\rm GR},\tau_A^{\rm GR})\) reproduce the
  entire solved middle block with no residual left behind.

  \item The apparent static $P_0/g$ ambiguity is not a surviving freedom.
  Sigma-collapse removes the one-parameter family and leaves one unique
  geometry-side completion,
  \begin{equation}
  \Delta l_{15}^{(g)}
  =
  \frac{\nu\bigl(408\nu^2+1232\nu-2080+63\pi^2\bigr)}{96}.
  \label{eq:conclusions-geometry-coeff}
  \end{equation}

  \item The apparent center-of-mass pure-kinetic remainder is not a new
  throat-response coefficient.  It is fixed automatically by the exact free
  two-body relativistic Hamiltonian and therefore appears only as the
  ordinary-Lagrangian compiler image
  \begin{equation}
  \Delta l_1
  =
  \frac{3\nu(3\nu-1)(4\nu-1)}{16}.
  \label{eq:conclusions-kinetic-coeff}
  \end{equation}
\end{enumerate}

When these facts are assembled, the full comparable-mass $c^{-6}$ coefficient is
assigned by three and only three lanes,
\begin{equation}
\boxed{
\Delta \hat L_3^{\rm GR}
=
\underbrace{\Delta l_1 v^8}_{\text{free-Hamiltonian compiler image}}
+
\underbrace{L_{P_2}^{\rm mid}}_{\text{exact grouped real }P_2\text{ closure}}
+
\underbrace{\Delta l_{15}^{(g)}U^4}_{\text{unique geometry completion}}.
}
\label{eq:conclusions-final-split}
\end{equation}
There is no fourth independent conservative carrier hiding behind the COM
bookkeeping.  Within the declared hierarchy, the final Hamiltonian statement is
therefore exact:
\begin{equation}
H_3^{\rm derived}=H_3^{\rm ADM},
\qquad
\Delta H_3\equiv H_3^{\rm derived}-H_3^{\rm ADM}=0.
\label{eq:conclusions-adm-equality}
\end{equation}

Conceptually, the paper does three things at once.  First, it upgrades the
series from the full conservative $2$PN ADM match to the first genuinely
nontrivial higher conservative grouped-$P_2$ gate.  Second, it proves that the
middle interaction burden is already carried by the grouped real quadrupole
lane, rather than by some qualitatively different hidden conservative carrier.
Third, it sharpens the bookkeeping of what is fixed, what is closure-level, and
what must still remain dynamical.  The accompanying referee archive supports
this reading in an auditable way: the $18$-stage SymPy master suite checks the
one-body gate, the cubic compiler, the imported target extraction, the
Hamiltonian-first lift, the grouped-$P_2$ closure, the sigma-collapse, the
pure-kinetic collapse, and the final theorem ledger inside the same declared
hierarchy.

\subsection{What remains open}
\label{sec:conclusions-open}

The negative statement is equally important.

First, the paper does not solve the fully dynamical moving-throat PDE and then
recover the $3$PN theorem as an automatic bulk consequence.  The present result
remains a full conservative derivation only within the declared reduction and
closure hierarchy already used by the lower-order papers.  That status is not a
weakness of the theorem; it is the correct statement of its scope.

Second, the paper does not yet derive the grouped real $P_2$ constitutive law or
the geometry-side completion microscopically from the completed moving-throat
branch.  It proves that the grouped real $P_2$ module is the correct
conservative carrier of the nine-slot middle block and that the residual static
slot is uniquely geometry-side, but it does not yet compute those lanes from the
underlying distributed throat dynamics.

Third, the dynamic pole/compliance observables carried forward from the earlier
papers remain genuinely open objects of the unsolved moving-throat response
problem.  Representative scales are
\begin{equation}
\Omega_{1\perp},\quad
\Omega_{10},\quad
\Omega_0,\quad
\Omega_{20},\quad
\Omega_{21},\quad
\Omega_{22},\quad
\Omega_g.
\label{eq:conclusions-open-poles}
\end{equation}
These are not leftover algebraic residues in the present $3$PN theorem.  They
are dynamical response data of the microscopic mouth/support/geometry system.

Fourth, the present paper does not close the already narrowed $2.5$PN outgoing
quadrupole normalization problem.  The remaining universal target may be stated
in the invariant form
\begin{equation}
\boxed{\hat m_0^{\,2}\Gamma_5=\frac{2G}{5c^5}.}
\label{eq:conclusions-normalization-target}
\end{equation}
The $3$PN theorem removes the conservative algebraic uncertainty \emph{beneath}
this bridge, but it does not by itself determine the final passive/outgoing
normalization on the physical moving-throat branch.

Finally, several larger sectors remain outside the honest scope of the present
paper: dissipative leakage and radiation reaction beyond the already narrowed
quadrupole route, spin couplings, strong-field completion, and higher-PN work
beyond the solved conservative $3$PN scope.  What remains open is therefore no
longer coefficient matching.  It is the deeper dynamical throat explanation of
data that are already algebraically assigned.

\subsection{Interface to the next calculation}
\label{sec:conclusions-next}

This gives the project a much cleaner handoff than before.  The next job is no
longer to reopen the $3$PN comparable-mass algebra, to search for another hidden
contact freedom, or to rerun a vague whole-theory higher-PN scan.  That work is
substantially finished in the present chart.  The next job is the
moving-throat grouped-quadrupole calculation itself.

On the conservative side, the physical branch must first determine the actual
low-frequency grouped real $P_2$ data and decide whether the minimal isotropic
verdict holds,
\begin{equation}
\boxed{a_2=0,
\qquad
b_2=0.}
\label{eq:conclusions-first-isotropy-gate}
\end{equation}
If that first gate fails, then the minimal isotropic quadrupole branch is
already dead before outgoing normalization is attempted.  If it passes, then the
next conservative task is the $O(\omega^4)$ grouped moment test and the
associated single-pole bookkeeping,
\begin{equation}
\boxed{
\bar u_4=4\bar u_2^{\,2},
\qquad
\Omega_Q^2=\frac{1}{4\bar u_2},
\qquad
\overline K_0^{\rm target}=\frac{2G}{45c^5\bar u_2^{5/2}}.
}
\label{eq:conclusions-second-isotropy-gate}
\end{equation}
Equivalently, once the physical grouped branch is known through its first even
moments, the later outgoing Burke--Thorne coefficient is no longer an
independent datum.

So the correct next calculation is sharply defined.  One must derive the grouped
real $P_2$ constitutive data and the geometry-side completion from the physical
moving-throat branch, determine whether the isotropic passive/outgoing branch is
actually realized there, and then evaluate the single invariant normalization
product in \eqref{eq:conclusions-normalization-target}.  That is the only place
where the present theorem still needs to be sharpened from a closure-level
conservative result into a fully microscopic one.

The strongest honest endpoint of the paper is therefore also the strongest honest
starting point for the next one: within the declared hierarchy, the conservative
$3$PN bottleneck is gone, the grouped real $P_2$ and geometry lanes are the
correct carriers, and the only serious universal bridge question left open is
the passive/outgoing normalization of that already identified quadrupole branch.

\section*{Acknowledgments}
This work was carried out by an independent researcher. Generative-AI tools were used
during the development of the manuscript to help translate conceptual ideas into
mathematical form, assist with symbolic manipulations, prototype and revise companion
Mathematica and SymPy verification scripts, and improve the exposition. The conceptual
direction, hypotheses, and overall structure of the work were directed by the author. The
resulting arguments and calculations were checked for internal consistency using the
accompanying symbolic scripts, but the work should be regarded as exploratory and would
benefit from review by subject-matter experts. Any errors, omissions, or misunderstandings
remain the author's.

\appendix

% The appendix map below follows 3pn_notes.md essentially one-to-one.

% Appendix content for 4d_3pn.tex

\section{Carry-forward formulas from the earlier papers}
\label{app:carry-forward}

This appendix collects the specific formulas and status statements imported from
 the parent $4$D paper and the full conservative $1$PN, $2$PN, and $2.5$PN
 packages and used without rederivation in the present $3$PN work.  Its role is
 technical rather than historical.  We do not repeat the full parent action, the
 full lower-order proofs, or the longer appendix-side derivations that fixed the
 wake, breathing, support, and quadrupole-response data.  Instead, we record only
 the load-bearing identities, reduced equations, and frozen coefficients that the
 present $3$PN theorem chain actually uses.

The claim-status boundary is the same one used throughout the earlier papers.
Exact projection identities from the parent $4$D theory are kept visibly exact.
The near-zone Poisson hook, the zero-mode Maxwell law, and the worldtube
 point-particle reduction are carried only as controlled reductions.  The
 conservative Newtonian+$1$PN and $2$PN ledgers are treated as frozen outputs of
 the earlier program within its declared closure hierarchy, not as live
 parameters to be retuned in the present paper.  The $2.5$PN package is used in
 the same way: not as a finished moving-throat PDE theorem, but as the theorem
 audit that narrows the surviving higher-order universal point-particle route to
 the orbital/worldtube STF quadrupole branch and repackages its conservative data
 as the grouped real $P_{20}\oplus P_{21}\oplus P_{22}$ bundle.

\subsection{Newtonian and \texorpdfstring{$1$}{1}PN carried data}
\label{app:carry-forward-1pn}

Before recording the Newtonian and $1$PN coefficient ledger itself, it is useful
 to state the minimal parent $4$D reduction dictionary through which that ledger
 is carried into the post-Newtonian sector.  The parent theory lives on a
 $(4+1)$-dimensional spacetime with coordinates
\[
X^M=(t,x,y,z,w),
\qquad
\mathbf{X}=(x,y,z,w)\in\mathbb R^4,
\qquad
\mathbf{x}=(x,y,z)\in\mathbb R^3.
\]
The bulk variables carried into the present paper are the matter field
 $\psi(\mathbf X,t)$, its density
\[
\rho(\mathbf X,t)=|\psi(\mathbf X,t)|^2,
\]
the optional localized gauge field $A_M(\mathbf X,t)$ with localization profile
 $Z(w)$, and the effective throat geometry variables
\[
a(t),\qquad L(t).
\]
When electric charge is named explicitly in this inherited downstream notation,
 it is always the corrected branch quantity
\[
q_*=\eta_Q e_*,
\qquad
q_{\rm eff}=\frac{q_*}{\sqrt{Z_{\rm int}}},
\qquad
Z_{\rm int}=\int_{-\infty}^{\infty} Z(w)\,dw,
\]
not the gravity-side scalar coefficient $\kappa_\rho$ used in the post-Newtonian
 ledger below.

Brane observables are defined by a fixed normalized projection kernel $W(w)$,
\[
\int_{-\infty}^{\infty} W(w)\,dw = 1,
\qquad
\mathcal P_W[Q](t,\mathbf x)
\equiv
\int_{-\infty}^{\infty} W(w)\,Q(t,\mathbf x,w)\,dw.
\]
In particular,
\[
\rho_{\rm br}=\mathcal P_W[\rho],
\qquad
\mathbf J_{\rm br}=\mathcal P_W[(j^x,j^y,j^z)],
\qquad
\mathbf v_{\rm br}=\frac{\mathbf J_{\rm br}}{\rho_{\rm br}}
\quad (\rho_{\rm br}>0).
\]
Projection is not the same thing as reduction.  Projection defines what a brane
 observer measures; reduction refers to the further elimination of the extra
 coordinate under an explicit ansatz or scaling regime.

From exact bulk continuity,
\[
\partial_t\rho+\partial_i j^i=0,
\qquad i\in\{x,y,z,w\},
\]
one obtains the exact open-system brane continuity law
\[
\partial_t\rho_{\rm br}+\nabla_3\!\cdot\mathbf J_{\rm br}=S_{\rm leak},
\]
with leakage source
\[
S_{\rm leak}
=
-\bigl[W(w)j^w\bigr]_{-\infty}^{+\infty}
+
\int_{-\infty}^{\infty} W'(w)j^w\,dw.
\]
After Helmholtz decomposition of the brane velocity,
\[
\mathbf v_{\rm br}=\nabla_3\phi_{\rm br}+\mathbf v_T,
\qquad
\nabla_3\!\cdot\mathbf v_T=0,
\]
projection continuity becomes the exact longitudinal identity
\[
\rho_{\rm br}\,\nabla_3^2\phi_{\rm br}
=
S_{\rm leak}
-\partial_t\rho_{\rm br}
-(\nabla_3\rho_{\rm br})\cdot(\nabla_3\phi_{\rm br}+\mathbf v_T).
\]
Under the usual quasi-static near-zone simplifications this becomes the carried
 Poisson hook
\[
\nabla_3^2\phi_{\rm br}\simeq \frac{S_{\rm eff}}{\rho_0},
\qquad
\Phi_N=\lambda_\Phi\,\phi_{\rm br},
\]
which is the origin of the Newtonian scalar sector used in the lower-order
 particle reduction.  The corresponding coherent-defect / small-body worldtube
 limit gives the point-particle force law
\[
M_A\,\ddot{\mathbf X}_A
=
-M_A\nabla\Phi_{\rm ext}(\mathbf X_A)
+O\!\left(\frac{a^2}{r^2}\,\frac{\Phi_{\rm ext}}{r}\right),
\]
with the first finite-size correction quadrupolar and therefore suppressed in the
 controlled small-worldtube regime.  For completeness, the same parent paper also
 carries a controlled zero-mode Maxwell reduction.  Under the far-field brane
 assumptions
\[
A_w\approx 0,
\qquad
\partial_w A_\mu\approx 0,
\qquad
J^w\approx 0,
\qquad
F_{\mu w}\approx 0,
\]
integration over $w$ gives the effective brane Maxwell law
\[
\partial_\mu F^{\mu\nu}=\mu_0^{\rm eff}J_{\rm eff}^\nu,
\qquad
\mu_0^{\rm eff}=\frac{\mu_0}{Z_{\rm int}}.
\]
This reduction suppresses the mixed channels only in the controlled far-field
 limit.  It does not erase $(A_w,J^w,F_{\mu w})$ from the microscopic ontology.

The Newtonian point-particle block carried into every later paper is therefore
\[
L_0
=
\frac12 m_A v_A^2
+\frac12 m_B v_B^2
+\frac{Gm_A m_B}{r}.
\]
At $1$PN the scalar/optical one-body reduction fixes the lower-order viability
 conditions
\[
\kappa_\rho=1,
\qquad
n=5,
\qquad
\kappa_{\rm add}=\frac12,
\qquad
\kappa_{\rm PV}=\frac32,
\qquad
\beta_{\rm 1PN}=3.
\]
The same response-side closure also freezes the breathing ledger
\[
E_w:E_f:E_{\rm PV}=11:2:5,
\qquad
\frac{d\ln a_*}{d\ln\rho}=-\frac{57}{64}.
\]
In the parity-even wake sector, the constructive isotropic Fourier-space closure
 fixes
\[
\alpha^2=\frac34,
\qquad
a_H=0,
\qquad
K_{\rm vec}=\frac{2}{\pi^2},
\]
which in turn give the EIH cross coefficients
\[
C_\parallel=-\frac72,
\qquad
C_L=-\frac12.
\]
With these values frozen, the full carried Newtonian+$1$PN two-body block is
\[
\begin{aligned}
L_1=
&\frac18(m_Av_A^4+m_Bv_B^4)
+\frac{Gm_A m_B}{r}
\left[
\frac32(v_A^2+v_B^2)
-\frac72 v_{AB}
-\frac12 v_{An}v_{Bn}
\right]
\\
&-\frac{G^2m_A m_B(m_A+m_B)}{2r^2}.
\end{aligned}
\]
Nothing in the present $3$PN paper is allowed to retune these lower-order values.
 They are carry-forward anchors rather than open coefficients.

\subsection{\texorpdfstring{$2$}{2}PN carried data and the \texorpdfstring{$P_0\oplus P_2$}{P 0 + P 2} support hierarchy}
\label{app:carry-forward-2pn}

The full conservative $2$PN paper closes the comparable-mass $c^{-4}$ ledger
 against the standard generic-frame ADM target within the same declared closure
 hierarchy, but the present $3$PN paper imports from it only the load-bearing
 formulas and the response-side structure that sharpened the ontology of the
 defect.  At the strict one-body level, the exact isotropic Schwarzschild gate
 through $2$PN is packaged by
\[
D_{\rm target}(u)=1-4u+8u^2+O(u^3),
\qquad
u\equiv\frac{U}{c^2},
\]
and the pure-static $2$PN gate fixes
\[
\lambda_\rho=\frac12,
\qquad
L_{2,\rm static}=\frac{G^3m_A m_B(m_A^2+m_B^2)}{4r^3}.
\]
The natural $2$PN self/static seed promoted from the exact one-body target is
\[
\begin{aligned}
L_{2,\mathrm{self+static}}
&= \frac{m_Av_A^6+m_Bv_B^6}{16}
 + \frac{7Gm_A m_B}{8r}(v_A^4+v_B^4)
\\
&\quad + \frac{2G^2m_A m_B}{r^2}(m_Bv_A^2+m_Av_B^2)
\\
&\quad + \frac{G^3m_A m_B(m_A^2+m_B^2)}{4r^3}.
\end{aligned}
\]
The full carried $2$PN coefficient inserted into the Newtonian+$1$PN ledger is
\[
\begin{aligned}
L_{2,\rm full}
={}&\frac{m_Av_A^6+m_Bv_B^6}{16}
\\
&+\frac{Gm_A m_B}{r}
\Big[
\frac78(v_A^4+v_B^4)
-\frac74 v_{AB}(v_A^2+v_B^2)
\\
&\hspace{4em}
\hspace{0pt}-\frac14 v_{An}v_{Bn}(v_A^2+v_B^2)
+\frac{11}{8}v_A^2v_B^2
\\
&\hspace{4em}
\frac14 v_{AB}^2
-\frac58(v_A^2v_{Bn}^2+v_B^2v_{An}^2)
+\frac32 v_{AB}v_{An}v_{Bn}
+\frac38 v_{An}^2v_{Bn}^2
\Big]
\\
&+\frac{G^2m_A m_B}{r^2}
\Big[
\Big(2m_B+\frac{11}{8}m_A\Big)v_A^2
+\Big(2m_A+\frac{11}{8}m_B\Big)v_B^2
\\
&\hspace{4em}
-\frac{15}{4}(m_A+m_B)v_{AB}
+\frac{15}{8}(m_Av_{An}^2+m_Bv_{Bn}^2)
\Big]
\\
&+\frac{G^3m_A m_B}{4r^3}(m_A^2+5m_A m_B+m_B^2).
\end{aligned}
\]
So by the start of the present paper, the full conservative Newtonian+$1$PN+$2$PN
 ledger is already frozen in the declared hierarchy.

Conceptually, the decisive ontology change at $2$PN is that the defect is no
 longer described adequately as a pure scalar monopole.  The conservative cross
 sector already splits into three response lanes:
\begin{enumerate}
  \item a carried odd dipole wake sector,
  \item a genuine even $P_0\oplus P_2$ mouth/support layer,
  \item and a separate geometry-closure channel.
\end{enumerate}
The appendix-side zero-frequency throat data make this explicit.  The odd dipole
 residues are frozen to
\[
R_{1\perp}=\frac72,
\qquad
R_{10}=4,
\]
while the canonical even support residues are
\[
R_0=R_{20}=R_{21}=R_{22}=1.
\]
The same solved static even operator fixes the source-weight vector
\[
J=\left(\frac{4}{\sqrt5},\frac54,0,0,0,0\right),
\qquad
J\cdot J=\frac{381}{80},
\]
and the pure-potential geometry-closure deficit
\[
\Delta_{\rm geom}=\frac{281}{80}.
\]
The monopole wall add-on
\[
\Delta K_{00}=\frac{109}{280}
\]
is kept distinct from $\Delta_{\rm geom}$; they are not the same object and
 should not be merged in later bookkeeping.

What remains open after the conservative $2$PN theorem is not these
 zero-frequency residues but the fully dynamical completion.  In particular, the
 pole or compliance scales
\[
\Omega_{1\perp},\ \Omega_{10},\ \Omega_0,
\ \Omega_{20},\ \Omega_{21},\ \Omega_{22},\ \Omega_g
\]
are not claimed as already derived from the completed moving-throat PDE.  This
 status boundary is central for the present paper.  The $3$PN calculation is
 allowed to use the fixed $2$PN residues and source data, but it is not allowed
 to treat the fully dynamical throat completion as already solved.

\subsection{\texorpdfstring{$2.5$}{2.5}PN narrowing to the grouped real \texorpdfstring{$P_2$}{P 2} target}
\label{app:carry-forward-25pn}

The $2.5$PN theorem audit is what makes the present $3$PN problem sharply
 focused rather than diffuse.  The decisive benchmark of that audit shows that
 the frozen conservative hierarchy by itself cannot generate a nonzero local
 $2.5$PN reaction theorem, while a minimal retarded quadrupole odd term lands on
 the standard Burke--Thorne / Iyer--Will family member
\[
\alpha=4,
\qquad
\beta=5.
\]
For the present conservative paper, however, the important carry-forward message
 is narrower: after the scalar, dipole, and quadrupole channel audits, the live
 higher-order universal point-particle route collapses to the orbital/worldtube
 STF quadrupole branch.  In grouped real language, that means the right
 conservative front end is the grouped real
\[
P_{20}\oplus P_{21}\oplus P_{22}
\]
bundle, not an unfocused collection of generic higher-order coefficients.

Accordingly, the grouped conservative response is organized as
\[
Y_{20}(\omega)=1+u_2^{(20)}\omega^2+u_4^{(20)}\omega^4+O(\omega^6),
\]
\[
Y_{21}(\omega)=1+u_2^{(21)}\omega^2+u_4^{(21)}\omega^4+O(\omega^6),
\]
\[
Y_{22}(\omega)=1+u_2^{(22)}\omega^2+u_4^{(22)}\omega^4+O(\omega^6).
\]
It is convenient to repackage these grouped triples into trace/anomaly variables.
 At $O(\omega^2)$ one defines
\[
\bar u_2=\frac{u_2^{(20)}+2u_2^{(21)}+2u_2^{(22)}}{5},
\qquad
a_2=\frac{2u_2^{(20)}-u_2^{(21)}-u_2^{(22)}}{10},
\qquad
b_2=\frac{u_2^{(21)}-u_2^{(22)}}{2},
\]
and at $O(\omega^4)$
\[
\bar u_4=\frac{u_4^{(20)}+2u_4^{(21)}+2u_4^{(22)}}{5},
\qquad
a_4=\frac{2u_4^{(20)}-u_4^{(21)}-u_4^{(22)}}{10},
\qquad
b_4=\frac{u_4^{(21)}-u_4^{(22)}}{2}.
\]
The exact inverse map is
\[
u_2^{(20)}=\bar u_2+4a_2,
\qquad
u_2^{(21)}=\bar u_2-a_2+b_2,
\qquad
u_2^{(22)}=\bar u_2-a_2-b_2,
\]
\[
u_4^{(20)}=\bar u_4+4a_4,
\qquad
u_4^{(21)}=\bar u_4-a_4+b_4,
\qquad
u_4^{(22)}=\bar u_4-a_4-b_4.
\]
So the first isotropy gate for the minimal branch is simply
\[
a_2=0,
\qquad
b_2=0,
\]
and the later $O(\omega^4)$ branch test is
\[
a_4=0,
\qquad
b_4=0.
\]
At the level relevant for the present $3$PN paper, the important point is that
 the grouped real conservative front end is already completely specified by the
 three $O(\omega^2)$ coefficients
\[
u_2^{(20)},\qquad u_2^{(21)},\qquad u_2^{(22)},
\]
or, equivalently, by $(\bar u_2,a_2,b_2)$.  That is why the $3$PN paper can ask
 whether the conservative middle block closes in grouped-$P_2$ language without
 pretending to have finished the full outgoing normalization problem.

For later reference, the minimal isotropic single-pole branch of the same $2.5$PN
 package implies the exact formulas
\[
\Omega_Q^2=\frac{1}{4u_2},
\qquad
\Gamma_5^{\rm norm}=9u_2^{5/2},
\qquad
\overline K_0^{\rm target}=\frac{2G}{45c^5u_2^{5/2}},
\]
together with the $O(\omega^4)$ branch identity
\[
u_4=4u_2^2.
\]
These are not part of the present conservative $3$PN theorem, but they explain
 why the $2.5$PN narrowing makes the grouped real $P_2$ front end the natural
 higher-order conservative payload.

% Appendix content for 4d_3pn.tex

\section{Exact \texorpdfstring{$3$}{3}PN one-body Schwarzschild expansion}
\label{app:onebody-expansion}

The strict one-body gate for the present $3$PN program is the isotropic
Schwarzschild worldline action. Its role is narrow but load-bearing. Before any
one-body repair ledger or comparable-mass lift is attempted, the reduced
hierarchy must know exactly what the test-mass sector is allowed to reduce to.
The point of the present appendix is therefore simply to record the exact
Schwarzschild expansion through total order $c^{-6}$ and isolate the new $3$PN
block.

Use the standard one-body potential variable
\begin{equation}
 u\equiv \frac{U}{c^2}=\frac{GM}{rc^2}.
\label{eq:app-onebody-u-def}
\end{equation}
The exact test-mass worldline Lagrangian in isotropic Schwarzschild coordinates
is
\begin{equation}
\frac{L_{\rm Schw}}{m}
=-c^2\sqrt{\left(\frac{1-u/2}{1+u/2}\right)^2-(1+u/2)^4\frac{v^2}{c^2}}.
\label{eq:app-onebody-schwarzschild-exact}
\end{equation}
Expanding \eqref{eq:app-onebody-schwarzschild-exact} through total order
$c^{-6}$ gives the full one-body target carried into the present paper,
\begin{align}
\frac{L_{\rm Schw}}{m}
={}&-c^2+\frac12v^2+U
+\frac{1}{c^2}\left(\frac18v^4+\frac32Uv^2-\frac12U^2\right)
\nonumber\\
&
+\frac{1}{c^4}\left(\frac1{16}v^6+\frac78Uv^4+2U^2v^2+\frac14U^3\right)
\nonumber\\
&
+\frac{1}{c^6}\left(
\frac5{128}v^8
+\frac{11}{16}Uv^6
+\frac{47}{16}U^2v^4
+\frac{13}{8}U^3v^2
-\frac18U^4
\right)
+O(c^{-8}).
\label{eq:app-onebody-schwarzschild-3pn}
\end{align}
The coefficient of $c^{-6}$ is therefore the exact $3$PN one-body destination,
\begin{equation}
\boxed{
\frac{L_{3\mathrm{PN},1\mathrm{body}}}{m}
=
\frac5{128}v^8
+\frac{11}{16}Uv^6
+\frac{47}{16}U^2v^4
+\frac{13}{8}U^3v^2
-\frac18U^4.
}
\label{eq:app-onebody-schwarzschild-3pn-target}
\end{equation}
This is the strict test-mass gate that every later comparable-mass construction
must reduce to. It is the exact analogue at $3$PN of the Schwarzschild one-body
gate already used successfully at $2$PN. The logic of the present paper is
therefore deliberately sequential: do not begin with the generic two-body lift;
first freeze the exact one-body target in
\eqref{eq:app-onebody-schwarzschild-3pn-target}, then ask what additional data
the carried denominator/self/static hierarchy needs in order to hit it exactly.
That repair algebra is the subject of the next appendix.

% Appendix content for 4d_3pn.tex

\section{One-body repair algebra}
\label{app:onebody-repair}

Appendix~\ref{app:onebody-expansion} freezes the exact isotropic Schwarzschild
one-body target through total order $c^{-6}$.  The natural next question is
whether the carried $2$PN denominator/self/static packaging can simply be pushed
one order higher, or whether the strict test-mass gate already forces genuinely
new $3$PN one-body data.  The point of the present appendix is to answer that
question algebraically and coefficient by coefficient.

The result is rigid.  One cubic denominator datum is needed to repair the
$U^3v^2$ slot, one cubic static datum is needed to repair the pure $U^4$ slot,
and one additional self-sector invariant is needed because the $U^2v^4$ slot
remains wrong even after the first two repairs are imposed.  So the full
one-body $3$PN repair ledger is already three-parameter at the strict
Schwarzschild gate.

\subsection{Cubic denominator datum}
\label{app:onebody-repair-denominator}

The first test is the minimal cubic continuation of the carried denominator-side
self sector,
\begin{equation}
L_{\rm red}
=-mc^2(1-u)\sqrt{1-\frac{v^2/c^2}{D(u)}},
\qquad
u\equiv\frac{U}{c^2},
\qquad
D(u)=1-4u+8u^2+d_3u^3+O(u^4).
\label{eq:app-onebody-repair-Lred}
\end{equation}
Expanding \eqref{eq:app-onebody-repair-Lred} through total order $c^{-6}$
gives
\begin{align}
\frac{L_{\rm red}}{m}
={}&-c^2+\frac12v^2+U
+\frac{1}{c^2}\left(\frac18v^4+\frac32Uv^2\right)
\nonumber\\
&
+\frac{1}{c^4}\left(\frac1{16}v^6+\frac78Uv^4+2U^2v^2\right)
\nonumber\\
&
+\frac{1}{c^6}\left(
\frac5{128}v^8
+\frac{11}{16}Uv^6
+3U^2v^4
+\left(-4-\frac{d_3}{2}\right)U^3v^2
\right)
+O(c^{-8}).
\label{eq:app-onebody-repair-Lred-series}
\end{align}
Two features of \eqref{eq:app-onebody-repair-Lred-series} are immediate.
First, the carried denominator logic already reproduces the universal $3$PN
coefficients of $v^8$ and $Uv^6$ automatically.  Second, it leaves the
$U^3v^2$ slot open through the new cubic coefficient $d_3$.

Matching the exact Schwarzschild target
\eqref{eq:app-onebody-schwarzschild-3pn-target} in the $U^3v^2$ channel gives
\begin{equation}
-4-\frac{d_3}{2}=\frac{13}{8}
\qquad\Longrightarrow\qquad
\boxed{d_3=-\frac{45}{4}.}
\label{eq:app-onebody-repair-d3}
\end{equation}
So the cubic denominator datum is fixed exactly by the $U^3v^2$ slot.

It is worth recording that the unmodified $2$PN invariant-denominator
continuation would have predicted a completely different cubic coefficient.  If
one simply extends the carried $2$PN denominator-side packaging without opening
any genuinely new cubic invariant, the audit gives
\begin{equation}
d_3^{\rm carry}=\frac{21783}{2432},
\label{eq:app-onebody-repair-d3-carry}
\end{equation}
which is nowhere near the exact value in
\eqref{eq:app-onebody-repair-d3}.  So already at the $U^3v^2$ gate, the strict
one-body problem forces a genuine new cubic invariant correction rather than a
trivial continuation of the old $2$PN denominator game.

\subsection{Static gate}
\label{app:onebody-repair-static}

The pure static slot gives an independent gate.  Extend the carried local
mass-scaling sector by one cubic static coefficient $\mu_{\rho,3}$,
\begin{equation}
\frac{L_{\rm static,1body}}{m}
=
U-\frac12\frac{U^2}{c^2}+\frac14\frac{U^3}{c^4}
-\frac{\mu_{\rho,3}}{2}\frac{U^4}{c^6}+O(c^{-8}).
\label{eq:app-onebody-repair-static-sector}
\end{equation}
The exact Schwarzschild target requires the static $3$PN term
\begin{equation}
-\frac18\frac{U^4}{c^6}.
\label{eq:app-onebody-repair-static-target}
\end{equation}
So the cubic static gate is fixed immediately by
\begin{equation}
-\frac{\mu_{\rho,3}}{2}=-\frac18
\qquad\Longrightarrow\qquad
\boxed{\mu_{\rho,3}=\frac14.}
\label{eq:app-onebody-repair-murho3}
\end{equation}
This is the direct $3$PN analogue of the static gate that already appeared at
$2$PN.  The important point is that the static channel remains a separate exact
gate of its own; it is not a residual bookkeeping convenience.

\subsection{The extra self datum and why one cubic denominator correction is not enough}
\label{app:onebody-repair-self}

Once the cubic denominator datum and the cubic static gate are admitted, there
is still one leftover mismatch.  The minimal repaired one-body candidate is
\begin{equation}
\frac{L_{\rm cand}}{m}
=
\frac{L_{\rm red}}{m}
-\frac12\frac{U^2}{c^2}
+\frac14\frac{U^3}{c^4}
-\frac{\mu_{\rho,3}}{2}\frac{U^4}{c^6}
+s_{24}\frac{U^2v^4}{c^6},
\label{eq:app-onebody-repair-candidate}
\end{equation}
where the first three displayed corrections simply restore the carried lower
static ledger and the final term is the minimal new explicit self-sector slot.
Subtracting \eqref{eq:app-onebody-repair-candidate} from the exact target gives
\begin{equation}
\left(\frac{L_{3\mathrm{PN},1\mathrm{body}}}{m}-\frac{L_{\rm cand}}{m}\right)_{c^{-6}}
=
\left(\frac{45}{8}+\frac{d_3}{2}\right)U^3v^2
+\left(\frac{47}{16}-3-s_{24}\right)U^2v^4
+\left(-\frac18+\frac{\mu_{\rho,3}}{2}\right)U^4.
\label{eq:app-onebody-repair-residual}
\end{equation}
So the one-body repair algebra separates cleanly by slot.

The first and third coefficients in
\eqref{eq:app-onebody-repair-residual} vanish once
\eqref{eq:app-onebody-repair-d3} and
\eqref{eq:app-onebody-repair-murho3} are imposed.  But even after those two
exact repairs, the denominator-style self sector still predicts
\begin{equation}
3U^2v^4=\frac{48}{16}U^2v^4,
\label{eq:app-onebody-repair-u2v4-prediction}
\end{equation}
whereas the exact Schwarzschild target demands
\begin{equation}
\frac{47}{16}U^2v^4.
\label{eq:app-onebody-repair-u2v4-target}
\end{equation}
So there is an uncancelled one-body discrepancy of
\begin{equation}
-\frac1{16}U^2v^4.
\label{eq:app-onebody-repair-u2v4-mismatch}
\end{equation}
This cannot be absorbed by $d_3$, because $d_3$ controls only the $U^3v^2$
slot, and it cannot be absorbed by $\mu_{\rho,3}$, because
$\mu_{\rho,3}$ controls only the pure static $U^4$ slot.  The minimal repair
is therefore to admit the explicit self invariant already written in
\eqref{eq:app-onebody-repair-candidate}, with coefficient fixed by
\begin{equation}
\frac{47}{16}-3-s_{24}=0
\qquad\Longrightarrow\qquad
\boxed{s_{24}=-\frac1{16}.}
\label{eq:app-onebody-repair-s24}
\end{equation}
Collecting everything, the exact one-body repair ledger is
\begin{equation}
\boxed{
\mu_{\rho,3}=\frac14,
\qquad
 d_3=-\frac{45}{4},
\qquad
 s_{24}=-\frac1{16}.
}
\label{eq:app-onebody-repair-ledger}
\end{equation}
This is the cleanest strict one-body lesson of the whole $3$PN analysis.  The
$3$PN gate does \emph{not} say ``continue the $2$PN denominator idea by one
more coefficient.''  It says that three distinct data are already needed at the
strict test-mass level:
\begin{enumerate}
    \item a cubic static gate,
    \item a cubic denominator datum,
    \item and one genuinely new self-sector slot.
\end{enumerate}
That is exactly why the full $3$PN program had to begin with a one-body audit
before the comparable-mass lift was even well posed.  Once the exact target and
its minimal repair ledger are frozen, the next step is to promote the strict
one-body structure into the natural two-body self/static seed, which is the
subject of the next appendix.

% Appendix content for 4d_3pn.tex

\section{Natural \texorpdfstring{$3$}{3}PN self/static seed derivation}
\label{app:self-static-seed}

Appendices~\ref{app:onebody-expansion} and \ref{app:onebody-repair} freeze
both ingredients that have to exist before any honest comparable-mass $3$PN lift
can begin: the exact isotropic Schwarzschild one-body target, and the minimal
one-body repair ledger needed to realize it inside the carried self/static
packaging.  Once those are fixed, the next object is not a generic two-body
ansatz.  It is the unique symmetric \emph{bodywise} promotion of the exact
one-body data into the two-body sector.

The purpose of the present appendix is to write that promotion explicitly and to
explain why it is the correct starting seed.  The logic is the same one that
already worked at $2$PN: keep the seed strictly bodywise and symmetric, so that
(i) the strict test-mass reduction is automatic, and (ii) no genuinely
comparable-mass interaction data are smuggled into the seed before the residual
solve even begins.

\subsection{Bodywise promotion of the exact one-body target}
\label{app:self-static-seed-bodywise}

Start from the exact one-body $3$PN target from
Appendix~\ref{app:onebody-expansion},
\begin{equation}
\frac{L_{3\mathrm{PN},1\mathrm{body}}}{m}
=
\frac5{128}v^8
+\frac{11}{16}Uv^6
+\frac{47}{16}U^2v^4
+\frac{13}{8}U^3v^2
-\frac18U^4.
\label{eq:app-self-static-seed-onebody-target}
\end{equation}
For body $A$ moving in the field of body $B$, the strict bodywise substitutions
are
\begin{equation}
U_A=\frac{Gm_B}{r},
\qquad
m\to m_A,
\qquad
v\to v_A.
\label{eq:app-self-static-seed-bodyA}
\end{equation}
For body $B$ moving in the field of body $A$, they are
\begin{equation}
U_B=\frac{Gm_A}{r},
\qquad
m\to m_B,
\qquad
v\to v_B.
\label{eq:app-self-static-seed-bodyB}
\end{equation}
So the admissible self/static seed is fixed by the rule
\begin{equation}
L_{3,\mathrm{seed}}
=
m_A\left[\frac{L_{3\mathrm{PN},1\mathrm{body}}}{m}\right]_{U\to U_A,\,v\to v_A}
+
m_B\left[\frac{L_{3\mathrm{PN},1\mathrm{body}}}{m}\right]_{U\to U_B,\,v\to v_B}.
\label{eq:app-self-static-seed-rule}
\end{equation}
Evaluating \eqref{eq:app-self-static-seed-rule} gives the natural symmetric
self/static $3$PN seed,
\begin{equation}
\boxed{
L_{3,\mathrm{seed}}
=
\frac{5}{128}(m_Av_A^8+m_Bv_B^8)
+\frac{11Gm_A m_B}{16r}(v_A^6+v_B^6)
}
\label{eq:app-self-static-seed-first-line}
\end{equation}
\begin{equation}
\boxed{
\qquad
+\frac{47G^2m_A m_B}{16r^2}(m_Bv_A^4+m_Av_B^4)
+\frac{13G^3m_A m_B}{8r^3}(m_B^2v_A^2+m_A^2v_B^2)
-\frac{G^4m_A m_B}{8r^4}(m_B^3+m_A^3),
}
\label{eq:app-self-static-seed-second-line}
\end{equation}
with the whole block understood as the coefficient of $c^{-6}$ in the full
conservative expansion.

Written term by term, the map from the one-body monomials in
\eqref{eq:app-self-static-seed-onebody-target} to their two-body bodywise images
is completely rigid:
\begin{align}
mv^8
&\longmapsto m_Av_A^8+m_Bv_B^8,
\label{eq:app-self-static-seed-map-v8}
\\
mUv^6
&\longmapsto \frac{Gm_A m_B}{r}(v_A^6+v_B^6),
\label{eq:app-self-static-seed-map-uv6}
\\
mU^2v^4
&\longmapsto \frac{G^2m_A m_B}{r^2}(m_Bv_A^4+m_Av_B^4),
\label{eq:app-self-static-seed-map-u2v4}
\\
mU^3v^2
&\longmapsto \frac{G^3m_A m_B}{r^3}(m_B^2v_A^2+m_A^2v_B^2),
\label{eq:app-self-static-seed-map-u3v2}
\\
mU^4
&\longmapsto \frac{G^4m_A m_B}{r^4}(m_B^3+m_A^3).
\label{eq:app-self-static-seed-map-u4}
\end{align}
No extra freedom appears in this dictionary once the seed is required to remain
bodywise and symmetric.

\subsection{Exact test-mass reduction and uniqueness}
\label{app:self-static-seed-testmass}

The reason the bodywise rule is the right one is that it preserves the strict
one-body gate automatically.  Take the test-mass limit with body $A$ light,
body $B$ heavy, and the heavy body unperturbed at leading order:
\begin{equation}
m_A\to \mu,
\qquad
m_B\to M,
\qquad
v_B\to 0,
\qquad
U_A\to \frac{GM}{r}.
\label{eq:app-self-static-seed-testmass-limit}
\end{equation}
Then \eqref{eq:app-self-static-seed-first-line}--\eqref{eq:app-self-static-seed-second-line}
reduce immediately to
\begin{equation}
L_{3,\mathrm{seed}}
\longrightarrow
\mu\left(
\frac5{128}v_A^8
+\frac{11}{16}U_Av_A^6
+\frac{47}{16}U_A^2v_A^4
+\frac{13}{8}U_A^3v_A^2
-\frac18U_A^4
\right),
\label{eq:app-self-static-seed-testmass-reduction}
\end{equation}
which is exactly the strict Schwarzschild target
\eqref{eq:app-self-static-seed-onebody-target} for the light body.
The heavy-body copy collapses because $v_B\to0$, so it does not contaminate the
light-body dynamics.

This already gives the uniqueness statement needed later in the paper.  Suppose
one asks for a $3$PN seed with the following three properties:
\begin{enumerate}
    \item it is symmetric under $A\leftrightarrow B$,
    \item it reproduces the exact one-body target in the strict test-mass limit,
    \item it contains only bodywise/self-static images of the one-body monomials,
          and no genuinely comparable-mass cross-velocity structures.
\end{enumerate}
Then each monomial in
\eqref{eq:app-self-static-seed-onebody-target} has exactly one admissible
symmetric image, namely the one listed in
\eqref{eq:app-self-static-seed-map-v8}--\eqref{eq:app-self-static-seed-map-u4}.
So the seed
\eqref{eq:app-self-static-seed-first-line}--\eqref{eq:app-self-static-seed-second-line}
is not just natural; within that bodywise/self-static class it is unique.

\subsection{Why this is the correct starting seed for the \texorpdfstring{$3$}{3}PN lift}
\label{app:self-static-seed-properties}

The seed has three structural properties that make it the correct starting point
for the full $3$PN lift.

First, it isolates the universal bodywise content and nothing else.  The seed is
built only from the exact one-body data already forced by the strict
Schwarzschild gate.  It therefore contains the universal free $v^8$ term, the
bodywise $Uv^6$, $U^2v^4$, and $U^3v^2$ self sectors, and the static $U^4$
sector, but it contains no genuinely comparable-mass middle-block structures of
the form
\[
v_A^2v_B^6,
\qquad
v_A^4v_B^4,
\qquad
\text{or analogous mixed invariant combinations.}
\]
Those belong to the comparable-mass residual, not to the seed.

Second, it makes the later residual split conceptually clean.  The full $3$PN
problem can therefore be organized as
\begin{equation}
L_3=L_{3,\mathrm{seed}}+\Delta L_3,
\label{eq:app-self-static-seed-split}
\end{equation}
where $\Delta L_3$ is the pure comparable-mass interaction residual to be solved
on the later $24/17/8/1$ generic-frame scaffold.  In other words, once the seed
is frozen, every remaining $3$PN interaction coefficient can be interpreted as
``beyond the exact one-body/self-static core'' rather than as a blur of bodywise
and comparable-mass information.

Third, the seed is sparse when carried into the later COM chart.  Before the
lower-order compiler feedback is applied, only the self/static lanes are
switched on directly.  That sparseness is precisely what makes the imported COM
residual finite and interpretable: the later exact $15$-slot residual vector is
not ``the whole $3$PN answer,'' but the part of the $3$PN answer that remains
once the unique bodywise seed has already been removed.

That is why the object derived here is the right seed and not merely a plausible
starting guess.  It is the unique symmetric promotion of the exact one-body
$3$PN data into the same self/static lane that already organized the successful
$2$PN construction, and it is the last universal input that should be frozen
before the genuine comparable-mass $3$PN solve begins.

% Appendix content for 4d_3pn.tex

\section{Exact cubic Legendre-transform identity at \texorpdfstring{$3$}{3}PN}
\label{app:cubic-legendre}

The $2$PN assembly already relied on the perturbative Legendre transform through
quadratic order.  The first genuinely new compiler ingredient at $3$PN is the
\emph{cubic-order} extension.  The purpose of the present appendix is to freeze
that identity in the compact matrix form actually used throughout the $3$PN
construction.

The role of this appendix is therefore very specific.  It does not solve the
full comparable-mass problem by itself.  What it does is fix the exact ordinary
Lagrangian $\to$ Hamiltonian compiler through order $\varepsilon^3$ with
\(\varepsilon=c^{-2}\).  Once that identity is available, the later reduced COM
and generic-frame lifts are controlled algebraically rather than guessed.

\subsection{Perturbative setup}
\label{app:cubic-legendre-setup}

Write the conservative ordinary Lagrangian as
\begin{equation}
L(v)=L_0(v)+\varepsilon L_1(v)+\varepsilon^2L_2(v)+\varepsilon^3L_3(v),
\qquad
\varepsilon=c^{-2},
\label{eq:app-cubic-legendre-L-expansion}
\end{equation}
with quadratic zeroth-order block
\begin{equation}
L_0(v)=\frac12 v^T M v,
\label{eq:app-cubic-legendre-L0}
\end{equation}
where the Newtonian mass matrix \(M\) is constant at the order relevant for the
perturbative inversion.  Throughout this appendix, gradients with respect to
\(v\) are taken as column vectors, and Hessians are taken as matrices.

The canonical momentum is
\begin{equation}
p\equiv \frac{\partial L}{\partial v}
= Mv
+\varepsilon\frac{\partial L_1}{\partial v}
+\varepsilon^2\frac{\partial L_2}{\partial v}
+\varepsilon^3\frac{\partial L_3}{\partial v}.
\label{eq:app-cubic-legendre-momentum}
\end{equation}
Define the Newtonian momentum-space velocity
\begin{equation}
v_0=M^{-1}p,
\label{eq:app-cubic-legendre-v0}
\end{equation}
and the lower-order derivative data evaluated at \(v_0\),
\begin{equation}
A_0\equiv \left.\frac{\partial L_1}{\partial v}\right|_{v_0},
\qquad
B_0\equiv \left.\frac{\partial L_2}{\partial v}\right|_{v_0},
\qquad
C_0\equiv \left.\frac{\partial^2L_1}{\partial v^2}\right|_{v_0}.
\label{eq:app-cubic-legendre-ABC}
\end{equation}

The perturbative momentum inversion is written as
\begin{equation}
v=v_0+\varepsilon v_1+\varepsilon^2v_2+O(\varepsilon^3).
\label{eq:app-cubic-legendre-v-series}
\end{equation}
For the cubic Hamiltonian coefficient \(H_3\), only \(v_1\) and \(v_2\) are
needed explicitly.

\subsection{Order-by-order momentum inversion}
\label{app:cubic-legendre-inversion}

Insert \eqref{eq:app-cubic-legendre-v-series} into
\eqref{eq:app-cubic-legendre-momentum} and expand around \(v_0\).  To the order
needed for \(H_3\), the Taylor expansions are
\begin{equation}
\frac{\partial L_1}{\partial v}(v)
=
A_0 + \varepsilon C_0 v_1 + O(\varepsilon^2),
\qquad
\frac{\partial L_2}{\partial v}(v)
=
B_0 + O(\varepsilon).
\label{eq:app-cubic-legendre-taylor}
\end{equation}
Using \(Mv_0=p\), the momentum equation becomes
\begin{equation}
0
=\varepsilon\bigl(Mv_1+A_0\bigr)
+\varepsilon^2\bigl(Mv_2+B_0+C_0v_1\bigr)
+O(\varepsilon^3).
\label{eq:app-cubic-legendre-momentum-orders}
\end{equation}
So the coefficients are fixed order by order:
\begin{equation}
Mv_1+A_0=0,
\qquad
Mv_2+B_0+C_0v_1=0.
\label{eq:app-cubic-legendre-order-eqs}
\end{equation}
Therefore
\begin{equation}
\boxed{
v_1=-M^{-1}A_0,
\qquad
v_2=M^{-1}C_0M^{-1}A_0-M^{-1}B_0.
}
\label{eq:app-cubic-legendre-v1v2}
\end{equation}
This is the exact perturbative inversion needed through cubic order.

\subsection{Exact cubic Hamiltonian identity}
\label{app:cubic-legendre-identity}

Now substitute \eqref{eq:app-cubic-legendre-v-series} and
\eqref{eq:app-cubic-legendre-v1v2} into
\begin{equation}
H(p)=p\cdot v-L(v).
\label{eq:app-cubic-legendre-H-def}
\end{equation}
The Hamiltonian expansion takes the form
\begin{equation}
H=H_0+\varepsilon H_1+\varepsilon^2H_2+\varepsilon^3H_3+O(\varepsilon^4),
\label{eq:app-cubic-legendre-H-series}
\end{equation}
with zeroth-order piece
\begin{equation}
H_0=\frac12 p^TM^{-1}p.
\label{eq:app-cubic-legendre-H0}
\end{equation}
Carrying the cubic expansion through exactly gives
\begin{equation}
H_1=-L_1(v_0),
\label{eq:app-cubic-legendre-H1}
\end{equation}
\begin{equation}
H_2=-L_2(v_0)+\frac12A_0^TM^{-1}A_0,
\label{eq:app-cubic-legendre-H2}
\end{equation}
and the key $3$PN compiler identity
\begin{equation}
\boxed{
H_3=-L_3(v_0)+A_0^TM^{-1}B_0-
\frac12A_0^TM^{-1}C_0M^{-1}A_0.
}
\label{eq:app-cubic-legendre-H3}
\end{equation}
Equation~\eqref{eq:app-cubic-legendre-H3} is the exact cubic analogue of the
quadratic compiler identity already used at $2$PN.  It is the formula actually
used to carry the ordinary $3$PN data into the Hamiltonian chart.

It is also useful to record the reduced COM specialization.  In the normalized
COM variables used later, the Newtonian quadratic block is the identity, so
\(M=I\) and therefore
\begin{equation}
v_0=p,
\qquad
H_3=-L_3(p)+A_0\cdot B_0-\frac12A_0^TC_0A_0.
\label{eq:app-cubic-legendre-H3-com}
\end{equation}
This is the compact form that the referee audit checks explicitly before the
later slot-by-slot COM compiler is built.

\subsection{Practical meaning for the \texorpdfstring{$3$}{3}PN lift}
\label{app:cubic-legendre-meaning}

The practical meaning of \eqref{eq:app-cubic-legendre-H3} is the same one that
made the $2$PN paper manageable, but now one order higher.  Once the lower-order
$1$PN/$2$PN ledger is frozen, the feedback tensors \(A_0\), \(B_0\), and
\(C_0\) are already known.  So the genuinely new $3$PN Hamiltonian content
enters only through the explicit \(-L_3(v_0)\) term, while the remaining cubic
feedback terms are completely fixed by carried lower-order data.

That observation is what turns the comparable-mass $3$PN problem into a finite
algebraic residual rather than a blind full derivation campaign.  In the later
reduced COM basis the map generated by
\eqref{eq:app-cubic-legendre-H3} diagonalizes, and in the fixed generic-frame
Hamiltonian chart the residual compiler becomes rigid enough that the remaining
$3$PN bookkeeping can be solved slot by slot.

% Appendix content for 4d_3pn.tex

\section{COM basis and diagonal compiler derivation}
\label{app:com-compiler}

Appendix~\ref{app:cubic-legendre} fixes the exact cubic ordinary-Lagrangian
$\to$ Hamiltonian compiler through order $c^{-6}$.  The purpose of the present
appendix is to specialize that compiler to the isotropic reduced
center-of-mass chart actually used in the $3$PN construction and to show that,
in the natural $15$-slot basis, the map reduces \emph{diagonally}.

That diagonalization is one of the main algebraic simplifications of the whole
$3$PN program.  Once the lower-order reduced COM blocks are frozen, the target
Hamiltonian coefficients $h_i$ determine the ordinary-Lagrangian coefficients
$l_i$ slot by slot, with no further nontrivial inversion.  This is the precise
sense in which the COM stage of the $3$PN lift is finite rather than open-ended.

\subsection{Reduced COM ordinary-Lagrangian basis}
\label{app:com-compiler-ordinary}

Write the reduced conservative COM ordinary Lagrangian in the standard isotropic
chart as
\begin{equation}
\hat L=\hat L_0+c^{-2}\hat L_1+c^{-4}\hat L_2+c^{-6}\hat L_3+O(c^{-8}),
\label{eq:app-com-compiler-L-expansion}
\end{equation}
with lower-order blocks already frozen by the carried Newtonian, $1$PN, and
$2$PN papers.  In the reduced COM variables used throughout the $3$PN session,
those blocks are
\begin{equation}
\hat L_0=\frac12 v^2+\frac1r,
\label{eq:app-com-compiler-L0}
\end{equation}
\begin{equation}
\hat L_1
=
\frac{1-3\nu}{8}v^4
+\frac1r\left[\frac{3+\nu}{2}v^2+\frac{\nu}{2}\dot r^2\right]
-\frac{1}{2r^2},
\label{eq:app-com-compiler-L1}
\end{equation}
\begin{align}
\hat L_2
={}&
\frac{1-5\nu+5\nu^2}{16}v^6
+\frac1r\Bigg[
\left(\frac78-\frac74\nu-\frac18\nu^2\right)v^4
+\left(\frac14\nu-\frac14\nu^2\right)v^2\dot r^2
+\frac38\nu^2\dot r^4
\Bigg]
\nonumber\\
&
+\frac1{r^2}\Bigg[
\left(2-\frac78\nu\right)v^2
+\frac{15}{8}\nu\dot r^2
\Bigg]
+\frac1{r^3}\left(\frac14+\frac34\nu\right).
\label{eq:app-com-compiler-L2}
\end{align}
Here
\begin{equation}
v^2=\dot r^2+v_\perp^2,
\qquad
\nu\equiv\frac{m_A m_B}{(m_A+m_B)^2},
\label{eq:app-com-compiler-v2nu}
\end{equation}
and the reduced COM normalization is the same one used throughout the main text.

The new $3$PN coefficient is then expanded in the natural $15$-slot ordinary
basis
\begin{align}
\hat L_3
={}&
 l_1 v^8 + l_2 v^6\dot r^2 + l_3 v^4\dot r^4 + l_4 v^2\dot r^6 + l_5\dot r^8
\nonumber\\
&
+\frac1r\Bigl(l_6 v^6 + l_7 v^4\dot r^2 + l_8 v^2\dot r^4 + l_9\dot r^6\Bigr)
\nonumber\\
&
+\frac1{r^2}\Bigl(l_{10} v^4 + l_{11} v^2\dot r^2 + l_{12}\dot r^4\Bigr)
+\frac1{r^3}\Bigl(l_{13} v^2 + l_{14}\dot r^2\Bigr)
+\frac{l_{15}}{r^4}.
\label{eq:app-com-compiler-L3-basis}
\end{align}
This is the reduced COM ordinary-Lagrangian ledger used everywhere in the $3$PN
solve.  Later exact comparable-mass results are always represented as
corrections to these $15$ slots beyond the frozen one-body/self-static seed.

\subsection{Reduced COM Hamiltonian basis}
\label{app:com-compiler-hamiltonian}

The matching isotropic reduced COM Hamiltonian expansion is
\begin{equation}
\hat H=\hat H_0+c^{-2}\hat H_1+c^{-4}\hat H_2+c^{-6}\hat H_3+O(c^{-8}),
\label{eq:app-com-compiler-H-expansion}
\end{equation}
and the $3$PN coefficient is written in the parallel $15$-slot Hamiltonian
basis
\begin{align}
\hat H_3
={}&
 h_1 p^8 + h_2 p^6p_r^2 + h_3 p^4p_r^4 + h_4 p^2p_r^6 + h_5 p_r^8
\nonumber\\
&
+\frac1r\Bigl(h_6 p^6 + h_7 p^4p_r^2 + h_8 p^2p_r^4 + h_9 p_r^6\Bigr)
\nonumber\\
&
+\frac1{r^2}\Bigl(h_{10} p^4 + h_{11} p^2p_r^2 + h_{12} p_r^4\Bigr)
+\frac1{r^3}\Bigl(h_{13} p^2 + h_{14} p_r^2\Bigr)
+\frac{h_{15}}{r^4}.
\label{eq:app-com-compiler-H3-basis}
\end{align}
Here
\begin{equation}
p^2=p_r^2+p_\perp^2
\label{eq:app-com-compiler-p2}
\end{equation}
is the reduced isotropic COM momentum invariant.  The choice of the monomials in
\eqref{eq:app-com-compiler-H3-basis} is not arbitrary.  It is precisely the
choice for which the cubic compiler from Appendix~\ref{app:cubic-legendre}
closes on the same $15$-slot monomial ladder and becomes diagonal.

The reduced COM specialization of the exact cubic compiler is
\begin{equation}
\hat H_3
=
-\hat L_3(v_0)
+A_0\cdot B_0
-\frac12 A_0^T C_0 A_0,
\qquad
v_0=p,
\label{eq:app-com-compiler-cubic}
\end{equation}
with carried lower-order derivative data
\begin{equation}
A_0\equiv\left.\frac{\partial \hat L_1}{\partial v}\right|_{v=p},
\qquad
B_0\equiv\left.\frac{\partial \hat L_2}{\partial v}\right|_{v=p},
\qquad
C_0\equiv\left.\frac{\partial^2 \hat L_1}{\partial v^2}\right|_{v=p}.
\label{eq:app-com-compiler-ABC}
\end{equation}
Because the carried lower-order blocks are isotropic polynomials in
\((v^2,\dot r^2,1/r)\), the feedback term
\(A_0\cdot B_0-\tfrac12A_0^TC_0A_0\) is an even polynomial in
\((p_r,p_\perp)\).  Re-expanding it in the basis
\((p^2,p_r^2,1/r)\) lands exactly on the same $15$ monomials as
\eqref{eq:app-com-compiler-H3-basis}.  That is the structural reason the final
map diagonalizes.

\subsection{Exact diagonal map and its inverse}
\label{app:com-compiler-map}

Applying \eqref{eq:app-com-compiler-cubic} to the basis
\eqref{eq:app-com-compiler-L3-basis}, and then re-expanding the result in the
Hamiltonian basis \eqref{eq:app-com-compiler-H3-basis}, gives the exact slotwise
map
\begin{equation}
\begin{aligned}
 h_1 &= -l_1+\frac{3}{16}\nu-\frac{21}{16}\nu^2+\frac94\nu^3,
 & h_2 &= -l_2,
\\
 h_3 &= -l_3,
 & h_4 &= -l_4,
\\
 h_5 &= -l_5.
\end{aligned}
\label{eq:app-com-compiler-forward-1}
\end{equation}
\begin{equation}
\begin{aligned}
 h_6 &= -l_6+\frac14+\frac78\nu-\frac{35}{8}\nu^2-\frac{21}{4}\nu^3,
 & h_7 &= -l_7+\frac{11}{8}\nu^2-\frac92\nu^3,
\\
 h_8 &= -l_8+\frac34\nu^2-\frac94\nu^3,
 & h_9 &= -l_9.
\end{aligned}
\label{eq:app-com-compiler-forward-2}
\end{equation}
\begin{equation}
\begin{aligned}
 h_{10} &= -l_{10}+\frac54+\frac{15}{8}\nu+\frac{123}{8}\nu^2+\frac{13}{4}\nu^3,
\\
 h_{11} &= -l_{11}+\frac78\nu+\frac{41}{8}\nu^2+\frac{31}{4}\nu^3,
\\
 h_{12} &= -l_{12}+\frac92\nu^2+4\nu^3.
\end{aligned}
\label{eq:app-com-compiler-forward-3}
\end{equation}
\begin{equation}
\begin{aligned}
 h_{13} &= -l_{13}-\frac32-\frac{59}{4}\nu-\frac{25}{4}\nu^2-\frac12\nu^3,
\\
 h_{14} &= -l_{14}+\frac74\nu-\frac{31}{4}\nu^2-\frac72\nu^3,
\\
 h_{15} &= -l_{15}.
\end{aligned}
\label{eq:app-com-compiler-forward-4}
\end{equation}
So the chosen ordinary basis diagonalizes the reduced COM cubic compiler exactly.
Equivalently,
\begin{equation}
\frac{\partial h_i}{\partial l_j}=-\delta_{ij}
\label{eq:app-com-compiler-jacobian}
\end{equation}
throughout the full $15$-slot basis.

The inverse map is therefore immediate:
\begin{equation}
\begin{aligned}
 l_1 &= \frac{3}{16}\nu-\frac{21}{16}\nu^2+\frac94\nu^3-h_1,
 & l_2 &= -h_2,
\\
 l_3 &= -h_3,
 & l_4 &= -h_4,
\\
 l_5 &= -h_5.
\end{aligned}
\label{eq:app-com-compiler-inverse-1}
\end{equation}
\begin{equation}
\begin{aligned}
 l_6 &= \frac14+\frac78\nu-\frac{35}{8}\nu^2-\frac{21}{4}\nu^3-h_6,
 & l_7 &= \frac{11}{8}\nu^2-\frac92\nu^3-h_7,
\\
 l_8 &= \frac34\nu^2-\frac94\nu^3-h_8,
 & l_9 &= -h_9.
\end{aligned}
\label{eq:app-com-compiler-inverse-2}
\end{equation}
\begin{equation}
\begin{aligned}
 l_{10} &= \frac54+\frac{15}{8}\nu+\frac{123}{8}\nu^2+\frac{13}{4}\nu^3-h_{10},
\\
 l_{11} &= \frac78\nu+\frac{41}{8}\nu^2+\frac{31}{4}\nu^3-h_{11},
\\
 l_{12} &= \frac92\nu^2+4\nu^3-h_{12}.
\end{aligned}
\label{eq:app-com-compiler-inverse-3}
\end{equation}
\begin{equation}
\begin{aligned}
 l_{13} &= -\frac32-\frac{59}{4}\nu-\frac{25}{4}\nu^2-\frac12\nu^3-h_{13},
\\
 l_{14} &= \frac74\nu-\frac{31}{4}\nu^2-\frac72\nu^3-h_{14},
\\
 l_{15} &= -h_{15}.
\end{aligned}
\label{eq:app-com-compiler-inverse-4}
\end{equation}
There is no hidden COM-only kernel here.  Once the target Hamiltonian
coefficients are supplied, the ordinary coefficients follow immediately and
uniquely.

That is the exact algebraic reason the later $3$PN reduced COM target import is
so efficient.  Appendix~\ref{app:gr-com-target} only has to provide the target
list \(h_i\); after that, the corresponding ordinary coefficients \(l_i\) are
already fixed by \eqref{eq:app-com-compiler-inverse-1}--\eqref{eq:app-com-compiler-inverse-4}.
In this sense the diagonal COM compiler is not just a calculational convenience.
It is the device that turns the reduced $3$PN lift from a nonlinear inversion
problem into a slotwise algebraic translation.

% Appendix content for 4d_3pn.tex

\section{Imported GR COM target coefficients}
\label{app:gr-com-target}

This appendix records the exact reduced center-of-mass $3$PN target in the
same isotropic $15$-slot Hamiltonian basis used in Appendix~F, translates that
Hamiltonian target back into the ordinary-Lagrangian chart by the exact inverse
compiler, and then isolates the residual beyond the frozen one-body/self-static
seed. This is the point at which the $3$PN comparable-mass problem becomes a
finite algebraic target rather than an open-ended generic two-body search.

\subsection{\texorpdfstring{$15$}{15}-slot target list}
\label{app:gr-com-target-list}

Using the standard reduced COM ADM Hamiltonian through $3$PN in the same
isotropic basis as Appendix~F, the exact GR target coefficients are
\begin{equation}
\begin{aligned}
 h_1&=\frac{-5+35\nu-70\nu^2+35\nu^3}{128},
 & h_2&=h_3=h_4=h_5=0.
\end{aligned}
\label{eq:app-gr-com-target-h-1-5}
\end{equation}
\begin{equation}
\begin{aligned}
 h_6&=\frac{-7+42\nu-53\nu^2-5\nu^3}{16},
 & h_7&=\frac{(2-3\nu)\nu^2}{16},
\\[0.5ex]
 h_8&=\frac{3(1-\nu)\nu^2}{16},
 & h_9&=-\frac{5\nu^3}{16}.
\end{aligned}
\label{eq:app-gr-com-target-h-6-9}
\end{equation}
\begin{equation}
\begin{aligned}
 h_{10}&=\frac{-27+136\nu+109\nu^2}{16},
 & h_{11}&=\frac{(17+30\nu)\nu}{16},
\\[0.5ex]
 h_{12}&=\frac{(5+43\nu)\nu}{12}.
\end{aligned}
\label{eq:app-gr-com-target-h-10-12}
\end{equation}
\begin{equation}
\begin{aligned}
 h_{13}&=\frac{-600+(3\pi^2-1340)\nu-552\nu^2}{192},
 & h_{14}&=-\frac{(340+3\pi^2+112\nu)\nu}{64}.
\end{aligned}
\label{eq:app-gr-com-target-h-13-14}
\end{equation}
\begin{equation}
 h_{15}=\frac{12+(872-63\pi^2)\nu}{96}.
\label{eq:app-gr-com-target-h-15}
\end{equation}
The strict test-mass limit reproduces the exact one-body Schwarzschild slots,
\begin{equation}
 h_1\to -\frac{5}{128},
\qquad
 h_6\to -\frac{7}{16},
\qquad
 h_{10}\to -\frac{27}{16},
\qquad
 h_{13}\to -\frac{25}{8},
\qquad
 h_{15}\to \frac18,
\label{eq:app-gr-com-target-test-mass-h}
\end{equation}
so the imported COM Hamiltonian target is compatible with the exact one-body
$3$PN gate from the start.

\subsection{Exact ordinary-Lagrangian solution}
\label{app:gr-com-target-solution}

Applying the inverse diagonal map from Appendix~F slot by slot gives the exact
COM ordinary-Lagrangian coefficients
\begin{equation}
\begin{aligned}
 l_1&=\frac{5-11\nu-98\nu^2+253\nu^3}{128},
 & l_2&=l_3=l_4=l_5=0.
\end{aligned}
\label{eq:app-gr-com-target-l-1-5}
\end{equation}
\begin{equation}
\begin{aligned}
 l_6&=\frac{11-28\nu-17\nu^2-79\nu^3}{16},
 & l_7&=\frac{20\nu^2-69\nu^3}{16},
\\[0.5ex]
 l_8&=\frac{9\nu^2-33\nu^3}{16},
 & l_9&=\frac{5\nu^3}{16}.
\end{aligned}
\label{eq:app-gr-com-target-l-6-9}
\end{equation}
\begin{equation}
\begin{aligned}
 l_{10}&=\frac{47-106\nu+137\nu^2+52\nu^3}{16},
 & l_{11}&=\frac{-3\nu+52\nu^2+124\nu^3}{16},
\\[0.5ex]
 l_{12}&=\frac{-5\nu+11\nu^2+48\nu^3}{12}.
\end{aligned}
\label{eq:app-gr-com-target-l-10-12}
\end{equation}
\begin{equation}
 l_{13}=\frac{13}{8}-\left(\frac{373}{48}+\frac{\pi^2}{64}\right)\nu
 -\frac{27}{8}\nu^2-\frac12\nu^3.
\label{eq:app-gr-com-target-l-13}
\end{equation}
\begin{equation}
\begin{aligned}
 l_{14}&=\frac{452+3\pi^2}{64}\nu-6\nu^2-\frac72\nu^3,
 & l_{15}&=-\frac18-\frac{872-63\pi^2}{96}\nu.
\end{aligned}
\label{eq:app-gr-com-target-l-14-15}
\end{equation}
Three structural facts are worth keeping explicit.

First, the pure quartic-radial slots remain absent,
\begin{equation}
 l_2=l_3=l_4=l_5=0.
\label{eq:app-gr-com-target-no-quartic-radial}
\end{equation}
Second, the strict one-body limit reproduces the carried self/static seed,
\begin{equation}
 l_1\to \frac{5}{128},
\qquad
 l_6\to \frac{11}{16},
\qquad
 l_{10}\to \frac{47}{16},
\qquad
 l_{13}\to \frac{13}{8},
\qquad
 l_{15}\to -\frac18.
\label{eq:app-gr-com-target-test-mass-l}
\end{equation}
Third, the $\pi^2$ dependence first enters only in the deep-potential
$r^{-3}$ and $r^{-4}$ slots, exactly where the genuinely new static $3$PN
structure sits.

\subsection{Residual beyond the self/static seed}
\label{app:gr-com-target-residual}

Let the frozen self/static seed occupy the nonzero COM ordinary slots
\begin{align}
 l_1^{\rm seed}&=\frac{5}{128}-\frac{35}{128}\nu+\frac{35}{64}\nu^2-\frac{35}{128}\nu^3,
\label{eq:app-gr-com-target-seed-l1}
\\[0.5ex]
 l_6^{\rm seed}&=\frac{11}{16}-\frac{33}{8}\nu+\frac{99}{16}\nu^2-\frac{11}{8}\nu^3,
\label{eq:app-gr-com-target-seed-l6}
\\[0.5ex]
 l_{10}^{\rm seed}&=\frac{47}{16}-\frac{235}{16}\nu+\frac{235}{16}\nu^2,
\label{eq:app-gr-com-target-seed-l10}
\\[0.5ex]
 l_{13}^{\rm seed}&=\frac{13}{8}-\frac{13}{2}\nu+\frac{13}{4}\nu^2,
\qquad
 l_{15}^{\rm seed}=-\frac18+\frac38\nu,
\label{eq:app-gr-com-target-seed-l13-l15}
\end{align}
with all other seed slots zero. Define the exact comparable-mass residual beyond
that seed by
\begin{equation}
 \Delta l_i\equiv l_i^{\rm GR}-l_i^{\rm seed}.
\label{eq:app-gr-com-target-delta-l-def}
\end{equation}
Then
\begin{equation}
\begin{aligned}
 \Delta l_1&=\frac{3\nu(3\nu-1)(4\nu-1)}{16},
 & \Delta l_2&=\Delta l_3=\Delta l_4=\Delta l_5=0.
\end{aligned}
\label{eq:app-gr-com-target-delta-l-1-5}
\end{equation}
\begin{equation}
\begin{aligned}
 \Delta l_6&=\frac{\nu(38-116\nu-57\nu^2)}{16},
 & \Delta l_7&=\frac{\nu^2(20-69\nu)}{16},
\\[0.5ex]
 \Delta l_8&=\frac{3\nu^2(3-11\nu)}{16},
 & \Delta l_9&=\frac{5\nu^3}{16}.
\end{aligned}
\label{eq:app-gr-com-target-delta-l-6-9}
\end{equation}
\begin{equation}
\begin{aligned}
 \Delta l_{10}&=\frac{\nu(129-98\nu+52\nu^2)}{16},
 & \Delta l_{11}&=\frac{\nu(-3+52\nu+124\nu^2)}{16},
\\[0.5ex]
 \Delta l_{12}&=\frac{\nu(-5+11\nu+48\nu^2)}{12}.
\end{aligned}
\label{eq:app-gr-com-target-delta-l-10-12}
\end{equation}
\begin{equation}
 \Delta l_{13}=-\frac{\nu(244+3\pi^2+1272\nu+96\nu^2)}{192}.
\label{eq:app-gr-com-target-delta-l-13}
\end{equation}
\begin{equation}
\begin{aligned}
 \Delta l_{14}&=\frac{\nu(452+3\pi^2-384\nu-224\nu^2)}{64},
 & \Delta l_{15}&=\frac{\nu(-908+63\pi^2)}{96}.
\end{aligned}
\label{eq:app-gr-com-target-delta-l-14-15}
\end{equation}
Because the COM compiler is diagonal with the same lower-order feedback for both
seed and target, the residual obeys the exact slotwise identity
\begin{equation}
 \Delta l_i=-\Delta h_i.
\label{eq:app-gr-com-target-delta-l-equals-minus-delta-h}
\end{equation}
This is the compact way to remember the entire COM target-import step: the
ordinary center-of-mass problem is already completely solved, and the remaining
content is a purely comparable-mass $15$-slot data vector. In the later
appendices this vector is what gets projected into the compact generic-frame
$24/17/8/1$ scaffold and, after that, into the grouped real $P_2$ inverse
problem.

% Appendix content for 4d_3pn.tex

\section{Generic-frame residual scaffold}
\label{app:generic-scaffold}

Once the exact one-body $3$PN gate, the natural bodywise self/static seed, and
 the reduced COM target are frozen, the next task is to state the honest
 generic-frame comparable-mass interaction scaffold before any contact/gauge
 quotient is imposed.  The guiding choice is deliberately conservative.  We do
 not try to guess the smallest basis too early.  Instead, we start from an
 exchange-symmetric \emph{overcomplete} residual scaffold that is guaranteed to
 preserve the exact one-body gate and then postpone every reduction of that
 scaffold until after the target Hamiltonian chart has been imported.

Write
\begin{equation}
\mathbf r\equiv \mathbf x_A-\mathbf x_B,
\qquad
r\equiv |\mathbf r|,
\qquad
\mathbf n\equiv \frac{\mathbf r}{r},
\label{eq:app-generic-scaffold-rn}
\end{equation}
and introduce the standard generic-frame invariants
\begin{equation}
 a=v_A^2,
 \qquad
 b=v_B^2,
 \qquad
 c=\mathbf v_A\!\cdot\!\mathbf v_B,
 \qquad
 d=v_{An}=\mathbf v_A\!\cdot\!\mathbf n,
 \qquad
 e=v_{Bn}=\mathbf v_B\!\cdot\!\mathbf n,
\label{eq:app-generic-scaffold-invariants}
\end{equation}
together with the mass labels
\begin{equation}
 p\equiv m_A,
 \qquad
 q\equiv m_B.
\label{eq:app-generic-scaffold-masses}
\end{equation}
The scaffold is symmetrized under body exchange,
\begin{equation}
A\leftrightarrow B:
\qquad
(a,b,d,e,p,q)\mapsto(b,a,e,d,q,p),
\qquad
c\mapsto c,
\label{eq:app-generic-scaffold-exchange}
\end{equation}
and every retained basis element is required to vanish in the strict test-mass
 branch
\begin{equation}
 p\to 0,
 \qquad
 q\to 1,
 \qquad
 b\to 0,
 \qquad
 c\to 0,
 \qquad
 e\to 0.
\label{eq:app-generic-scaffold-testmass-branch}
\end{equation}
That last condition is what keeps the exact one-body gate cleanly separated from
 the genuine comparable-mass lift.

The resulting interaction scaffold has block sizes
\begin{equation}
24,
\qquad
17,
\qquad
8,
\qquad
1,
\label{eq:app-generic-scaffold-block-counts}
\end{equation}
for the $G/r$, $G^2/r^2$, $G^3/r^3$, and $G^4/r^4$ blocks, respectively, so the
 full generic-frame $3$PN interaction ledger contains $50$ directions before any
 later quotient by COM blindness or ordinary contact freedom.

\subsection{\texorpdfstring{$G/r$}{G/r} sextic block}
\label{app:generic-scaffold-gr}

At order $G/r$ the new comparable-mass residual is an exchange-symmetric sextic
 polynomial in the invariants \eqref{eq:app-generic-scaffold-invariants}.  The
 overcomplete sextic basis has $24$ elements,
\begin{align}
Q_{01}&=a^2b+ab^2,
&
Q_{02}&=a^2c+b^2c,
&
Q_{03}&=a^2de+b^2de,
&
Q_{04}&=a^2e^2+b^2d^2,
\label{eq:app-generic-scaffold-Q-1-4}
\\[0.5ex]
Q_{05}&=abc,
&
Q_{06}&=abd^2+abe^2,
&
Q_{07}&=abde,
&
Q_{08}&=ac^2+bc^2,
\label{eq:app-generic-scaffold-Q-5-8}
\\[0.5ex]
Q_{09}&=acd^2+bce^2,
&
Q_{10}&=acde+bcde,
&
Q_{11}&=ace^2+bcd^2,
&
Q_{12}&=ad^2e^2+bd^2e^2,
\label{eq:app-generic-scaffold-Q-9-12}
\\[0.5ex]
Q_{13}&=ad^3e+bde^3,
&
Q_{14}&=ade^3+bd^3e,
&
Q_{15}&=ae^4+bd^4,
&
Q_{16}&=c^2d^2+c^2e^2,
\label{eq:app-generic-scaffold-Q-13-16}
\\[0.5ex]
Q_{17}&=c^2de,
&
Q_{18}&=c^3,
&
Q_{19}&=cd^2e^2,
&
Q_{20}&=cd^3e+cde^3,
\label{eq:app-generic-scaffold-Q-17-20}
\\[0.5ex]
Q_{21}&=cd^4+ce^4,
&
Q_{22}&=d^3e^3,
&
Q_{23}&=d^4e^2+d^2e^4,
&
Q_{24}&=d^5e+de^5.
\label{eq:app-generic-scaffold-Q-21-24}
\end{align}
So the corresponding residual contribution is
\begin{equation}
\frac{Gm_A m_B}{r}\sum_{i=1}^{24} q_i Q_i.
\label{eq:app-generic-scaffold-gr-residual}
\end{equation}
This is already broad enough to span the later reduced COM image of the
 velocity-six interaction block; later appendices only remove redundancy, not
 missing structure.

\subsection{\texorpdfstring{$G^2/r^2$}{G squared/r squared} quartic block}
\label{app:generic-scaffold-g2}

At order $G^2/r^2$ one extra mass power must be carried explicitly, so the
 comparable-mass residual is organized in an exchange-symmetric quartic
 mass-weighted basis with $17$ elements,
\begin{align}
T_{01}&=a^2p+b^2q,
&
T_{02}&=abp+abq,
&
T_{03}&=acp+bcq,
&
T_{04}&=acq+bcp,
\label{eq:app-generic-scaffold-T-1-4}
\\[0.5ex]
T_{05}&=ad^2p+be^2q,
&
T_{06}&=adep+bdeq,
&
T_{07}&=adeq+bdep,
&
T_{08}&=ae^2p+bd^2q,
\label{eq:app-generic-scaffold-T-5-8}
\\[0.5ex]
T_{09}&=ae^2q+bd^2p,
&
T_{10}&=c^2p+c^2q,
&
T_{11}&=cd^2p+ce^2q,
&
T_{12}&=cd^2q+ce^2p,
\label{eq:app-generic-scaffold-T-9-12}
\\[0.5ex]
T_{13}&=cdep+cdeq,
&
T_{14}&=d^2e^2p+d^2e^2q,
&
T_{15}&=d^3ep+de^3q,
\label{eq:app-generic-scaffold-T-13-15}
\\[0.5ex]
T_{16}&=d^3eq+de^3p,
&
T_{17}&=d^4p+e^4q.
\label{eq:app-generic-scaffold-T-16-17}
\end{align}
The quartic residual block is therefore
\begin{equation}
\frac{G^2m_A m_B}{r^2}\sum_{j=1}^{17} t_j T_j.
\label{eq:app-generic-scaffold-g2-residual}
\end{equation}
This is the first block in which the later seed-obstruction analysis becomes
 nontrivial.  The basis is intentionally broader than the final solved
 representative because the purpose of the present appendix is to preserve the
 honest algebraic search space before the middle-block repair is known.

\subsection{\texorpdfstring{$G^3/r^3$}{G cubed/r cubed} quadratic block}
\label{app:generic-scaffold-g3}

At order $G^3/r^3$ the residual is quadratic in the velocity invariants and
 carries two extra explicit mass powers.  The exchange-symmetric basis has eight
 elements,
\begin{align}
S_{01}&=ap^2+bq^2,
&
S_{02}&=apq+bpq,
&
S_{03}&=cp^2+cq^2,
&
S_{04}&=cpq,
\label{eq:app-generic-scaffold-S-1-4}
\\[0.5ex]
S_{05}&=d^2p^2+e^2q^2,
&
S_{06}&=d^2pq+e^2pq,
&
S_{07}&=dep^2+deq^2,
&
S_{08}&=depq.
\label{eq:app-generic-scaffold-S-5-8}
\end{align}
So the quadratic residual block is
\begin{equation}
\frac{G^3m_A m_B}{r^3}\sum_{k=1}^{8} s_k S_k.
\label{eq:app-generic-scaffold-g3-residual}
\end{equation}
This is the other block in which the later COM projection reveals a genuine
 middle-block obstruction relative to the naive seed placement.  Again, that
 obstruction is a bookkeeping issue, not an indication that the present block is
 too small.

\subsection{\texorpdfstring{$G^4/r^4$}{G to the 4/r to the 4} static block}
\label{app:generic-scaffold-g4}

At order $G^4/r^4$ there is only one exchange-symmetric static cross polynomial,
 namely
\begin{equation}
U_{01}=p^2q+pq^2.
\label{eq:app-generic-scaffold-U01}
\end{equation}
So the static residual block is simply
\begin{equation}
\frac{G^4m_A m_B}{r^4} u_1 U_{01}.
\label{eq:app-generic-scaffold-g4-residual}
\end{equation}
Combining all four interaction blocks, the honest generic-frame comparable-mass
 scaffold is
\begin{equation}
L_{3,\mathrm{full}}=L_{3,\mathrm{seed}}+L_{3,\mathrm{res}},
\label{eq:app-generic-scaffold-full-split}
\end{equation}
with
\begin{equation}
L_{3,\mathrm{res}}
=
\frac{Gm_A m_B}{r}\sum_{i=1}^{24} q_i Q_i
+\frac{G^2m_A m_B}{r^2}\sum_{j=1}^{17} t_j T_j
+\frac{G^3m_A m_B}{r^3}\sum_{k=1}^{8} s_k S_k
+\frac{G^4m_A m_B}{r^4} u_1 U_{01}.
\label{eq:app-generic-scaffold-full-residual}
\end{equation}
The point of \eqref{eq:app-generic-scaffold-full-residual} is not that the final
 $3$PN answer needs all $50$ coefficients.  It does not.  The point is that the
 full generic-frame lift should begin from a scaffold that is manifestly
 exchange-symmetric, manifestly compatible with the exact one-body gate, and
 manifestly finite before any later quotient by seed repair, COM blindness, or
 contact/gauge reduction is imposed.  Later appendices will show that once the
 fixed ADM Hamiltonian chart is used, the remaining algebraic freedom collapses
 dramatically; but the right place to impose that collapse is later, not here.

% Appendix content for 4d_3pn.tex

\section{COM projection matrices, seed obstruction, and seed repair}
\label{app:com-projection-seed-repair}

The reduced COM target of the previous appendices is already exact, and the
exchange-symmetric generic-frame scaffold from the comparable-mass setup is also
already finite.  The next honest question is therefore not whether a lift
exists in principle, but how much of that lift is already fixed by center-of-mass
projection and where the first genuine mismatch occurs.  This appendix records
that first exact answer.

The result is both cleaner and smaller than one might have expected.  COM
projection fixes a large part of the generic-frame interaction lift
immediately.  The first mismatch is not in the one-body gate, not in the
$G/r$ sextic block, and not in the static $G^4/r^4$ slot.  It sits entirely in
small $\nu^3$ tails in the middle $G^2/r^2$ and $G^3/r^3$ blocks, and it is
removed exactly by a tiny $\nu$-dressed seed correction that vanishes in the
strict one-body limit.  So the first generic-frame $3$PN obstruction is a
placement/bookkeeping issue, not a failure of the full scaffold.

\subsection{COM reduction maps and the blockwise projection matrices}
\label{app:com-projection-maps}

Use the generic-frame invariants
\begin{equation}
 a=v_A^2,
 \qquad
 b=v_B^2,
 \qquad
 c=\mathbf v_A\!\cdot\!\mathbf v_B,
 \qquad
 d=v_{An},
 \qquad
 e=v_{Bn},
\label{eq:app-com-proj-abcde}
\end{equation}
with masses
\begin{equation}
 p=m_A,
 \qquad
 q=m_B,
 \qquad
 M=p+q,
 \qquad
 \nu\equiv\frac{pq}{M^2}.
\label{eq:app-com-proj-pqnu}
\end{equation}
In the COM frame,
\begin{equation}
\mathbf v_A=X_B\mathbf v,
\qquad
\mathbf v_B=-X_A\mathbf v,
\qquad
X_A+X_B=1,
\qquad
\nu=X_AX_B,
\label{eq:app-com-proj-XAXB}
\end{equation}
so
\begin{equation}
 a=X_B^2v^2,
 \qquad
 b=X_A^2v^2,
 \qquad
 c=-\nu v^2,
\label{eq:app-com-proj-abc}
\end{equation}
\begin{equation}
 d=X_B\dot r,
 \qquad
 e=-X_A\dot r.
\label{eq:app-com-proj-de}
\end{equation}
Therefore each exchange-symmetric generic-frame basis monomial has a definite
image in the reduced COM slot basis.

Write the interaction scaffold, after subtraction of the carried one-body/self-static
seed, as
\begin{equation}
L_{3,\rm res}
=
\frac{Gm_A m_B}{r}\sum_{i=1}^{24} q_i Q_i
+\frac{G^2m_A m_B}{r^2}\sum_{j=1}^{17} t_j T_j
+\frac{G^3m_A m_B}{r^3}\sum_{k=1}^{8} s_k S_k
+\frac{G^4m_A m_B}{r^4}u_1 U_{01},
\label{eq:app-com-proj-residual-scaffold}
\end{equation}
where $\{Q_i\}$, $\{T_j\}$, $\{S_k\}$, and $U_{01}$ are the compact
exchange-symmetric basis elements introduced in the generic-frame scaffold
appendix.  Let
\begin{equation}
\mathbf q=(q_1,\dots,q_{24})^T,
\qquad
\mathbf t=(t_1,\dots,t_{17})^T,
\qquad
\mathbf s=(s_1,\dots,s_8)^T.
\label{eq:app-com-proj-coefficient-vectors}
\end{equation}
Then the COM reduction induces exact coefficient-extraction maps
\begin{equation}
\Pi_Q(\mathbf q)=M_Q\mathbf q,
\qquad
M_Q\in\mathrm{Mat}_{12\times24}(\mathbb Q),
\label{eq:app-com-proj-MQ}
\end{equation}
\begin{equation}
\Pi_T(\mathbf t)=M_T\mathbf t,
\qquad
M_T\in\mathrm{Mat}_{6\times17}(\mathbb Q),
\label{eq:app-com-proj-MT}
\end{equation}
\begin{equation}
\Pi_S(\mathbf s)=M_S\mathbf s,
\qquad
M_S\in\mathrm{Mat}_{4\times8}(\mathbb Q),
\label{eq:app-com-proj-MS}
\end{equation}
\begin{equation}
\Pi_U(u_1)=M_U u_1,
\qquad
M_U\in\mathrm{Mat}_{1\times1}(\mathbb Q[\pi^2]).
\label{eq:app-com-proj-MU}
\end{equation}
Here the entries of the projection matrices are defined by coefficient extraction
in the reduced COM slots:
\begin{equation}
\Pi_Q(\mathbf q)
\equiv
\bigl([
u]\delta l_6,[\nu^2]\delta l_6,[\nu^3]\delta l_6,\dots,[\nu^3]\delta l_9\bigr)^T,
\label{eq:app-com-proj-PiQ}
\end{equation}
\begin{equation}
\Pi_T(\mathbf t)
\equiv
\bigl([
u]\delta l_{10},[\nu^2]\delta l_{10},[\nu]\delta l_{11},[\nu^2]\delta l_{11},[\nu]\delta l_{12},[\nu^2]\delta l_{12}\bigr)^T,
\label{eq:app-com-proj-PiT}
\end{equation}
\begin{equation}
\Pi_S(\mathbf s)
\equiv
\bigl([
u]\delta l_{13},[\nu^2]\delta l_{13},[\nu]\delta l_{14},[\nu^2]\delta l_{14}\bigr)^T,
\label{eq:app-com-proj-PiS}
\end{equation}
\begin{equation}
\Pi_U(u_1)
\equiv
\bigl([
u]\delta l_{15}\bigr).
\label{eq:app-com-proj-PiU}
\end{equation}
So the phrase ``COM projection matrix'' means the exact linear map that sends a
block of generic-frame coefficients to the polynomial coefficients of the
reduced COM slots.

The audit establishes the exact ranks
\begin{equation}
\operatorname{rank}M_Q=12,
\qquad
\operatorname{rank}M_T=6,
\qquad
\operatorname{rank}M_S=4,
\qquad
\operatorname{rank}M_U=1,
\label{eq:app-com-proj-ranks}
\end{equation}
with nullities
\begin{equation}
\dim\ker M_Q=12,
\qquad
\dim\ker M_T=11,
\qquad
\dim\ker M_S=4,
\qquad
\dim\ker M_U=0.
\label{eq:app-com-proj-nullities}
\end{equation}
Equivalently, the total COM-blind interaction freedom is
\begin{equation}
\boxed{12+11+4+0=27.}
\label{eq:app-com-proj-total-nullity}
\end{equation}
In the compact basis ordering used by the audit, the pivot columns may be taken
as
\begin{equation}
Q:\{1,2,3,4,5,7,12,13,14,22,23,24\},
\label{eq:app-com-proj-Q-pivots}
\end{equation}
\begin{equation}
T:\{1,2,5,6,14,16\},
\qquad
S:\{1,2,5,6\},
\label{eq:app-com-proj-TS-pivots}
\end{equation}
while the $U$ block is of course already unique.  These pivot data are not used
later as physical objects; they are simply a compact way to say that the COM
image is already completely explicit before any gauge/contact quotient is taken.

The structural interpretation of \eqref{eq:app-com-proj-ranks} is immediate.
The $Q$ block can support $\nu$, $\nu^2$, and $\nu^3$ in each of its four COM
slots; the $T$ block supports only $\nu$ and $\nu^2$ in each of its three slots;
the $S$ block likewise supports only $\nu$ and $\nu^2$ in each of its two slots,
with coefficients allowed in $\mathbb Q(\pi^2)$; and the $U$ block reduces only
to a coefficient proportional to $\nu$.  This is exactly the mass-degree fact
that produces the first obstruction.

\subsection{The exact middle-block obstruction relative to the natural seed}
\label{app:com-proj-obstruction}

Subtract the natural $3$PN one-body/self-static seed from the exact COM target.
The $G/r$ sextic block and the $G^4/r^4$ static slot are already compatible with
the images of $M_Q$ and $M_U$.  The only incompatibility sits in the middle two
blocks.  Explicitly, the exact COM residuals are
\begin{equation}
\Delta l_{10}
=
\nu\left(\frac{129}{16}-\frac{98}{16}\nu\right)+\boxed{\frac{13}{4}\nu^3},
\label{eq:app-com-proj-dl10}
\end{equation}
\begin{equation}
\Delta l_{11}
=
\nu\left(-\frac{3}{16}+\frac{52}{16}\nu\right)+\boxed{\frac{31}{4}\nu^3},
\label{eq:app-com-proj-dl11}
\end{equation}
\begin{equation}
\Delta l_{12}
=
\nu\left(-\frac{5}{12}+\frac{11}{12}\nu\right)+\boxed{4\nu^3},
\label{eq:app-com-proj-dl12}
\end{equation}
for the $G^2/r^2$ block, while for the $G^3/r^3$ block one finds
\begin{equation}
\Delta l_{13}
=
\nu\left[-\frac{244+3\pi^2}{192}-\frac{1272}{192}\nu\right]+\boxed{-\frac12\nu^3},
\label{eq:app-com-proj-dl13}
\end{equation}
\begin{equation}
\Delta l_{14}
=
\nu\left[\frac{452+3\pi^2}{64}-6\nu\right]+\boxed{-\frac72\nu^3}.
\label{eq:app-com-proj-dl14}
\end{equation}
Because the $T$ and $S$ projection matrices carry only $\nu$ and $\nu^2$ support,
none of the boxed $\nu^3$ pieces in
\eqref{eq:app-com-proj-dl10}--\eqref{eq:app-com-proj-dl14} can arise from the
present constant-coefficient generic-frame middle blocks.  The exact obstruction
vector is therefore
\begin{equation}
\boxed{
\delta l_{10}^{\rm obs}=\frac{13}{4}\nu^3,
\qquad
\delta l_{11}^{\rm obs}=\frac{31}{4}\nu^3,
\qquad
\delta l_{12}^{\rm obs}=4\nu^3,}
\label{eq:app-com-proj-T-obstruction}
\end{equation}
\begin{equation}
\boxed{
\delta l_{13}^{\rm obs}=-\frac12\nu^3,
\qquad
\delta l_{14}^{\rm obs}=-\frac72\nu^3.}
\label{eq:app-com-proj-S-obstruction}
\end{equation}
This is the first exact generic-frame lift obstruction.  The natural seed is
already exact in the strict one-body sector, but it is not yet aligned with the
mass-degree structure of the middle interaction blocks.

\subsection{Minimal \texorpdfstring{$\nu$}{nu}-dressed seed repair}
\label{app:com-proj-seed-repair}

The obstruction is repaired by a remarkably small correction sector.  Define
\begin{equation}
\boxed{
\Delta T_\nu
=
\frac{pq}{4(p+q)}\Bigl(13ab+31cde+16d^2e^2\Bigr),}
\label{eq:app-com-proj-DeltaTnu}
\end{equation}
\begin{equation}
\boxed{
\Delta S_\nu
=
\frac{p^2q^2}{2(p+q)^2}\Bigl(c+7de\Bigr).}
\label{eq:app-com-proj-DeltaSnu}
\end{equation}
These terms have exactly the two properties one wants.

First, they vanish in the strict one-body limit.  On the branch
\begin{equation}
 m_A\to 0,
 \qquad
 \mathbf v_B\to 0,
 \qquad
 p\to 0,
 \qquad
 b,c,e\to 0,
\label{eq:app-com-proj-testmass-branch}
\end{equation}
one has
\begin{equation}
\boxed{\Delta T_\nu\to 0,
\qquad
\Delta S_\nu\to 0.}
\label{eq:app-com-proj-repair-onebody}
\end{equation}
So the seed repair never touches the exact one-body gate.

Second, under COM reduction these terms reproduce precisely the missing
obstruction tails.  In slot language,
\begin{equation}
\Delta T_\nu\stackrel{\rm COM}{\longmapsto}
\left(\frac{13}{4}\nu^3,\ \frac{31}{4}\nu^3,\ 4\nu^3\right)
\quad\text{in }(l_{10},l_{11},l_{12}),
\label{eq:app-com-proj-DeltaTnu-image}
\end{equation}
\begin{equation}
\Delta S_\nu\stackrel{\rm COM}{\longmapsto}
\left(-\frac12\nu^3,\ -\frac72\nu^3\right)
\quad\text{in }(l_{13},l_{14}).
\label{eq:app-com-proj-DeltaSnu-image}
\end{equation}
So the refined seed is simply
\begin{equation}
\boxed{
L_{3,\rm seed}^{\rm ref}
=
L_{3,\rm seed}
+\frac{G^2pq}{r^2}\,\Delta T_\nu
+\frac{G^3pq}{r^3}\,\Delta S_\nu.}
\label{eq:app-com-proj-refined-seed}
\end{equation}
This is the exact $\nu$-dressed seed repair advertised in the main text.  It is
minimal in the literal sense: it touches only the obstructed middle blocks, it
vanishes in the strict one-body limit, and it removes all of the COM mismatch.

\subsection{One exact repaired representative in each interaction block}
\label{app:com-proj-representatives}

After the refined seed is frozen, each block admits an exact generic-frame
representative with the correct COM image.  The representative is not unique ---
that is the subject of the next appendix --- but existence is now explicit.
Using the compact basis monomials from the generic-frame scaffold appendix, one
may take
\begin{equation}
L_{Q,\rm part}
=
\frac{Gm_A m_B}{r}
\Bigl[
\frac94 Q_1
-\frac{19}{8}Q_2
+\frac54 Q_4
+\frac{61}{16}Q_5
+\frac{29}{16}Q_7
+\frac{9}{16}Q_{12}
+\frac{15}{32}Q_{14}
-\frac{5}{16}Q_{22}
\Bigr],
\label{eq:app-com-proj-Qpart}
\end{equation}
\begin{equation}
\begin{aligned}
L_{T,\rm part}
={}&
\frac{G^2m_A m_B}{r^2}
\Bigl[
\frac{129}{16}(a^2p+b^2q)
+\frac{289}{16}(abp+abq)
-\frac{3}{16}(ad^2p+be^2q)
\Bigr]
\\
&
+\frac{G^2m_A m_B}{r^2}
\Bigl[
-\frac{43}{16}(adep+bdeq)
-\frac13(d^2e^2p+d^2e^2q)
+\frac{5}{12}(d^3eq+de^3p)
\Bigr],
\end{aligned}
\label{eq:app-com-proj-Tpart}
\end{equation}
\begin{equation}
\begin{aligned}
L_{S,\rm part}
={}&
\frac{G^3m_A m_B}{r^3}
\Bigl[
-\frac{880+3\pi^2}{192}(ap^2+bq^2)
-\frac{244+3\pi^2}{192}(apq+bpq)
\Bigr]
\\
&
+\frac{G^3m_A m_B}{r^3}
\Bigl[
\frac{260+3\pi^2}{64}(d^2p^2+e^2q^2)
+\frac{452+3\pi^2}{64}(d^2pq+e^2pq)
\Bigr],
\end{aligned}
\label{eq:app-com-proj-Spart}
\end{equation}
and the unique static block
\begin{equation}
L_{U,\rm part}
=
\frac{G^4m_A m_B}{r^4}
\frac{63\pi^2-908}{96}(p^2q+pq^2).
\label{eq:app-com-proj-Upart}
\end{equation}
By construction, the COM reductions of
\eqref{eq:app-com-proj-Qpart}--\eqref{eq:app-com-proj-Upart} reproduce the exact
refined targets in every block.  The important point is not that these formulas
are especially canonical --- they are not --- but that, after the seed repair,
blockwise COM compatibility is completely explicit.

\subsection{What the projection stage has fixed --- and what it has not}
\label{app:com-proj-ledger}

The full content of the present appendix can now be summarized in one line:
\begin{equation}
\boxed{
\begin{aligned}
&\text{COM projection plus the minimal $\nu$-dressed seed repair removes every}\\
&\text{blockwise compatibility obstruction, but it still leaves a }27\text{-dimensional}\\
&\text{COM-blind family.}
\end{aligned}}
\label{eq:app-com-proj-summary-line}
\end{equation}
Equivalently,
\begin{equation}
Q:\ 24\to 12,
\qquad
T:\ 17\to 6,
\qquad
S:\ 8\to 4,
\qquad
U:\ 1\to 1,
\label{eq:app-com-proj-block-map-summary}
\end{equation}
so after repair the unfixed interaction nullities remain
\begin{equation}
12,\qquad 11,\qquad 4,\qquad 0,
\label{eq:app-com-proj-nullity-summary}
\end{equation}
for a total of
\begin{equation}
\boxed{27}
\label{eq:app-com-proj-total-nullity-repeat}
\end{equation}
COM-blind directions.

So the present appendix does not finish the full generic-frame theorem.  What it
finishes is the exact COM-constraint layer.  It shows that
\begin{enumerate}
  \item the $Q$ block was already compatible,
  \item the $U$ block is already unique once the $\pi^2$ coefficient is admitted,
  \item the only true mismatch sat in five tiny $\nu^3$ tails in the middle
  blocks,
  \item those tails are absorbed exactly by the refined seed
  \eqref{eq:app-com-proj-refined-seed}, and
  \item after that repair one explicit COM-compatible representative exists in
  every interaction block.
\end{enumerate}
What remains is therefore no longer a compatibility problem.  It is the
genuine COM-null\allowbreak/\allowbreak gauge\allowbreak/\allowbreak contact
quotient problem taken up in the next appendix.

% Appendix content for 4d_3pn.tex

\section{The COM-null ideal and canonical quotient slice}
\label{app:com-null}

After the exact $\nu$-dressed seed repair, the generic-frame $3$PN lift is no
longer obstructed by any middle-block mismatch.  What remains is a quotient
problem.  The repaired $24/17/8/1$ interaction scaffold still contains
nontrivial generic-frame directions that vanish identically after COM
reduction, and the next honest task is to identify those directions exactly and
choose a convenient representative of their quotient class.

This appendix records that algebraic structure.  The main conclusion is simple.
The remaining unfixed family is not unexplained physics.  It is precisely the
COM-blind family generated by the standard COM identities.  Once that family is
recognized as an ideal, it becomes natural to choose a canonical sparse
representative by quotienting it away.  That choice is exact algebraically and
already sufficient to write down a clean generic-frame representative of the
full solved COM target.  What it is \emph{not} yet is the final ordinary
gauge/contact theorem; that sharper statement is postponed to the next
appendix.

Continue to write the repaired generic-frame interaction scaffold as
\begin{equation}
L_{3,\rm res}
=
\frac{Gpq}{r}\sum_{i=1}^{24} q_i Q_i
+\frac{G^2pq}{r^2}\sum_{j=1}^{17} t_j T_j
+\frac{G^3pq}{r^3}\sum_{k=1}^{8} s_k S_k
+\frac{G^4pq}{r^4}u_1 U_{01},
\label{eq:app-com-null-residual-scaffold}
\end{equation}
where the compact exchange-symmetric basis elements $Q_i$, $T_j$, $S_k$, and
$U_{01}$ were fixed in the generic-frame scaffold appendix, and where the
coefficient vectors are
\begin{equation}
\mathbf q=(q_1,\dots,q_{24})^T,
\qquad
\mathbf t=(t_1,\dots,t_{17})^T,
\qquad
\mathbf s=(s_1,\dots,s_8)^T.
\label{eq:app-com-null-coeff-vectors}
\end{equation}
The previous appendix showed that after seed repair the COM image is completely
fixed block by block.  The present appendix identifies exactly what survives in
the kernel of that projection.

\subsection{The COM-null ideal}
\label{app:com-null-ideal}

In the COM frame the body velocities satisfy three linear-momentum identities,
\begin{equation}
\boxed{
\mathcal C_1\equiv pa+qc=0,
\qquad
\mathcal C_2\equiv pc+qb=0,
\qquad
\mathcal C_3\equiv pd+qe=0,}
\label{eq:app-com-null-C123}
\end{equation}
and three collinearity identities,
\begin{equation}
\boxed{
\mathcal C_4\equiv ab-c^2=0,
\qquad
\mathcal C_5\equiv ae-cd=0,
\qquad
\mathcal C_6\equiv bd-ce=0.}
\label{eq:app-com-null-C456}
\end{equation}
These generate the COM-null ideal
\begin{equation}
\boxed{
\mathfrak I_{\rm COM}
=
\langle \mathcal C_1,\mathcal C_2,\mathcal C_3,\mathcal C_4,\mathcal C_5,\mathcal C_6\rangle.}
\label{eq:app-com-null-ideal-def}
\end{equation}
It is also convenient to name the smaller linear-momentum ideal
\begin{equation}
\mathfrak I_{\rm lin}
=
\langle \mathcal C_1,\mathcal C_2,\mathcal C_3\rangle.
\label{eq:app-com-null-linear-ideal}
\end{equation}
Then the full $27$-dimensional COM-blind family is exactly the generic-frame
freedom invisible in the quotient algebra
\begin{equation}
\mathbb Q[a,b,c,d,e,p,q]/\mathfrak I_{\rm COM}.
\label{eq:app-com-null-quotient-algebra}
\end{equation}
A useful refinement already visible in the referee audit is that the $T$- and
$S$-block null directions lie in the smaller ideal $\mathfrak I_{\rm lin}$,
whereas the $Q$-block null directions require the full ideal
$\mathfrak I_{\rm COM}$.

The sparse null generators returned by the exact COM-projection audit may be
chosen as follows.

\paragraph{T-block null generators.}
A sparse basis of $\ker M_T$ is
{\small
\begin{align*}
N_{T1}&=abp+abq+acp+bcq,
&
N_{T2}&=a^2p+acq+b^2q+bcp,
\\
N_{T3}&=ad^2p+adeq+bdep+be^2q,
&
N_{T4}&=adep+ae^2p+bd^2q+bdeq,
\\
N_{T5}&=adep+ae^2q+bd^2p+bdeq,
&
N_{T6}&=-abp-abq+c^2p+c^2q,
\\
N_{T7}&=-adep-bdeq+cd^2p+ce^2q,
&
N_{T8}&=ad^2p+be^2q+cd^2q+ce^2p,
\\
N_{T9}&=adep+bdeq+cdep+cdeq,
&
N_{T10}&=d^3ep+d^2e^2p+d^2e^2q+de^3q,
\\
N_{T11}&=d^4p+d^3eq+de^3p+e^4q.
\end{align*}}
The exact audit verifies
\begin{equation}
N_{Ti}\in\mathfrak I_{\rm lin}
\qquad (i=1,\dots,11).
\label{eq:app-com-null-T-membership}
\end{equation}

\paragraph{S-block null generators.}
A sparse basis of $\ker M_S$ is
{\small
\begin{align*}
N_{S1}&=apq+bpq+cp^2+cq^2,
&
N_{S2}&=\frac12 a p^2+\frac12 b q^2+cpq,
\\
N_{S3}&=d^2pq+dep^2+deq^2+e^2pq,
&
N_{S4}&=\frac12 d^2p^2+depq+\frac12 e^2q^2.
\end{align*}}
Again the audit gives
\begin{equation}
N_{Si}\in\mathfrak I_{\rm lin}
\qquad (i=1,\dots,4).
\label{eq:app-com-null-S-membership}
\end{equation}

\paragraph{Q-block null generators.}
A sparse basis of $\ker M_Q$ is
{\small
\begin{align*}
N_{Q1}&=-a^2e^2+ab d^2+ab e^2-b^2d^2,
&
N_{Q2}&=-a^2b-ab^2+ac^2+bc^2,
\\
N_{Q3}&=-a^2de+ac d^2-b^2de+bc e^2,
&
N_{Q4}&=-a^2e^2+acde-b^2d^2+bcde,
\\
N_{Q5}&=-2abde+ac e^2+bc d^2,
&
N_{Q6}&=-ad^2e^2+ae^4+bd^4-bd^2e^2,
\\
N_{Q7}&=-a^2e^2-b^2d^2+c^2d^2+c^2e^2,
&
N_{Q8}&=-abde+c^2de,
\\
N_{Q9}&=-abc+c^3,
&
N_{Q10}&=-\frac12 ad e^3-\frac12 bd^3e+cd^2e^2,
\\
N_{Q11}&=-ad^2e^2-bd^2e^2+cd^3e+cde^3,
&
N_{Q12}&=-ad^3e-bde^3+cd^4+ce^4.
\end{align*}}
For these generators the exact statement is only
\begin{equation}
N_{Qi}\in\mathfrak I_{\rm COM}
\qquad (i=1,\dots,12),
\label{eq:app-com-null-Q-membership}
\end{equation}
not membership in the smaller linear ideal.

The conceptual consequence of
\eqref{eq:app-com-null-T-membership}--\eqref{eq:app-com-null-Q-membership} is
important.  The entire COM-blind family is generated by exact algebraic COM
identities, but only the $T$ and $S$ kernels are already linear-momentum-null.
The $Q$ block genuinely needs the collinearity relations as well.

\subsection{What COM fixes and what it does not fix}
\label{app:com-null-fixes}

After the minimal $\nu$-dressed repair, COM matching determines one exact image
in each interaction block, but it does not fix the full generic-frame lift.
Block by block, the exact nullities are
\begin{equation}
\boxed{
\dim\ker M_Q=12,
\qquad
\dim\ker M_T=11,
\qquad
\dim\ker M_S=4,
\qquad
\dim\ker M_U=0.}
\label{eq:app-com-null-block-nullities}
\end{equation}
Equivalently, the total COM-blind interaction freedom is
\begin{equation}
\boxed{12+11+4+0=27.}
\label{eq:app-com-null-total-nullity}
\end{equation}
So COM reduction fixes the image layer but leaves a $27$-dimensional quotient
family unfixed.  In that precise sense, COM data alone determine the target
only modulo the ideal $\mathfrak I_{\rm COM}$.

This also clarifies what the previous appendix actually achieved.  The
$\nu$-dressed repair removed the only genuine image obstruction.  Once that is
done, the leftover freedom is entirely quotient freedom.  There is no remaining
middle-block compatibility failure hiding inside the $27$ dimensions of
\eqref{eq:app-com-null-total-nullity}; what remains is the choice of a
representative in a fixed quotient class.

\subsection{Canonical quotient slice}
\label{app:com-null-slice}

The natural sparse algebraic choice is therefore to quotient by the COM-null
ideal and set those $27$ null coefficients to zero.  Equivalently, one chooses
a complement to the nullspace of the COM projection matrices and keeps only that
complement.  This defines a canonical quotient slice: a distinguished algebraic
representative of the full COM equivalence class.

The usefulness of this choice is immediate.  It turns the abstract statement
``the generic-frame lift is determined only up to COM-blind directions'' into an
explicit sparse answer.  Write that canonical representative as
\begin{equation}
\boxed{
L_{3,\rm can}
=
L_{3,\rm seed}
+\frac{Gpq}{r}Q_{\rm can}
+\frac{G^2pq}{r^2}T_{\rm can}
+\frac{G^3pq}{r^3}S_{\rm can}
+\frac{G^4}{r^4}U_{\rm can}.}
\label{eq:app-com-null-L3can}
\end{equation}
By construction,
\begin{equation}
\Pi_Q(Q_{\rm can})=(\Delta l_6,\Delta l_7,\Delta l_8,\Delta l_9),
\qquad
\Pi_T(T_{\rm can})=(\Delta l_{10},\Delta l_{11},\Delta l_{12}),
\label{eq:app-com-null-projection-QT}
\end{equation}
\begin{equation}
\Pi_S(S_{\rm can})=(\Delta l_{13},\Delta l_{14}),
\qquad
\Pi_U(U_{\rm can})=(\Delta l_{15}).
\label{eq:app-com-null-projection-SU}
\end{equation}
The next four subsections simply record the corresponding block formulas.

\subsection{\texorpdfstring{$G/r$}{G/r} sextic block}
\label{app:com-null-slice-gr}

The canonical sextic representative may be taken as
\begin{align}
Q_{\rm can}={}&
\frac94(a^2b+ab^2)
-\frac{19}{8}(a^2c+b^2c)
+\frac54(a^2e^2+b^2d^2)
+\frac{61}{16}abc
+\frac{29}{16}abde
\notag\\
&
+\frac{9}{16}(ad^2e^2+bd^2e^2)
+\frac{15}{32}(ade^3+bd^3e)
-\frac{5}{16}d^3e^3.
\label{eq:app-com-null-Qcan}
\end{align}
This is the unique sparse representative of the sextic block in the chosen
quotient slice.

\subsection{\texorpdfstring{$G^2/r^2$}{G squared/r squared} quartic block}
\label{app:com-null-slice-g2}

The canonical quartic representative may be taken as
\begin{align}
T_{\rm can}={}&
\frac{129}{16}(a^2p+b^2q)
+\frac{289}{16}(abp+abq)
-\frac{3}{16}(ad^2p+be^2q)
-\frac{43}{16}(adep+bdeq)
\notag\\
&
-\frac13(d^2e^2p+d^2e^2q)
+\frac{5}{12}(d^3eq+de^3p)
+\frac{pq}{4(p+q)}\bigl(13ab+31cde+16d^2e^2\bigr).
\label{eq:app-com-null-Tcan}
\end{align}
The final term is precisely the compact way in which the $\nu$-dressed repair is
incorporated into the canonical quotient representative.

\subsection{\texorpdfstring{$G^3/r^3$}{G cubed/r cubed} quadratic block}
\label{app:com-null-slice-g3}

The canonical quadratic representative may be taken as
\begin{align}
S_{\rm can}={}&
-\frac{880+3\pi^2}{192}(ap^2+bq^2)
-\frac{244+3\pi^2}{192}(apq+bpq)
+\frac{260+3\pi^2}{64}(d^2p^2+e^2q^2)
\notag\\
&
+\frac{452+3\pi^2}{64}(d^2pq+e^2pq)
+\frac{p^2q^2}{2(p+q)^2}(c+7de).
\label{eq:app-com-null-Scan}
\end{align}
Again the last term is the compact generic-frame image of the repaired
middle-block placement.

\subsection{\texorpdfstring{$G^4/r^4$}{G to the 4/r to the 4} static block}
\label{app:com-null-slice-g4}

The static block is already unique once the COM target is fixed.  Its canonical
representative is therefore simply
\begin{equation}
\boxed{
U_{\rm can}=\frac{63\pi^2-908}{96}(p^2q+pq^2).}
\label{eq:app-com-null-Ucan}
\end{equation}
Substituting
\eqref{eq:app-com-null-Qcan}--\eqref{eq:app-com-null-Ucan} into
\eqref{eq:app-com-null-L3can} gives one explicit sparse generic-frame answer
whose COM reduction reproduces the exact solved $3$PN target in every slot.

\subsection{Status of the canonical slice}
\label{app:com-null-status}

The canonical quotient slice is an exact algebraic representative, and it is
extremely useful because it turns the whole COM-blind family into a sparse
explicit formula.  But it is still only a quotient choice.  Nothing in the
construction above says that every COM-blind direction is an ordinary contact
transformation, or even that the full $27$-dimensional quotient family is
physically gauge.

That distinction is the real point of this appendix.  The canonical slice gives
the right algebraic reference chart, but not yet the final physical uniqueness
theorem.  The next appendix sharpens the result by computing the actual
scaffold-preserving ordinary contact/gauge orbit and showing that it is only an
$11$-dimensional subspace of the larger $27$-dimensional COM-blind family.
Equivalently,
\begin{quote}
COM blindness and ordinary contact/gauge freedom are not identical.
\end{quote}
So Appendix~K is not a reworking of the same material.  It is the additional
step that turns the algebraic quotient picture into the true ordinary
contact/gauge statement.

% Appendix content for 4d_3pn.tex

\section{Contact generator and gauge-orbit derivation}
\label{app:gauge-orbit}

The previous appendix identified the full $27$-dimensional COM-blind family and
used it to define a canonical sparse quotient representative.  That algebraic
slice is exact and extremely useful, but by itself it does not yet answer the
sharper physical bookkeeping question: how much of that COM-blind freedom is
actually generated by ordinary $3$PN contact transformations?  The purpose of
this appendix is to answer that question exactly.

The result is stricter than the naive quotient picture.  Once the generic-frame
$3$PN interaction lift is required both to stay inside the compact
$24/17/8/1$ scaffold and to preserve the already solved COM target, the true
ordinary contact/gauge family is only $11$-dimensional.  So COM blindness and
ordinary contact freedom are not identical.  The earlier canonical quotient
slice is therefore the right sparse algebraic representative, but it is larger
than the actual ordinary gauge orbit.

Continue to write the repaired generic-frame residual scaffold as
\begin{equation}
L_{3,\rm res}
=
\frac{Gpq}{r}\sum_{i=1}^{24} q_i Q_i
+\frac{G^2pq}{r^2}\sum_{j=1}^{17} t_j T_j
+\frac{G^3pq}{r^3}\sum_{k=1}^{8} s_k S_k
+\frac{G^4pq}{r^4}u_1 U_{01},
\label{eq:app-gauge-orbit-residual-scaffold}
\end{equation}
with the compact exchange-symmetric scaffold basis fixed in the earlier
appendices.  In the invariant shorthand used throughout the generic-frame lift,
\begin{equation}
 a=v_A^2,
 \qquad
 b=v_B^2,
 \qquad
 c=\mathbf v_A\!\cdot\!\mathbf v_B,
 \qquad
 d=v_{An},
 \qquad
 e=v_{Bn},
\label{eq:app-gauge-orbit-abcde}
\end{equation}
while
\begin{equation}
 p=m_A,
 \qquad
 q=m_B,
 \qquad
 r=|\mathbf x_A-\mathbf x_B|.
\label{eq:app-gauge-orbit-pqr}
\end{equation}
The scaffold block sizes are
\begin{equation}
\boxed{\dim Q=24,
\qquad
\dim T=17,
\qquad
\dim S=8,
\qquad
\dim U=1.}
\label{eq:app-gauge-orbit-scaffold-sizes}
\end{equation}
The present appendix reconstructs the ordinary total-derivative orbit inside
that fixed scaffold.

\subsection{Raw odd-generator families}
\label{app:gauge-orbit-raw}

At $3$PN an ordinary contact transformation is generated by a total derivative
of an odd scalar.  The crucial exchange-parity rule differs from the even-scalar
rule used for the interaction scaffold, because $\mathbf n$ changes sign under
label interchange.  The correct odd rule is
\begin{equation}
A\leftrightarrow B,
\qquad
c\mapsto c,
\qquad
d\mapsto -e,
\qquad
e\mapsto -d,
\qquad
p\leftrightarrow q.
\label{eq:app-gauge-orbit-odd-exchange}
\end{equation}
The referee audit shows that the raw odd-scalar generator families relevant at
formal $3$PN are
\begin{equation}
 r\times(\text{degree-}7\ \text{odd scalars}),
 \qquad
 G\times(\text{degree-}5\ \text{odd scalars}),
\label{eq:app-gauge-orbit-raw-families-1}
\end{equation}
\begin{equation}
 \frac{G^2}{r}\times(\text{degree-}3\ \text{odd scalars}),
 \qquad
 \frac{G^3}{r^2}\times(\text{degree-}1\ \text{odd scalars}).
\label{eq:app-gauge-orbit-raw-families-2}
\end{equation}
Counting exchange-symmetric odd monomials with the strict test-mass-vanishing
rule gives the exact raw dimensions
\begin{equation}
\boxed{31,
\qquad
12,
\qquad
8,
\qquad
2,}
\label{eq:app-gauge-orbit-raw-counts}
\end{equation}
for a total raw odd family of
\begin{equation}
\boxed{31+12+8+2=53.}
\label{eq:app-gauge-orbit-total-raw}
\end{equation}
So before any scaffold or COM condition is imposed there are fifty-three
candidate odd generators whose Newtonian-order-reduced total derivatives could,
in principle, shift the ordinary $3$PN representative.

\subsection{Scaffold-preserving and COM-preserving reduction}
\label{app:gauge-orbit-reduction}

Not every total derivative generated from the raw odd family stays inside the
compact $24/17/8/1$ scaffold.  Let $F$ be any raw odd scalar and define the
induced ordinary shift after Newtonian order reduction by
\begin{equation}
\delta L_{\rm cont}[F]
\equiv
\frac{dF}{dt}\Big|_{\rm Newt\ red.}.
\label{eq:app-gauge-orbit-dFdt}
\end{equation}
Demanding that $\delta L_{\rm cont}[F]$ project entirely onto the compact
$Q/T/S/U$ scaffold imposes a linear constraint system of rank $31$ on the
$53$-dimensional raw family.  The resulting scaffold-preserving kernel therefore
has dimension
\begin{equation}
\boxed{53-31=22.}
\label{eq:app-gauge-orbit-scaffold-kernel}
\end{equation}
The second condition is that the contact shift must leave the solved COM target
unchanged.  Imposing zero COM image on that $22$-dimensional scaffold-preserving
family produces an exact COM image rank of $11$, so the actual ordinary
contact/gauge orbit has dimension
\begin{equation}
\boxed{22-11=11.}
\label{eq:app-gauge-orbit-contact-dim}
\end{equation}
Equivalently, the exact reduction chain is
\begin{equation}
\boxed{53\longrightarrow 22\longrightarrow 11.}
\label{eq:app-gauge-orbit-dimension-chain}
\end{equation}
This is the sharp bookkeeping statement that the previous appendix did not yet
make explicit: the full COM-blind quotient is larger than the true ordinary
contact family.

\subsection{Sparse basis for the actual \texorpdfstring{$11$}{11}-generator orbit}
\label{app:gauge-orbit-basis}

It is convenient to isolate the part of the COM-null ideal that actually enters
the ordinary contact family.  In addition to the earlier linear generators
\begin{equation}
\mathcal C_1\equiv pa+qc,
\qquad
\mathcal C_2\equiv pc+qb,
\label{eq:app-gauge-orbit-C12}
\end{equation}
the four generators that appear in the sparse contact basis are
\begin{equation}
\boxed{
\mathcal C_3\equiv pd+qe,
\qquad
\mathcal C_4\equiv ab-c^2,
\qquad
\mathcal C_5\equiv ae-cd,
\qquad
\mathcal C_6\equiv bd-ce.}
\label{eq:app-gauge-orbit-C3456}
\end{equation}
A sparse basis for the actual scaffold-preserving and COM-preserving ordinary
contact orbit may then be chosen as
\begin{align}
\Gamma_1&=G\,(-a\mathcal C_5+b\mathcal C_6),
&
\Gamma_2&=G\,(b\mathcal C_5-a\mathcal C_6),
\label{eq:app-gauge-orbit-Gamma12}
\\
\Gamma_3&=G\,(d+e)(e\mathcal C_5-d\mathcal C_6),
&
\Gamma_4&=-G\,(d-e)\mathcal C_4,
\label{eq:app-gauge-orbit-Gamma34}
\\
\Gamma_5&=-Gde(\mathcal C_5-\mathcal C_6),
&
\Gamma_6&=G\,(-d^2\mathcal C_5+e^2\mathcal C_6),
\label{eq:app-gauge-orbit-Gamma56}
\\
\Gamma_7&=\frac{G^2}{r}(q\mathcal C_5-p\mathcal C_6),
&
\Gamma_8&=\frac{G^2}{r}(a-b)\mathcal C_3,
\label{eq:app-gauge-orbit-Gamma78}
\\
\Gamma_9&=\frac{G^2}{r}(-p\mathcal C_5+q\mathcal C_6),
&
\Gamma_{10}&=\frac{G^2}{r}(d^2-e^2)\mathcal C_3,
\label{eq:app-gauge-orbit-Gamma910}
\\
\Gamma_{11}&=-\frac{G^3}{r^2}(p-q)\mathcal C_3.
\label{eq:app-gauge-orbit-Gamma11}
\end{align}
The factorized form in
\eqref{eq:app-gauge-orbit-Gamma12}--\eqref{eq:app-gauge-orbit-Gamma11} is more
informative than the fully expanded version printed by the referee audit,
because it makes three structural facts immediate.

First, every surviving generator factors through
\begin{equation}
\Gamma_i\in\langle \mathcal C_3,\mathcal C_4,\mathcal C_5,\mathcal C_6\rangle,
\qquad i=1,\dots,11,
\label{eq:app-gauge-orbit-factorization}
\end{equation}
so every generator already vanishes in the COM quotient.
Second, no generator uses $\mathcal C_1$ or $\mathcal C_2$, so the actual
ordinary contact orbit lives only in the smaller subideal
\begin{equation}
\boxed{
\mathfrak I_{\rm contact}
=
\langle \mathcal C_3,\mathcal C_4,\mathcal C_5,\mathcal C_6\rangle
\subset
\mathfrak I_{\rm COM}.}
\label{eq:app-gauge-orbit-contact-ideal}
\end{equation}
Third, because none of the generators carries a purely static $G^4/r^4$ piece,
the simple ordinary contact family leaves the static block untouched.

\subsection{Image inside the COM-null family}
\label{app:gauge-orbit-image}

After Newtonian order reduction and projection back onto the compact scaffold,
the image of the eleven generators above has exact rank
\begin{equation}
\boxed{\operatorname{rank}\,\mathrm{Im}(\Gamma_1,\dots,\Gamma_{11})=11.}
\label{eq:app-gauge-orbit-image-rank}
\end{equation}
By contrast, the full COM-null family identified in the previous appendix has
rank
\begin{equation}
\boxed{\operatorname{rank}\,\mathfrak I_{\rm COM}=27.}
\label{eq:app-gauge-orbit-com-null-rank}
\end{equation}
So the residual algebraic quotient left after dividing out the actual ordinary
contact orbit has dimension
\begin{equation}
\boxed{27-11=16.}
\label{eq:app-gauge-orbit-residual-quotient}
\end{equation}
Equivalently, if $M_{\rm full}$ denotes the full COM projection matrix on the
compact scaffold, then the exact referee statement is
\begin{equation}
M_{\rm full}\,\mathrm{Im}(\Gamma_1,\dots,\Gamma_{11})=0,
\label{eq:app-gauge-orbit-zero-com-image}
\end{equation}
but the converse fails because the COM-null kernel is larger than the contact
image.

The blockwise contact-image ranks are
\begin{equation}
\boxed{\operatorname{rank}_Q=6,
\qquad
\operatorname{rank}_T=9,
\qquad
\operatorname{rank}_S=4,
\qquad
\operatorname{rank}_U=0.}
\label{eq:app-gauge-orbit-block-ranks}
\end{equation}
The vanishing $U$-block rank is the exact statement that the simple ordinary
contact family does not move the static $G^4/r^4$ slot at all.  More broadly,
\eqref{eq:app-gauge-orbit-block-ranks} explains why the canonical quotient slice
from the previous appendix is larger than the actual ordinary contact orbit: the
contact image only sweeps out one $11$-dimensional suborbit inside the full
$27$-dimensional COM-blind family.

This is the precise conclusion needed before the generic-frame uniqueness theorem
can be stated without overclaiming.  The canonical quotient slice is still the
right sparse algebraic representative, but it is not exhausted by ordinary gauge
alone.  That is why the final uniqueness claim must be deferred to the
Hamiltonian-first fixed-chart argument of the next appendix rather than being
read off from ordinary contact freedom at this stage.

% Appendix content for 4d_3pn.tex

\section{Generic-frame target import and Hamiltonian-first lift}
\label{app:hamiltonian-first}

This appendix records the last conceptual clarification needed before the full
fixed-chart generic-frame uniqueness theorem can be stated.  The ordinary
generic-frame target itself is not the problem.  After the one-body gate, the
natural self/static seed, the reduced COM target, the COM-null ideal, and the
ordinary contact family have all been separated cleanly, the remaining question
is only how the imported generic-frame representative should be compared with
COM data.  The answer is now exact: the apparent inconsistency comes entirely
from doing the reduction in the wrong order.  Naive ordinary-Lagrangian COM
substitution is not the right quotient principle; Hamiltonian-first lift is.

\subsection{Ordinary generic-frame target import}
\label{app:hamiltonian-first-import}

Use again the exchange-symmetric generic-frame invariants
\begin{equation}
 a=v_A^2,
 \qquad
 b=v_B^2,
 \qquad
 c=\mathbf v_A\!\cdot\!\mathbf v_B,
 \qquad
 d=v_{An}=\mathbf v_A\!\cdot\!\mathbf n,
 \qquad
 e=v_{Bn}=\mathbf v_B\!\cdot\!\mathbf n,
\label{eq:app-hamiltonian-first-abcde}
\end{equation}
with
\begin{equation}
 p\equiv m_A,
 \qquad
 q\equiv m_B,
 \qquad
 \mathbf r\equiv \mathbf x_A-\mathbf x_B,
 \qquad
 r\equiv |\mathbf r|,
 \qquad
 \mathbf n\equiv \frac{\mathbf r}{r}.
\label{eq:app-hamiltonian-first-pqrn}
\end{equation}
The full generic-frame ordinary $3$PN target used in the paper is imported from
 the standard literature with the ambiguity fixed by
\begin{equation}
\lambda=-\frac{3}{11}\omega_s-\frac{1987}{3080},
\qquad
\omega_s=0,
\qquad
\Longrightarrow
\qquad
\lambda=-\frac{1987}{3080}.
\label{eq:app-hamiltonian-first-lambda}
\end{equation}

It is convenient to split the imported target into the frozen natural
one-body/self-static seed and the genuinely comparable-mass residual.  Write
\begin{equation}
L_{3,\mathrm{gen}}^{\mathrm{imp}}
=
L_{3,\mathrm{seed}}
+\frac{Gm_A m_B}{r}\,Q_{\rm res}
+\frac{G^2m_A m_B}{r^2}\,T_{\rm res}
+\frac{G^3m_A m_B}{r^3}\,S_{\rm res}
+\frac{G^4m_A m_B}{r^4}\,U_{\rm res},
\label{eq:app-hamiltonian-first-import-split}
\end{equation}
where the seed polynomials are
\begin{equation}
\begin{aligned}
Q_{\rm seed}&=\frac{11}{16}(a^3+b^3),
& T_{\rm seed}&=\frac{47}{16}(qa^2+pb^2),
\\
S_{\rm seed}&=\frac{13}{8}(q^2a+p^2b),
& U_{\rm seed}&=-\frac18(q^3+p^3).
\end{aligned}
\label{eq:app-hamiltonian-first-seed-polynomials}
\end{equation}
The exact import audit shows that after subtraction of
\eqref{eq:app-hamiltonian-first-seed-polynomials} the residual lands entirely
inside the compact generic-frame interaction scaffold,
\begin{equation}
Q_{\rm res}\in\mathrm{span}\{Q_i\}_{i=1}^{24},
\qquad
T_{\rm res}\in\mathrm{span}\{T_j\}_{j=1}^{17},
\qquad
S_{\rm res}\in\mathrm{span}\{S_k\}_{k=1}^{8},
\qquad
U_{\rm res}\in\mathrm{span}\{U_{01}\},
\label{eq:app-hamiltonian-first-scaffold-membership}
\end{equation}
so the imported target may be written exactly as
\begin{equation}
\begin{aligned}
L_{3,\mathrm{gen}}^{\mathrm{imp}}
=\;&L_{3,\mathrm{seed}}
+\frac{Gm_A m_B}{r}\sum_{i=1}^{24} q_i^{\rm imp} Q_i
+\frac{G^2m_A m_B}{r^2}\sum_{j=1}^{17} t_j^{\rm imp} T_j
\\
&+\frac{G^3m_A m_B}{r^3}\sum_{k=1}^{8} s_k^{\rm imp} S_k
+\frac{G^4m_A m_B}{r^4}u_1^{\rm imp}U_{01}.
\end{aligned}
\label{eq:app-hamiltonian-first-import-scaffold}
\end{equation}
No extra generic-frame direction outside the $24/17/8/1$ scaffold is required.
The import also preserves the strict one-body gate, because on the test-mass
branch
\begin{equation}
 p\to 0,
 \qquad
 q\to 1,
 \qquad
 b\to 0,
 \qquad
 c\to 0,
 \qquad
 e\to 0,
\label{eq:app-hamiltonian-first-testmass-branch}
\end{equation}
one has
\begin{equation}
Q_{\rm res}\to 0,
\qquad
T_{\rm res}\to 0,
\qquad
S_{\rm res}\to 0,
\qquad
U_{\rm res}\to 0.
\label{eq:app-hamiltonian-first-testmass-vanish}
\end{equation}
So the imported ordinary generic-frame target is exact, compatible with the
one-body gate, and already expressible in the finite interaction scaffold.

\subsection{Naive COM reduction mismatch}
\label{app:hamiltonian-first-mismatch}

The decisive subtlety appears only after one tries to compare the imported
ordinary generic-frame target directly with the previously solved COM ordinary
ledger.  In the COM frame,
\begin{equation}
\mathbf v_A=X_B\mathbf v,
\qquad
\mathbf v_B=-X_A\mathbf v,
\qquad
X_A+X_B=1,
\qquad
\nu=X_AX_B,
\label{eq:app-hamiltonian-first-com-relations}
\end{equation}
so one might try to obtain the reduced COM ordinary target by simple
substitution into \eqref{eq:app-hamiltonian-first-import-scaffold}.  That
procedure fails.

If one performs this naive ordinary reduction, the resulting COM coefficients do
not coincide with the exact COM ordinary target of the earlier appendices.  The
exact mismatches are
\begin{equation}
\delta l_6=\frac{\nu(80\nu^2+48\nu-17)}{16},
\qquad
\delta l_7=\frac{3\nu(24\nu^2-10\nu+1)}{16},
\qquad
\delta l_8=\frac{3\nu^2(4\nu-1)}{8},
\qquad
\delta l_9=0,
\label{eq:app-hamiltonian-first-dl6-9}
\end{equation}
\begin{equation}
\delta l_{10}=\frac{\nu(-52\nu^2-123\nu+34)}{16},
\qquad
\delta l_{11}=\frac{\nu(-124\nu^2-97\nu+32)}{16},
\qquad
\delta l_{12}=\nu^2(1-4\nu),
\label{eq:app-hamiltonian-first-dl10-12}
\end{equation}
\begin{equation}
\delta l_{13}=\frac{\nu(4\nu^2+27\nu-7)}{8},
\qquad
\delta l_{14}=\frac{\nu(28\nu^2+69\nu-19)}{8},
\qquad
\delta l_{15}=0.
\label{eq:app-hamiltonian-first-dl13-15}
\end{equation}
So naive ordinary COM reduction agrees only in the $l_9$ and $l_{15}$ slots.
This is not a contradiction and it is not a defect of the imported target.
Rather, it is the familiar $3$PN reduction-ordering subtlety: the reduced
relative ordinary Lagrangian does not follow from the general-frame ordinary
representative by naive substitution alone.  Therefore naive ordinary COM
reduction cannot be used as the final generic-to-COM quotient principle.

\subsection{Hamiltonian-first recovery theorem}
\label{app:hamiltonian-first-recovery}

The correct cure is to apply the exact cubic Legendre transform in the generic
frame first, and only then reduce the resulting Hamiltonian to COM.  Carry
forward the exact cubic compiler from Appendix E,
\begin{equation}
H_3=-L_3(v_0)+A_0^TM^{-1}B_0-\frac12A_0^TM^{-1}C_0M^{-1}A_0,
\label{eq:app-hamiltonian-first-cubic-compiler}
\end{equation}
with
\begin{equation}
v_0=M^{-1}P,
\qquad
A_0=\left.\frac{\partial L_1}{\partial v}\right|_{v_0},
\qquad
B_0=\left.\frac{\partial L_2}{\partial v}\right|_{v_0},
\qquad
C_0=\left.\frac{\partial^2L_1}{\partial v^2}\right|_{v_0}.
\label{eq:app-hamiltonian-first-ABCM}
\end{equation}
Because the lower-order $1$PN/$2$PN ledger is already frozen, the feedback terms
in \eqref{eq:app-hamiltonian-first-cubic-compiler} are completely determined
before the new generic-frame $3$PN residual is inserted.  One therefore lifts
\eqref{eq:app-hamiltonian-first-import-scaffold} to the generic-frame
Hamiltonian first and imposes the COM relations only afterward.

That procedure gives the exact theorem
\begin{equation}
\boxed{H_{3,\rm COM}^{\rm(lift)}=H_{3,\rm COM}^{\rm target}.}
\label{eq:app-hamiltonian-first-theorem}
\end{equation}
Equivalently,
\begin{equation}
\boxed{
\text{generic-frame Legendre lift first}
\Longrightarrow
H_{3,\rm COM}^{\rm lift}=H_{3,\rm COM}^{\rm GR}.}
\label{eq:app-hamiltonian-first-theorem-boxed}
\end{equation}
So the earlier mismatch is diagnosed completely:
\begin{equation}
\text{naive COM reduction of ordinary }L_3
\neq
\text{correct COM ordinary target},
\label{eq:app-hamiltonian-first-ordering-diagnosis}
\end{equation}
but the imported generic-frame ordinary target itself is exact and perfectly
compatible with the one-body gate.  The error was the order of operations, not
any defect in the target.

\subsection{What Appendix L fixes for the later lift}
\label{app:hamiltonian-first-fixes}

Appendix L settles the last conceptual issue in the generic-frame import stage.
Its output can be summarized in four statements.

First, the imported generic-frame ordinary target is exact, lives entirely
inside the compact $24/17/8/1$ interaction scaffold, and continues to vanish on
 the strict test-mass branch.  So the generic-frame target itself is fully
consistent with the lower-order backbone.

Second, naive ordinary-Lagrangian COM reduction is not the right quotient
principle.  The exact nonzero mismatches
\eqref{eq:app-hamiltonian-first-dl6-9}--\eqref{eq:app-hamiltonian-first-dl13-15}
show that direct ordinary COM substitution loses information at precisely the
level where the $3$PN chart subtlety matters.

Third, Hamiltonian-first lift cures the problem exactly.  Once the imported
ordinary target is lifted through the cubic compiler before COM reduction, the
standard COM $3$PN Hamiltonian is recovered without residue.  This vindicates
both the imported target and the lower-order compiler chain.

Fourth, after this appendix the remaining ambiguity question is no longer
``is the target right?'' but rather ``which ordinary representative/canonical
slice survives in the fixed Hamiltonian chart?''  That is precisely the point at
which the next appendix takes over.  From there onward the lift is no longer a
COM-quotient problem; it becomes the fixed-chart generic-frame Hamiltonian
compiler and uniqueness theorem.

% Appendix content for 4d_3pn.tex

\section{Generic-frame Hamiltonian compiler and uniqueness}
\label{app:generic-frame-hamiltonian-compiler-uniqueness}

This appendix records the point at which the generic-frame $3$PN lift stops
being a COM-quotient problem and becomes an exact compiler theorem.
The previous appendix already showed that the imported
ordinary generic-frame target is exact and that the correct comparison with COM
data is Hamiltonian-first rather than naive ordinary substitution.  What
remains is therefore sharper: does the previously identified $27$-dimensional
COM-blind family represent a genuine $3$PN ambiguity, or only a blindness of
COM reduction?  In the fixed ADM Hamiltonian chart the answer is now exact: it
is only a blindness.

The earlier appendices had already done four things:
\begin{enumerate}
    \item frozen the exact one-body $3$PN gate and the natural one-body/self-static seed;
    \item solved the reduced COM $15$-slot Hamiltonian target and its ordinary-side image;
    \item built the $50$-slot generic-frame interaction scaffold with block sizes
    \begin{equation}
        24,\qquad 17,\qquad 8,\qquad 1;
        \label{eq:app-gfhc-block-counts-intro}
    \end{equation}
    \item and identified the $27$-dimensional COM-blind family together with the smaller
    $11$-generator ordinary contact orbit.
\end{enumerate}
The present appendix proves that once the lower-order ledger and the natural
$3$PN seed are frozen, the full generic-frame Hamiltonian compiler has no
kernel and therefore removes the whole COM-null family at once.

\subsection{Fixed generic-frame momentum basis}
\label{app:generic-frame-hamiltonian-compiler-uniqueness-basis}

Use the same formal invariant symbols as in the ordinary generic-frame
scaffold, but reinterpret them on the Hamiltonian side as Newtonian-order
momentum invariants:
\begin{equation}
 a=\frac{P_A^2}{m_A^2},
 \qquad
 b=\frac{P_B^2}{m_B^2},
 \qquad
 c=\frac{P_A\!\cdot\!P_B}{m_A m_B},
 \qquad
 d=\frac{\mathbf n\!\cdot\!P_A}{m_A},
 \qquad
 e=\frac{\mathbf n\!\cdot\!P_B}{m_B}.
\label{eq:app-gfhc-abcde}
\end{equation}
Write again
\begin{equation}
 p\equiv m_A,
 \qquad
 q\equiv m_B,
 \qquad
 \mathbf r\equiv \mathbf x_A-\mathbf x_B,
 \qquad
 r\equiv |\mathbf r|,
 \qquad
 \mathbf n\equiv \frac{\mathbf r}{r}.
\label{eq:app-gfhc-pqrn}
\end{equation}
In these variables the generic-frame Hamiltonian interaction scaffold has the
same formal block structure as the ordinary one:
\begin{itemize}
    \item a $24$-slot sextic block at order $G/r$,
    \item a $17$-slot quartic block at order $G^2/r^2$,
    \item an $8$-slot quadratic block at order $G^3/r^3$,
    \item and a $1$-slot static block at order $G^4/r^4$.
\end{itemize}
So the full generic-frame $3$PN interaction ledger is still a $50$-slot basis.
Nothing new is being hidden in a different chart; only the kinematic
interpretation of the invariants changes.

Write the ordinary interaction residual beyond the already frozen one-body/self-
static seed as
\begin{equation}
\Delta L_3
=
\frac{Gpq}{r}\sum_{i=1}^{24} q_i Q_i
+\frac{G^2pq}{r^2}\sum_{j=1}^{17} t_j T_j
+\frac{G^3pq}{r^3}\sum_{k=1}^{8} s_k S_k
+\frac{G^4pq}{r^4}u_1 U_{01},
\label{eq:app-gfhc-deltaL-expansion}
\end{equation}
and write the Hamiltonian residual in the same formal invariant chart as
\begin{equation}
\Delta H_3
=
\frac{Gpq}{r}\sum_{i=1}^{24} \hat q_i Q_i
+\frac{G^2pq}{r^2}\sum_{j=1}^{17} \hat t_j T_j
+\frac{G^3pq}{r^3}\sum_{k=1}^{8} \hat s_k S_k
+\frac{G^4pq}{r^4}\hat u_1 U_{01}.
\label{eq:app-gfhc-deltaH-expansion}
\end{equation}
The problem of uniqueness is therefore exactly the problem of how the $50$
ordinary coefficients in \eqref{eq:app-gfhc-deltaL-expansion} are mapped into
the $50$ Hamiltonian coefficients in \eqref{eq:app-gfhc-deltaH-expansion} once
all lower-order data are frozen.

\subsection{Exact compiler theorem}
\label{app:generic-frame-hamiltonian-compiler-uniqueness-theorem}

Carry forward the exact cubic Legendre identity proved earlier:
\begin{equation}
H_3=-L_3(v_0)+A_0^TM^{-1}B_0-\frac12A_0^TM^{-1}C_0M^{-1}A_0,
\label{eq:app-gfhc-cubic-compiler}
\end{equation}
with
\begin{equation}
 v_0=M^{-1}P,
 \qquad
 A_0=\left.\frac{\partial L_1}{\partial v}\right|_{v_0},
 \qquad
 B_0=\left.\frac{\partial L_2}{\partial v}\right|_{v_0},
 \qquad
 C_0=\left.\frac{\partial^2L_1}{\partial v^2}\right|_{v_0}.
\label{eq:app-gfhc-cubic-data}
\end{equation}
Now split the full ordinary $3$PN block as
\begin{equation}
L_3=L_{3,\mathrm{seed}}+\Delta L_3,
\label{eq:app-gfhc-seed-plus-residual}
\end{equation}
where $L_{3,\mathrm{seed}}$ is the natural one-body/self-static seed already
frozen in the previous appendices.  The decisive point is that the second and
third terms in \eqref{eq:app-gfhc-cubic-compiler} depend only on the already
frozen lower-order ledger $(L_1,L_2)$.  They do not depend on the new $3$PN
interaction residual $\Delta L_3$.  Therefore the full generic-frame residual
compiler is exact and trivial:
\begin{equation}
\boxed{\Delta H_3=-\Delta L_3(v_0).}
\label{eq:app-gfhc-residual-compiler}
\end{equation}
Comparing \eqref{eq:app-gfhc-deltaL-expansion} and
\eqref{eq:app-gfhc-deltaH-expansion}, this is literally the slotwise sign flip
\begin{equation}
\boxed{
\hat q_i=-q_i,
\qquad
\hat t_j=-t_j,
\qquad
\hat s_k=-s_k,
\qquad
\hat u_1=-u_1.}
\label{eq:app-gfhc-slotwise-signflip}
\end{equation}
Equivalently, if the $50$ ordinary residual coordinates are assembled into one
column vector $\Delta l$ and the $50$ Hamiltonian residual coordinates into
$\Delta h$, then
\begin{equation}
\boxed{
\Delta h=C_{\rm gen}\,\Delta l,
\qquad
C_{\rm gen}=-I_{50}.}
\label{eq:app-gfhc-compiler-matrix}
\end{equation}
It is useful to display the blockwise form explicitly:
\begin{equation}
C_Q=-I_{24},
\qquad
C_T=-I_{17},
\qquad
C_S=-I_{8},
\qquad
C_U=-I_1,
\qquad
C_{\rm gen}=\mathrm{diag}(C_Q,C_T,C_S,C_U).
\label{eq:app-gfhc-blockwise-compiler}
\end{equation}
The referee audit therefore gives the exact linear-algebraic status
\begin{equation}
\operatorname{rank}C_{\rm gen}=50,
\qquad
\ker C_{\rm gen}=\{0\}.
\label{eq:app-gfhc-rank-nullity}
\end{equation}
So the fixed-chart generic-frame Hamiltonian compiler has full rank and no
kernel.  Once the lower-order blocks and the natural $3$PN seed are frozen, the
generic-frame lift is no longer a nonlinear inversion problem.  The Hamiltonian
residual is simply the negative of the ordinary residual in the same invariant
chart.

\subsection{Direct recovery of the ordinary representative}
\label{app:generic-frame-hamiltonian-compiler-uniqueness-recovery}

Let the exact imported ordinary target from the previous appendix be
written on the compact interaction scaffold as the coefficient vectors
\begin{equation}
\mathbf q^{\rm imp}\in\mathbb R^{24},
\qquad
\mathbf t^{\rm imp}\in\mathbb R^{17},
\qquad
\mathbf s^{\rm imp}\in\mathbb R^{8},
\qquad
u_1^{\rm imp}\in\mathbb R,
\label{eq:app-gfhc-imported-vectors}
\end{equation}
so that
\begin{equation}
\Delta L_3^{\rm imp}
=
\frac{Gpq}{r}\sum_{i=1}^{24} q_i^{\rm imp}Q_i
+\frac{G^2pq}{r^2}\sum_{j=1}^{17} t_j^{\rm imp}T_j
+\frac{G^3pq}{r^3}\sum_{k=1}^{8} s_k^{\rm imp}S_k
+\frac{G^4pq}{r^4}u_1^{\rm imp}U_{01}.
\label{eq:app-gfhc-imported-ordinary-target}
\end{equation}
Then the fixed-chart Hamiltonian target is obtained immediately by the sign flip
of \eqref{eq:app-gfhc-slotwise-signflip}:
\begin{equation}
\mathbf q_H^{\rm target}=-\mathbf q^{\rm imp},
\qquad
\mathbf t_H^{\rm target}=-\mathbf t^{\rm imp},
\qquad
\mathbf s_H^{\rm target}=-\mathbf s^{\rm imp},
\qquad
u_{1,H}^{\rm target}=-u_1^{\rm imp}.
\label{eq:app-gfhc-hamiltonian-target-signflip}
\end{equation}
Because the inverse compiler is the same matrix,
\begin{equation}
C_{\rm gen}^{-1}=C_{\rm gen}=-I_{50},
\label{eq:app-gfhc-inverse-compiler}
\end{equation}
one recovers the ordinary representative directly from the fixed-chart
Hamiltonian target:
\begin{equation}
\mathbf q^{\rm imp}=-\mathbf q_H^{\rm target},
\qquad
\mathbf t^{\rm imp}=-\mathbf t_H^{\rm target},
\qquad
\mathbf s^{\rm imp}=-\mathbf s_H^{\rm target},
\qquad
u_1^{\rm imp}=-u_{1,H}^{\rm target}.
\label{eq:app-gfhc-ordinary-recovery}
\end{equation}
So once the fixed ADM Hamiltonian chart is chosen, there is no further algebraic
quotienting step.  The exact ordinary representative is read off immediately by
sign reversal of the Hamiltonian residual coordinates.

\subsection{Fixed-chart uniqueness and the fate of the COM-null family}
\label{app:generic-frame-hamiltonian-compiler-uniqueness-fixed-chart}

This theorem resolves the earlier $27$-dimensional COM-blind family exactly.
Before the full Hamiltonian chart was imposed, the ordinary generic-frame lift
had
\begin{itemize}
    \item a $27$-dimensional family invisible to COM reduction,
    \item inside it an $11$-generator ordinary contact/gauge orbit,
    \item and therefore a residual $16$-dimensional quotient after the true
    contact family was removed.
\end{itemize}
That was a real COM-side obstruction, but it was only a COM-side obstruction.
Because the full generic-frame Hamiltonian compiler has zero kernel, none of
those $27$ directions is Hamiltonian-blind in the fixed ADM chart.  So the full
generic-frame Hamiltonian target removes the whole COM-null family at once.
Equivalently,
\begin{equation}
\boxed{\text{COM blindness is not Hamiltonian blindness.}}
\label{eq:app-gfhc-com-blind-not-hamiltonian-blind}
\end{equation}
The $11$-generator ordinary contact family is not an exception to this.  Inside
the fixed ADM Hamiltonian chart, those directions do not lie in the compiler
kernel; they correspond instead to moving to a different canonical chart.
Once that chart is frozen, the Hamiltonian target selects a single ordinary
representative.

The uniqueness theorem may therefore be written in its sharpest form as
follows.  Suppose $\Delta L_3$ and $\Delta L_3'$ are two generic-frame ordinary
residuals that lead to the same fixed-chart Hamiltonian target.  Then by
\eqref{eq:app-gfhc-residual-compiler},
\begin{equation}
0=\Delta H_3-\Delta H_3'=-\bigl(\Delta L_3-\Delta L_3'\bigr),
\label{eq:app-gfhc-uniqueness-proof}
\end{equation}
so
\begin{equation}
\boxed{\Delta L_3'=\Delta L_3.}
\label{eq:app-gfhc-ordinary-uniqueness}
\end{equation}
In other words, once the fixed ADM Hamiltonian chart is chosen, the generic-
frame ordinary representative is unique.

At that point the $3$PN generic-frame lift is no longer the bottleneck.  The
live conservative work moves entirely into the physical interpretation of the
solved residual: the grouped real $P_2$ middle block, the unique static
geometry completion, and the pure-kinetic compiler image.

\subsection{What this appendix fixes for the later split}
\label{app:generic-frame-hamiltonian-compiler-uniqueness-summary}

This appendix settles the last generic-frame algebra question in four
steps.

First, it fixes the Hamiltonian-side invariant chart in exactly the same
$24/17/8/1$ interaction scaffold used on the ordinary side, so no hidden degree
count is being changed when one passes to momentum variables.

Second, it proves that the exact cubic Legendre transform acts on the full
$50$-slot interaction residual by pure sign flip,
\(C_{\rm gen}=-I_{50}\), once the lower-order ledger and the natural $3$PN seed
are frozen.

Third, it shows that the compiler has zero kernel.  The previously found
$27$-dimensional COM-blind family, and the smaller $11$-generator ordinary
contact orbit sitting inside it, therefore do not survive as true algebraic
freedoms in the fixed ADM chart.

Fourth, it gives the fixed-chart uniqueness theorem
\eqref{eq:app-gfhc-ordinary-uniqueness}.  From this point onward, the live
conservative questions are no longer about which generic-frame representative is
allowed.  They are only about how the unique residual decomposes physically into
its grouped real $P_2$, geometry, and pure-kinetic lanes.

% Appendix content for 4d_3pn.tex

\section{Grouped real \texorpdfstring{$P_2$}{P 2} target pack and co-rotating front-end scaffold}
\label{app:grouped-p2-target-pack}

This appendix freezes the exact conservative grouped-$P_2$ inverse problem that
stands between the solved comparable-mass $3$PN target and any future
moving-throat constitutive interpretation.  Its job is twofold.  First, it
states the exact $3$PN data vector that any grouped real
\(
P_{20}\oplus P_{21}\oplus P_{22}
\)
conservative module must face once the frozen one-body/self-static seed is
subtracted.  Second, it writes the first exact co-rotating time-local
\(O(\omega^2)\) front-end scaffold in grouped language, before any throat-side
closure is guessed.

The point of the appendix is therefore not to complete the $3$PN grouped
module---that comes only after the richer local lift is admitted---but to make
clear exactly what information the grouped branch carries, how sharply it is
constrained, and what the smallest kinematic front end looks like before any
constitutive dictionary is imposed.

\subsection{Exact \texorpdfstring{$3$}{3}PN data vector seen by the grouped front end}
\label{app:grouped-p2-target-pack-data}

From Appendix~\ref{app:gr-com-target}, the exact comparable-mass residual beyond
the frozen one-body/self-static seed is
\begin{equation}
\begin{aligned}
 \Delta l_1&=\frac{3\nu(3\nu-1)(4\nu-1)}{16},
 & \Delta l_2&=\Delta l_3=\Delta l_4=\Delta l_5=0.
\end{aligned}
\label{eq:app-grouped-p2-target-pack-dl-1-5}
\end{equation}
\begin{equation}
\begin{aligned}
 \Delta l_6&=\frac{\nu(38-116\nu-57\nu^2)}{16},
 & \Delta l_7&=\frac{\nu^2(20-69\nu)}{16},
\\[0.5ex]
 \Delta l_8&=\frac{3\nu^2(3-11\nu)}{16},
 & \Delta l_9&=\frac{5\nu^3}{16}.
\end{aligned}
\label{eq:app-grouped-p2-target-pack-dl-6-9}
\end{equation}
\begin{equation}
\begin{aligned}
 \Delta l_{10}&=\frac{\nu(129-98\nu+52\nu^2)}{16},
 & \Delta l_{11}&=\frac{\nu(-3+52\nu+124\nu^2)}{16},
\\[0.5ex]
 \Delta l_{12}&=\frac{\nu(-5+11\nu+48\nu^2)}{12}.
\end{aligned}
\label{eq:app-grouped-p2-target-pack-dl-10-12}
\end{equation}
\begin{equation}
 \Delta l_{13}=-\frac{\nu(244+3\pi^2+1272\nu+96\nu^2)}{192}.
\label{eq:app-grouped-p2-target-pack-dl-13}
\end{equation}
\begin{equation}
\begin{aligned}
 \Delta l_{14}&=\frac{\nu(452+3\pi^2-384\nu-224\nu^2)}{64},
 & \Delta l_{15}&=\frac{\nu(-908+63\pi^2)}{96}.
\end{aligned}
\label{eq:app-grouped-p2-target-pack-dl-14-15}
\end{equation}
The grouped conservative front end can only hope to carry the middle block
\((\Delta l_6,\dots,\Delta l_{14})\).  The pure-kinetic slot \(\Delta l_1\) and the
static scalar/geometry slot \(\Delta l_{15}\) therefore sit outside the first
minimal grouped inverse problem from the beginning.

At $3$PN the grouped conservative payload through order \(O(\omega^2)\) contains
only three real data,
\begin{equation}
 u_2^{(20)},\qquad u_2^{(21)},\qquad u_2^{(22)}.
\label{eq:app-grouped-p2-target-pack-u2-data}
\end{equation}
So the grouped front end faces an exact $11$-slot residual vector while carrying
only three real conservative unknowns.  Before any isotropy assumption is made,
this is already an eight-fold overdetermined test.  On the minimal isotropic
branch it becomes a ten-fold overdetermined test.  That sharp mismatch of data
counts is the real reason the grouped-$P_2$ route is such a useful pre-PDE
conservative gate: once a constitutive dictionary is written, the test is
brutally exact rather than qualitative.

\subsection{Grouped trace, anisotropy variables, and the isotropy gate}
\label{app:grouped-p2-target-pack-trace}

It is often more transparent to rewrite the three grouped coefficients in terms
of one grouped trace and two anisotropy defects.  Define
\begin{equation}
 \bar u_2=\frac{u_2^{(20)}+2u_2^{(21)}+2u_2^{(22)}}{5},
 \qquad
 a_2=\frac{2u_2^{(20)}-u_2^{(21)}-u_2^{(22)}}{10},
 \qquad
 b_2=\frac{u_2^{(21)}-u_2^{(22)}}{2}.
\label{eq:app-grouped-p2-target-pack-trace-defects}
\end{equation}
The inverse map is exact:
\begin{equation}
 u_2^{(20)}=\bar u_2+4a_2,
 \qquad
 u_2^{(21)}=\bar u_2-a_2+b_2,
 \qquad
 u_2^{(22)}=\bar u_2-a_2-b_2.
\label{eq:app-grouped-p2-target-pack-inverse}
\end{equation}
Equivalently, the grouped weighted-sum identity is automatic,
\begin{equation}
 \bigl(u_2^{(20)}-\bar u_2\bigr)
 +2\bigl(u_2^{(21)}-\bar u_2\bigr)
 +2\bigl(u_2^{(22)}-\bar u_2\bigr)=0.
\label{eq:app-grouped-p2-target-pack-weighted-sum}
\end{equation}
The natural grouped anisotropy norm is therefore
\begin{equation}
 \mathcal A_2^2=4a_2^2+\frac{4}{5}b_2^2.
\label{eq:app-grouped-p2-target-pack-anisotropy}
\end{equation}
So the exact isotropy gate for the minimal branch is simply
\begin{equation}
 \boxed{a_2=0,
 \qquad
 b_2=0,}
\label{eq:app-grouped-p2-target-pack-isotropy}
\end{equation}
or equivalently \(\mathcal A_2^2=0\).  If that gate passes, then all three
$O(\omega^2)$ grouped coefficients coincide and one may write
\begin{equation}
 u_2^{(20)}=u_2^{(21)}=u_2^{(22)}\equiv u_2.
\label{eq:app-grouped-p2-target-pack-isotropic-u2}
\end{equation}
This is the same isotropy test that later feeds directly into the narrowed
$2.5$PN quadrupole bridge.

\subsection{Co-rotating real-\texorpdfstring{$P_2$}{P 2} source coefficients}
\label{app:grouped-p2-target-pack-corotating}

Work in the instantaneous co-rotating pair frame and decompose the relative
velocity as
\begin{equation}
 \mathbf v=\mathbf u+d\,\mathbf n,
 \qquad
 \mathbf u\!\cdot\!\mathbf n=0,
 \qquad
 v^2=u^2+d^2,
 \qquad
 U=\frac{GM}{r}.
\label{eq:app-grouped-p2-target-pack-pair-frame}
\end{equation}
In this frame the grouped real source coefficients of the normalized STF orbital
quadrupole
\(\ddot I_{ij}^{\rm orb}/(2\mu)\)
are
\begin{align}
 C_{20}&=\sqrt{\frac23}\left(d^2-\frac{u^2}{2}-U\right),
\label{eq:app-grouped-p2-target-pack-C20}
\\[0.5ex]
 C_{21c}&=\sqrt2\,d\,u_x,
 & C_{21s}&=\sqrt2\,d\,u_y,
\label{eq:app-grouped-p2-target-pack-C21}
\\[0.5ex]
 C_{22c}&=\frac{u_x^2-u_y^2}{\sqrt2},
 & C_{22s}&=\sqrt2\,u_xu_y.
\label{eq:app-grouped-p2-target-pack-C22}
\end{align}
Using the exact co-rotating pair-frame kinematics
\begin{equation}
 \dot d=\frac{u^2-U}{r},
 \qquad
 \dot u_x=-\frac{d}{r}u_x,
 \qquad
 \dot u_y=-\frac{d}{r}u_y,
 \qquad
 \dot U=-\frac{d}{r}U,
\label{eq:app-grouped-p2-target-pack-corotating-kinematics}
\end{equation}
one obtains the exact grouped derivative norms
\begin{equation}
 \dot C_{20}^2=\frac{2d^2}{3r^2}(3u^2-U)^2,
\label{eq:app-grouped-p2-target-pack-dotC20}
\end{equation}
\begin{equation}
 \dot C_{21c}^2+\dot C_{21s}^2
 =\frac{2u^2}{r^2}(u^2-d^2-U)^2,
\label{eq:app-grouped-p2-target-pack-dotC21}
\end{equation}
\begin{equation}
 \dot C_{22c}^2+\dot C_{22s}^2
 =\frac{2d^2u^4}{r^2}.
\label{eq:app-grouped-p2-target-pack-dotC22}
\end{equation}
These formulas are exact kinematic identities of the grouped source map itself.
No throat-side constitutive closure has yet been used.

\subsection{First exact time-local grouped-\texorpdfstring{$P_2$}{P 2} front-end scaffold}
\label{app:grouped-p2-target-pack-front-end}

The smallest time-local grouped conservative ansatz built from the exact source
norms is
\begin{equation}
L_{P_2,\omega^2}^{\rm front}
=\frac12c_2^{(20)}\dot C_{20}^2
+\frac12c_2^{(21)}\bigl(\dot C_{21c}^2+\dot C_{21s}^2\bigr)
+\frac12c_2^{(22)}\bigl(\dot C_{22c}^2+\dot C_{22s}^2\bigr),
\label{eq:app-grouped-p2-target-pack-front-raw}
\end{equation}
with three real grouped coefficients
\(c_2^{(20)},c_2^{(21)},c_2^{(22)}\).
Substituting
\eqref{eq:app-grouped-p2-target-pack-dotC20}--\eqref{eq:app-grouped-p2-target-pack-dotC22}
gives the exact co-rotating front-end scaffold
\begin{equation}
\boxed{
L_{P_2,\omega^2}^{\rm front}
=
\frac{c_2^{(20)}}{3r^2}d^2(3u^2-U)^2
+
\frac{c_2^{(21)}}{r^2}u^2(u^2-d^2-U)^2
+
\frac{c_2^{(22)}}{r^2}d^2u^4.
}
\label{eq:app-grouped-p2-target-pack-front-substituted}
\end{equation}
This is still only a kinematic front end.  But it turns the conservative
grouped-$P_2$ problem into a concrete inverse problem: exact middle-block target
data on one side, and three grouped \(O(\omega^2)\) coefficients on the other.

On the isotropic branch,
\begin{equation}
 a_2=b_2=0,
 \qquad
 c_2^{(20)}=c_2^{(21)}=c_2^{(22)}\equiv \bar c_2,
\label{eq:app-grouped-p2-target-pack-isotropic-c2}
\end{equation}
and the front-end scaffold collapses exactly to
\begin{equation}
\boxed{
L_{P_2,\omega^2}^{\rm iso}
=
\frac{\bar c_2}{3r^2}
\Bigl(
3v^6-3d^2v^4-6Uv^4+12Ud^2v^2-6Ud^4+3U^2v^2-2U^2d^2
\Bigr).
}
\label{eq:app-grouped-p2-target-pack-isotropic-front}
\end{equation}
This isotropic collapse is useful later because it makes the degree of
overdetermination completely explicit: once isotropy is imposed, one conservative
trace datum is being asked to account for the full $3$PN middle block.

That is why Appendix~\ref{app:grouped-p2-target-pack} is the right place to stop
before any throat-side constitutive guess is made.  The exact $3$PN grouped data,
its isotropy gate, the co-rotating source coefficients, and the first exact
front-end scaffold are now fully explicit.  The next question---taken up in the
following appendix---is whether the smallest local constitutive lift of this
front end is already enough.  It is not, and that precise negative result is what
forces the richer grouped local family used later in the paper.

\section{Failure of the minimal local grouped-\texorpdfstring{$P_2$}{P 2} demotion}
\label{app:minimal_local_grouped_p2_demotion_failure}

This appendix records the first exact negative result of the grouped real $P_2$ program at $3$PN. The point is not that the grouped-$P_2$ ontology fails. The point is narrower and more useful: the \emph{smallest} local constitutive lift of the grouped front end is too rigid to reproduce the exact GR middle block. This no-go is what forces the richer grouped compiler used in the main text.

\subsection{Exact grouped target pack and the overdetermined \texorpdfstring{$3$}{3}PN gate}

At $3$PN, the grouped conservative front end is controlled by exactly three real $O(\omega^2)$ data,
\begin{equation}
  u_2^{(20)},\qquad u_2^{(21)},\qquad u_2^{(22)},
\end{equation}
or equivalently by the grouped trace and anisotropy defects
\begin{equation}
  \bar u_2 = \frac{u_2^{(20)} + 2u_2^{(21)} + 2u_2^{(22)}}{5},
  \qquad
  a_2 = \frac{2u_2^{(20)} - u_2^{(21)} - u_2^{(22)}}{10},
  \qquad
  b_2 = \frac{u_2^{(21)} - u_2^{(22)}}{2}.
\end{equation}
So the exact isotropy gate for the minimal branch is
\begin{equation}
  a_2 = 0,
  \qquad
  b_2 = 0.
\end{equation}
The structural reason this is already a hard theorem gate is that the exact $3$PN comparable-mass COM target presents an $11$-slot residual vector in the interaction sector, whereas the grouped front end offers only three conservative unknowns. On the isotropic branch that shrinks to one parameter. In that sense, $3$PN is an overdetermined algebraic test of the grouped real $P_2$ idea rather than a loose plausibility exercise.

The exact co-rotating time-local front end built from the real STF orbital source coefficients takes the form
\begin{equation}
  L_{P_2,\omega^2}^{\rm front}
  =
  \frac{c_2^{(20)}}{3r^2}\,\dot r^2\,(3u_\perp^2-U)^2
  +
  \frac{c_2^{(21)}}{r^2}\,u_\perp^2\,(u_\perp^2-\dot r^2-U)^2
  +
  \frac{c_2^{(22)}}{r^2}\,\dot r^2\,u_\perp^4,
\end{equation}
where $u_\perp^2=v^2-\dot r^2$. This is the first exact grouped kinematic scaffold. The question addressed here is whether its \emph{minimal local demotion} can already match the exact GR middle block.

\subsection{Why the direct grouped front end is too rigid}

The direct grouped front end has uniform virial weight five, so by itself it is too high to be a formal $3$PN ordinary-Lagrangian block. The unique local isotropic one-step demotion that lowers the weight by exactly one unit is multiplication by one power of $U^{-1}\sim r$. After that demotion, the image lands only in the nine interaction slots
\begin{equation}
  (l_6,\dots,l_{14}),
\end{equation}
with no support at all for the pure-kinetic slots $(l_1,\dots,l_5)$ and no support for the static slot $l_{15}$.

This is already informative: even before comparison with the exact GR target, the minimally demoted grouped lift can at most hope to carry the middle block. It cannot possibly account for the pure-kinetic or static sectors. The real obstruction, however, is stronger than that. Because the demoted image depends on only three grouped coefficients, its nine-slot image has rank three and therefore obeys six exact linear relations.

\subsection{The isotropic image and the six exact slot relations}

On the isotropic branch $a_2=b_2=0$, the demoted grouped image is especially rigid. In the standard isotropic COM basis it gives
\begin{equation}
  l_6 = \bar u_2,
  \qquad
  l_7 = -\bar u_2,
  \qquad
  l_8 = l_9 = 0,
\end{equation}
\begin{equation}
  l_{10} = -2\bar u_2,
  \qquad
  l_{11} = 4\bar u_2,
  \qquad
  l_{12} = -2\bar u_2,
  \qquad
  l_{13} = \bar u_2,
  \qquad
  l_{14} = -\frac{2}{3}\bar u_2.
\end{equation}
So even on the most favorable branch, the minimally demoted image is far too constrained to be a generic middle-block carrier.

More generally, because the nine slots depend on only three grouped coefficients, the minimal image obeys six exact relations. A convenient basis is
\begin{align}
  2l_6 + 2l_7 + l_8 &= 0,
  \\
  l_9 &= l_6 + l_7,
  \\
  l_{10} &= -2l_6,
  \\
  l_{11} + l_{12} &= 2l_6,
  \\
  l_{13} &= l_6,
  \\
  l_{14} &= -\frac{1}{6}l_{11}.
\end{align}
Any $3$PN target produced by the minimally demoted local grouped-$P_2$ scaffold must satisfy all six simultaneously.

\subsection{The exact GR target violates all six relations}

Substituting the exact GR residual slots $(\Delta l_6,\dots,\Delta l_{14})$ into these relations gives the obstruction polynomials
\begin{equation}
  2\Delta l_6 + 2\Delta l_7 + \Delta l_8
  =
  \frac{\nu(-285\nu^2-183\nu+76)}{16},
\end{equation}
\begin{equation}
  -\Delta l_6 - \Delta l_7 + \Delta l_9
  =
  \frac{\nu(131\nu^2+96\nu-38)}{16},
\end{equation}
\begin{equation}
  \Delta l_{10} + 2\Delta l_6
  =
  \frac{\nu(-62\nu^2-330\nu+205)}{16},
\end{equation}
\begin{equation}
  \Delta l_{11} + \Delta l_{12} - 2\Delta l_6
  =
  \frac{\nu(906\nu^2+896\nu-257)}{48},
\end{equation}
\begin{equation}
  \Delta l_{13} - \Delta l_6
  =
  \frac{\nu(588\nu^2+120\nu-700-3\pi^2)}{192},
\end{equation}
\begin{equation}
  \Delta l_{14} + \frac{1}{6}\Delta l_{11}
  =
  \frac{\nu(-424\nu^2-1048\nu+9\pi^2+1350)}{192}.
\end{equation}
All six are nonzero for generic $\nu$. In particular, the obstruction does not disappear at equal mass. The point is not one exceptional coefficient mismatch; it is a structural rank failure of the minimally demoted grouped map.

So the exact conclusion is
\begin{equation}
  \boxed{\begin{gathered}
  \text{the minimally demoted local grouped-$P_2$ image}\\
  \text{does not contain the exact GR $3$PN middle block}
  \end{gathered}}
\end{equation}
This is the first exact grouped-$P_2$ no-go of the $3$PN program.

\subsection{Interpretation and consequence for the next step}

This no-go does \emph{not} kill the grouped ontology. It kills only the first, overly rigid constitutive guess. The grouped real $P_2$ bundle is still the right front end; what fails is the minimal local demotion of the direct channel norms. The practical lesson is therefore very sharp:
\begin{enumerate}
  \item the direct grouped $O(\omega^2)$ scaffold is the correct kinematic starting point,
  \item the unique one-step local demotion by one inverse orbital weight is too small,
  \item the true $3$PN grouped constitutive lift must introduce additional middle-block mixing,
  \item and separate mechanisms are still needed for the pure-kinetic and static sectors.
\end{enumerate}
This is exactly why the next appendix enlarges the grouped local family from the three demoted $T_A$ directions to the full nine-basis $(T_A,S_A,V_A)$ compiler. The failure recorded here is therefore a useful theorem gate: it prevents the grouped program from being pursued in an underfit form and directly motivates the richer exact middle-block closure used in the main theorem chain.

\section{Richer grouped-\texorpdfstring{$P_2$}{P 2} compiler and inverse map}
\label{app:richer_grouped_p2_compiler_inverse_map}

This appendix records the exact positive grouped-$P_2$ result that replaces the no-go of the minimal local lift. The point is not that the grouped ontology was too small. The point is that the first local constitutive choice was too rigid. Once the natural local grouped family is enlarged in the unique minimal way described below, the entire nine-slot $3$PN middle block becomes exactly reachable, with an explicit inverse compiler from the solved GR target to the grouped coefficients.

Throughout this appendix we use the reduced COM notation
\[
d \equiv \dot r,
\qquad
u_\perp^2 \equiv v^2-d^2,
\qquad
U \equiv \frac{GM}{r},
\]
and the standard middle-block slot labels $(l_6,\dots,l_{14})$.

\subsection{Why the minimal grouped lift had to be enlarged}
\label{app:richer_grouped_p2_compiler_inverse_map:motivation}

The minimal local grouped-$P_2$ lift built only from the demoted dynamic family $T_A$ lands in the correct nine interaction slots, but its image has rank three and obeys six exact slot relations. The solved GR middle block violates all six, so the minimal lift is too small. The correct interpretation is not that the grouped branch fails, but that the constitutive family is underfit.

The first demoted grouped family is
\[
T_A \equiv U^{-1}\dot C_A^{\,2},
\qquad A\in\{20,21,22\},
\]
where the grouped source-square norms are
\[
C_{20}^2=\frac16\bigl(3d^2-v^2-2U\bigr)^2,
\qquad
C_{21}^2=2d^2\nu_\perp^2,
\qquad
C_{22}^2=\frac12\nu_\perp^4.
\]
In explicit local form this gives
\[
T_{20}=\frac13 U d^2\bigl(3\nu_\perp^2-U\bigr)^2,
\qquad
T_{21}=U\nu_\perp^2\bigl(\nu_\perp^2-d^2-U\bigr)^2,
\qquad
T_{22}=U d^2\nu_\perp^4.
\]

The exact positive result is that the grouped ontology itself is already large enough; only the constitutive family must be enlarged by the two natural local grouped companions
\[
S_A\equiv U^2 C_A^2,
\qquad
V_A\equiv \frac{v^2}{U}S_A=Uv^2C_A^2,
\qquad A\in\{20,21,22\}.
\]
So the richer grouped real $P_2$ family is
\[
(T_{20},T_{21},T_{22},S_{20},S_{21},S_{22},V_{20},V_{21},V_{22}).
\]

\subsection{Exact middle-block compiler}
\label{app:richer_grouped_p2_compiler_inverse_map:compiler}

Write the richer grouped lift as
\[
L_{P_2}^{\rm mid}
=
\sum_{A\in\{20,21,22\}}
\bigl(\lambda_A T_A+\sigma_A S_A+\tau_A V_A\bigr),
\]
with grouped coefficient triples
\[
(\lambda_{20},\lambda_{21},\lambda_{22}),
\qquad
(\sigma_{20},\sigma_{21},\sigma_{22}),
\qquad
(\tau_{20},\tau_{21},\tau_{22}).
\]
Expanding the nine basis elements in the COM monomial basis for the middle block,
\[
(Uv^6,
\;Uv^4d^2,
\;Uv^2d^4,
\;Ud^6,
\;U^2v^4,
\;U^2v^2d^2,
\;U^2d^4,
\;U^3v^2,
\;U^3d^2),
\]
produces an exact $9\times 9$ linear map
\[
M_{\rm mid}:
(\lambda_{20},\lambda_{21},\lambda_{22},
\sigma_{20},\sigma_{21},\sigma_{22},
\tau_{20},\tau_{21},\tau_{22})
\mapsto
(l_6,\dots,l_{14}).
\]
The decisive fact is that this compiler is exactly invertible:
\[
\boxed{\det M_{\rm mid}=-\frac{4}{27}.}
\]
So the richer local grouped family closes the full nine-slot middle block exactly.

\subsection{Exact inverse map}
\label{app:richer_grouped_p2_compiler_inverse_map:inverse}

Because $M_{\rm mid}$ is invertible, the exact inverse compiler from the middle-block target
\[
(L_6,\dots,L_{14})
\]
to the grouped coefficients is algebraic and unique. In the notation above one finds
\[
\boxed{
\lambda_{20}
=
6(L_{10}+L_{11}+L_{12})
+3(L_{13}+L_{14})
+9(L_6+L_7+L_8+L_9),
}
\]
\[
\boxed{
\lambda_{21}
=
-(L_{10}+L_{11}+L_{12})
+L_{13}
-3(L_6+L_7+L_8+L_9),
}
\]
\[
\boxed{
\lambda_{22}
=
-22(L_{10}+L_{11}+L_{12})
-5L_{13}-9L_{14}
-39(L_6+L_7+L_8)-38L_9,
}
\]
\[
\boxed{
\sigma_{20}
=
\frac32(L_{10}+L_{11}+L_{12})
+3(L_6+L_7+L_8+L_9),
}
\]
\[
\boxed{
\sigma_{21}
=
\frac{17}{2}L_{10}+8L_{11}+\frac{15}{2}L_{12}
+2L_{13}+3L_{14}
+\frac{27}{2}(L_6+L_7+L_8+L_9),
}
\]
\[
\boxed{
\sigma_{22}
=
-\frac52L_{10}-\frac92L_{11}-\frac92L_{12}
+4L_{13}-15(L_6+L_7+L_8+L_9),
}
\]
\[
\boxed{
\tau_{20}
=
\frac32(L_6+L_7+L_8+L_9),
}
\]
\[
\boxed{
\tau_{21}
=
\frac12(L_{10}+L_{11}+L_{12}-L_{13})
+3L_6+\frac52L_7+2L_8+\frac32L_9,
}
\]
\[
\boxed{
\tau_{22}
=
2(L_{10}+L_{11}+L_{12}-L_{13})
+\frac{15}{2}L_6+\frac{11}{2}(L_7+L_8+L_9).
}
\]
This is the exact grouped-$P_2$ inverse map on the $3$PN middle block.

\subsection{Exact fit to the solved GR middle block}
\label{app:richer_grouped_p2_compiler_inverse_map:gr_fit}

Now substitute the solved GR middle-block residual
\[
(L_6,\dots,L_{14})=(\Delta l_6,\dots,\Delta l_{14}).
\]
The inverse map then gives one unique grouped coefficient set:
\[
\lambda_{20}^{\rm GR}
=
\frac{\nu\,(2672-3154\nu-276\nu^2+3\pi^2)}{32},
\]
\[
\lambda_{21}^{\rm GR}
=
\frac{\nu\,(-3044+2236\nu+2568\nu^2-3\pi^2)}{192},
\]
\[
\lambda_{22}^{\rm GR}
=
\frac{\nu\,(-30136+32858\nu+7650\nu^2-33\pi^2)}{96},
\]
\[
\sigma_{20}^{\rm GR}
=
\frac{\nu\,(293-308\nu-102\nu^2)}{16},
\]
\[
\sigma_{21}^{\rm GR}
=
\frac{\nu\,(22006-23778\nu-4188\nu^2+21\pi^2)}{192},
\]
\[
\sigma_{22}^{\rm GR}
=
\frac{\nu\,(-2791+2478\nu+3906\nu^2-3\pi^2)}{48},
\]
\[
\tau_{20}^{\rm GR}
=
\frac{3\nu\,(38-87\nu-154\nu^2)}{32},
\]
\[
\tau_{21}^{\rm GR}
=
\frac{\nu\,(4412-5824\nu-6672\nu^2+3\pi^2)}{384},
\]
\[
\tau_{22}^{\rm GR}
=
\frac{\nu\,(3386-3367\nu-2790\nu^2+3\pi^2)}{96}.
\]
With exactly these coefficients, the richer grouped lift reproduces the full solved GR middle block with no residual left:
\[
(l_6,\dots,l_{14})=(\Delta l_6,\dots,\Delta l_{14}).
\]

\subsection{Structural corollaries}
\label{app:richer_grouped_p2_compiler_inverse_map:corollaries}

Two corollaries matter for the carry-forward theorem ledger.

First, every basis element in $(T_A,S_A,V_A)$ carries at least one power of $U$. Therefore the richer grouped compiler still forces
\[
\boxed{l_1=l_2=l_3=l_4=l_5=0}
\]
identically. So the pure-kinetic residual remains outside the grouped-$P_2$ module.

Second, the richer grouped compiler predicts a specific static companion rather than leaving the static slot arbitrary:
\[
\boxed{l_{15}^{(P_2\,\rm pred)}=l_{10}+l_{11}+l_{12}+2(l_6+l_7+l_8+l_9).}
\]
Evaluated on the solved GR middle block, this gives
\[
\boxed{l_{15}^{(P_2\,\rm pred)}=\frac{\nu(293-308\nu-102\nu^2)}{24}.}
\]
This is already close to the exact GR static residual, but it is not equal to it. So after exact grouped middle-block closure, the only remaining grouped-side obstruction is a one-dimensional static remainder, which the main text shows is uniquely geometry-side.

The correct carry-forward reading is therefore:
\[
\boxed{\text{the richer grouped real }P_2\text{ family closes the full nine-slot }3\text{PN middle block exactly,}}
\]
\[
\boxed{\text{but the static completion and the outgoing/isotropy questions remain separate.}}
\]

\section{Static \texorpdfstring{$P_0$}{P 0}/geometry counterterm and sigma-collapse}
\label{app:static_p0_geometry_counterterm_sigma_collapse}

This appendix records the exact scalar-side conclusion that remains after the richer grouped real $P_2$ compiler has closed the full nine-slot $3$PN middle block. The positive grouped result leaves only a one-dimensional static remainder, but the right carrier of that remainder is not obvious if one looks only at the COM chart. At first sight one finds an apparent one-parameter split between a body-local scalar $P_0$ family and a pair/geometry family. The point of this appendix is that the apparent freedom is algebraically fake. Once the exact generic-frame mass identity is used, the whole family collapses to a unique geometry-side completion.

Throughout this appendix we use the generic-frame mass labels
\[
p\equiv m_A,
\qquad
q\equiv m_B,
\qquad
\nu\equiv \frac{pq}{(p+q)^2},
\qquad
U\equiv \frac{G(p+q)}{r}.
\]

\subsection{Residual static slot after grouped closure}
\label{app:static_p0_geometry_counterterm_sigma_collapse:residual_slot}

After exact grouped middle-block closure, the solved COM residual takes the form
\[
\Delta \hat L_3^{\rm GR}
=
\Delta l_1 v^8
+
L_{P_2}^{\rm mid}[\lambda_A^{\rm GR},\sigma_A^{\rm GR},\tau_A^{\rm GR}]
+
\Delta l_{15}^{(0/g)}U^4.
\]
So the grouped real $P_2$ module has done all it can do. The only remaining conservative scalar job is the one-dimensional static remainder
\[
\boxed{
\Delta l_{15}^{(0/g)}
=
\frac{\nu\bigl(408\nu^2+1232\nu-2080+63\pi^2\bigr)}{96}.
}
\]
In the fixed ADM chart this is the full scalar gap left after the grouped-$P_2$ middle block is closed.

\subsection{Why the old constant-coefficient \texorpdfstring{$U$}{U}-block is too small}
\label{app:static_p0_geometry_counterterm_sigma_collapse:U_block_nogo}

Inside the old constant-coefficient generic-frame interaction scaffold, the only static basis polynomial is
\[
U_1=p^2q+pq^2
\]
inside the common prefactor $G^4pq/r^4$. In the COM chart this always reduces to a coefficient proportional to
\[
\nu U^4,
\]
so any such constant-coefficient static slot is linear in $\nu$ after the overall $\nu$ is factored out.

But the exact $3$PN static remainder above contains $\nu^2$ and $\nu^3$ pieces. Therefore the old constant-coefficient $U$-block is too small to carry the actual scalar gap. This is the exact no-go behind the later counterterm interpretation: the remaining scalar completion cannot be represented by just ``one more constant static basis coordinate'' in the original generic-frame scaffold.

\subsection{Apparent body-local \texorpdfstring{$P_0$}{P 0} versus pair/geometry family}
\label{app:static_p0_geometry_counterterm_sigma_collapse:apparent_family}

Once one allows a body-local scalar carrier as well as the pair/geometry carrier, the static completion initially looks like a genuine one-parameter family. The two natural mass polynomials are
\[
U_0=p^3+q^3,
\qquad
U_g=p^2q+pq^2.
\]
The body-local family behaves like a scalar $P_0$ wall/support term, while the pair family is the natural geometry-side carrier because it is intrinsically pair-coupled and vanishes in the strict test-mass branch.

At the generic-frame level one can therefore write the regular one-parameter family
\[
\Delta L_{3,\rm ct}^{(0)}
=
\frac{G^4pq}{r^4}\,\sigma\,\nu\,(p^3+q^3),
\]
\[
\Delta L_{3,\rm ct}^{(g)}
=
\frac{G^4pq}{r^4}
\left[
\frac{408\nu^2+1232\nu-2080+63\pi^2}{96}
-\sigma(1-3\nu)
\right](p^2q+pq^2).
\]
Their COM images combine to reproduce the exact scalar gap for every regular value of the split parameter $\sigma$. So if one stops at the COM bookkeeping stage, the scalar completion appears to admit a free $P_0$-versus-geometry admixture.

A useful warning sign is obtained by asking for a ``pure-$P_0$'' slice. Solving for a split that removes the pair/geometry piece gives
\[
\sigma_{\rm pure\,P_0}(\nu)
=
\frac{-408\nu^2-1232\nu-63\pi^2+2080}{96(3\nu-1)},
\]
which is singular at $\nu=1/3$. So even before the exact collapse is imposed, the apparent pure-$P_0$ option is already unnatural.

\subsection{Exact sigma-collapse and unique geometry completion}
\label{app:static_p0_geometry_counterterm_sigma_collapse:sigma_collapse}

The apparent $\sigma$-freedom is not physical. The exact mass identity
\[
\boxed{
\nu(p^3+q^3)=(1-3\nu)(p^2q+pq^2)
}
\]
collapses the whole family identically. Indeed,
\[
\sigma\,\nu(p^3+q^3)=\sigma(1-3\nu)(p^2q+pq^2),
\]
so the $\sigma$-dependence cancels between the body-local and pair pieces. The full static counterterm therefore reduces to the unique geometry-side expression
\[
\boxed{
\Delta L_{3,\rm ct}^{\rm scalar}
=
\frac{G^4pq}{96r^4}
\bigl(408\nu^2+1232\nu-2080+63\pi^2\bigr)
(p^2q+pq^2).
}
\]
There is no independent regular body-local $P_0$ completion left in the fixed ADM chart.

The same conclusion can be read directly on the target coefficients. The imported generic-frame static target coefficient is
\[
c_U^{\rm target}=-\frac{227}{24}+\frac{21\pi^2}{32},
\]
while the richer grouped-$P_2$ middle-block fit already forces the static companion
\[
c_U^{(P_2\,\rm pred)}=\frac{293-308\nu-102\nu^2}{24}.
\]
Their exact difference is the unique geometry coefficient
\[
\boxed{
c_U^{(g)}=
\frac{408\nu^2+1232\nu-2080+63\pi^2}{96}.
}
\]
On the COM side this reduces exactly to
\[
\boxed{
\Delta l_{15}^{(g)}=
\frac{\nu\bigl(408\nu^2+1232\nu-2080+63\pi^2\bigr)}{96}.
}
\]
So the sigma-collapse is not a matter of taste or chart choice. It is the exact statement that, after the grouped-$P_2$ middle block is closed, the remaining $3$PN static scalar piece lands uniquely in the pair/geometry lane.

\subsection{Push through the generic-frame Hamiltonian compiler}
\label{app:static_p0_geometry_counterterm_sigma_collapse:hamiltonian_image}

Once the lower-order ledger is frozen, the generic-frame $3$PN Hamiltonian compiler remains the exact sign flip
\[
\Delta H_3=-\Delta L_3.
\]
So the unique scalar counterterm lifts immediately to
\[
\boxed{
\Delta H_{3,\rm ct}^{\rm scalar}=-\Delta L_{3,\rm ct}^{\rm scalar}.
}
\]
There is no further algebraic freedom to resolve. The carry-forward interpretation is therefore sharp:
\[
 \boxed{\begin{gathered}
 \text{after the grouped-}P_2\text{ middle block is closed,}\\
 \text{the remaining static }3\text{PN scalar piece is uniquely geometry-side}
 \end{gathered}}
\]

This is the scalar-side analog of the middle-block grouped-$P_2$ result. The algebraic bottleneck is gone; what remains is the physical question of what moving-throat wall/geometry dynamics microscopically realizes this geometry-side scalar completion.

\section{Pure-kinetic collapse derivation}
\label{app:pure_kinetic_collapse_derivation}

This appendix records the final kinetic audit behind the conservative $3$PN theorem. After the richer grouped real $P_2$ compiler has closed the full nine-slot middle block and the static scalar remainder has collapsed uniquely to the geometry side, one apparently independent center-of-mass kinetic slot is still left over. At first sight that leftover looks like one more independent $3$PN throat-response datum. The exact free-body audit shows that this reading is wrong. The residual is not a new comparable-mass interaction family at all. It is only the ordinary-Lagrangian counterimage of the universal free relativistic two-body center-of-mass Hamiltonian.

The point of this appendix is therefore very narrow: to show explicitly that the apparent pure-kinetic remainder
\[
\Delta l_1=\frac{3\nu(3\nu-1)(4\nu-1)}{16}
\]
is generated automatically by the exact COM compiler map, and that no additional generic-frame pure-kinetic interaction module should be introduced in the fixed ADM chart.

\subsection{Apparent residual after grouped and static closure}
\label{app:pure_kinetic_collapse_derivation:apparent_residual}

After exact grouped middle-block closure and unique static completion, the remaining COM kinetic slots are
\[
\boxed{
\Delta l_1=\frac{3\nu(3\nu-1)(4\nu-1)}{16},
\qquad
\Delta l_2=\Delta l_3=\Delta l_4=\Delta l_5=0.
}
\]
In the isotropic $15$-slot COM basis, this is the whole apparent pure-kinetic residue. Taken by itself, it looks as though one more independent $3$PN kinetic throat coefficient must still be supplied.

That interpretation would only be correct if a genuinely new kinetic datum first appeared as a new generic-frame interaction family. The exact free-body audit shows that no such family exists. The pure-kinetic issue is therefore not a new interaction-sector problem. It is a compiler problem.

\subsection{Exact free generic-frame block and its naive COM reduction}
\label{app:pure_kinetic_collapse_derivation:generic_frame_block}

At generic-frame level, the exact free relativistic two-body ordinary Lagrangian is simply the sum of the two separate free-body square roots,
\[
L_{\rm free}^{\rm gen}
=
-m_Ac^2\sqrt{1-\frac{v_A^2}{c^2}}
-
 m_Bc^2\sqrt{1-\frac{v_B^2}{c^2}}.
\]
Expanding through $3$PN gives
\[
\begin{aligned}
L_{\rm free}^{\rm gen}
=
&-(m_A+m_B)c^2
+\frac12\bigl(m_Av_A^2+m_Bv_B^2\bigr)
+\frac{1}{8c^2}\bigl(m_Av_A^4+m_Bv_B^4\bigr)
\\
&+\frac{1}{16c^4}\bigl(m_Av_A^6+m_Bv_B^6\bigr)
+\frac{5}{128c^6}\bigl(m_Av_A^8+m_Bv_B^8\bigr)
+O(c^{-8}).
\end{aligned}
\]
So the exact generic-frame $3$PN pure-kinetic block is
\[
\boxed{
L_{3,\rm free}^{\rm gen}
=
\frac{5}{128}\bigl(m_Av_A^8+m_Bv_B^8\bigr).
}
\]
There is no mixed comparable-mass pure-kinetic interaction monomial here. In particular, the free generic-frame block contains no independent $m_Am_B$-type octic interaction that could serve as a new $3$PN constitutive datum.

Now impose the naive COM substitution
\[
\mathbf v_A=\eta_B\mathbf v,
\qquad
\mathbf v_B=-\eta_A\mathbf v,
\qquad
\eta_A\equiv\frac{m_A}{M},
\qquad
\eta_B\equiv\frac{m_B}{M},
\qquad
M\equiv m_A+m_B,
\]
with reduced mass
\[
\mu\equiv\frac{m_Am_B}{M}
\]
and symmetric mass ratio
\[
\nu\equiv\frac{\mu}{M}=\frac{m_Am_B}{M^2}.
\]
Then the carried COM seed coefficient is
\[
\frac{1}{\mu}\,\frac{5}{128}\bigl(m_A\eta_B^8+m_B\eta_A^8\bigr)v^8
=
 l_1^{\rm seed}v^8,
\]
with
\[
\boxed{
 l_1^{\rm seed}
 =
 \frac{5}{128}-\frac{35}{128}\nu+\frac{35}{64}\nu^2-\frac{35}{128}\nu^3.
}
\]
So the earlier seed coefficient was nothing more than the naive COM reduction of the universal free generic-frame relativistic Lagrangian.

\subsection{Exact free COM Hamiltonian target}
\label{app:pure_kinetic_collapse_derivation:free_com_target}

The correct free reference object is instead the exact relativistic two-body COM Hamiltonian,
\[
H_{\rm free}^{\rm COM}
=
\sqrt{m_A^2c^4+P^2c^2}
+
\sqrt{m_B^2c^4+P^2c^2},
\]
with reduced momentum variable defined by
\[
P=\mu p.
\]
Subtracting the rest energy and dividing by $\mu$ gives the reduced expansion
\[
\frac{H_{\rm free}^{\rm COM}-Mc^2}{\mu}
=
\frac12p^2
-
\frac{1-3\nu}{8c^2}p^4
+
\frac{1-5\nu+5\nu^2}{16c^4}p^6
+
\frac{h_1^{\rm free}}{c^6}p^8
+O(c^{-8}),
\]
where the exact $3$PN free coefficient is
\[
\boxed{
 h_1^{\rm free}
 =
 \frac{-5+35\nu-70\nu^2+35\nu^3}{128},
 \qquad
 h_2=h_3=h_4=h_5=0.
}
\]
This is not a new observable of the throat. It is just the universal free relativistic two-body Hamiltonian written in reduced COM variables. So the relevant comparison is not between the ordinary seed and a new kinetic throat datum. It is between the ordinary seed and the already-fixed free COM Hamiltonian target.

\subsection{Exact COM compiler map and collapse of the apparent residual}
\label{app:pure_kinetic_collapse_derivation:compiler_collapse}

The exact COM compiler acts on the first pure-kinetic slot by
\[
\boxed{
 l_1=F_1(\nu)-h_1,
 \qquad
 F_1(\nu)=\frac{3}{16}\nu-\frac{21}{16}\nu^2+\frac94\nu^3.
}
\]
Applying this map to the naive free ordinary seed gives its Hamiltonian image,
\[
\boxed{
 h_1^{\rm seed}
 =
 F_1(\nu)-l_1^{\rm seed}
 =
 -\frac{5}{128}+\frac{59}{128}\nu-\frac{119}{64}\nu^2+\frac{323}{128}\nu^3.
}
\]
So the shift between the naive seed and the true free COM target is
\[
\begin{aligned}
\Delta l_1
&=
 l_1^{\rm full}-l_1^{\rm seed}
=
 h_1^{\rm seed}-h_1^{\rm free}
\\
&=
\frac{3}{16}\nu-\frac{21}{16}\nu^2+\frac94\nu^3
=
\boxed{\frac{3\nu(3\nu-1)(4\nu-1)}{16}}.
\end{aligned}
\]
Equivalently, the full ordinary coefficient required by the exact free COM Hamiltonian is
\[
\boxed{
 l_1^{\rm full}
 =
 F_1(\nu)-h_1^{\rm free}
 =
 \frac{5}{128}-\frac{11}{128}\nu-\frac{49}{64}\nu^2+\frac{253}{128}\nu^3.
}
\]
Thus the whole apparent leftover is exactly the ordinary-Lagrangian compensation required to reproduce the already-fixed free relativistic COM Hamiltonian.

The essential consequence is immediate:
\[
\boxed{
\text{the apparent COM pure-kinetic remainder is not a new dynamical throat datum.}
}
\]
It is only the ordinary-side image of the universal free Hamiltonian under the exact COM compiler.

\subsection{Fixed-chart interpretation}
\label{app:pure_kinetic_collapse_derivation:fixed_chart_interpretation}

The theorem-level reading is therefore sharp.

First, the exact generic-frame free-body ordinary block has no mixed pure-kinetic interaction monomial:
\[
L_{3,\rm free}^{\rm gen}=\frac{5}{128}(m_Av_A^8+m_Bv_B^8).
\]
Second, naive COM reduction of that block gives the carried seed coefficient $l_1^{\rm seed}$ and nothing more. Third, the exact free relativistic two-body COM Hamiltonian fixes the true target coefficient $h_1^{\rm free}$ with
\[
 h_2=h_3=h_4=h_5=0.
\]
So once the exact COM compiler is imposed, the residual
\[
\Delta l_1=\frac{3\nu(3\nu-1)(4\nu-1)}{16}
\]
is generated automatically, and no extra generic-frame pure-kinetic interaction module should be added in the fixed ADM chart.

This is the final kinetic-side analogue of the grouped-$P_2$ and sigma-collapse results. The algebraic bottleneck is gone. After grouped middle-block closure, unique static geometry completion, and pure-kinetic collapse, the conservative $3$PN split really does stop with three lanes only: the compiler image of free kinematics, the exact grouped real $P_2$ module, and the unique geometry-side static completion.

\section{Consolidated theorem ledger}
\label{app:theorem-ledger}

This appendix is the shortest single place to recover the final conservative $3$PN status after the earlier appendices are assembled. Its purpose is not to re-derive the chain once more, but to freeze the exact theorem package in one compact location so that later work does not reopen already settled bookkeeping.

The right reading is the same as in the earlier full $1$PN and $2$PN papers: this is an exact conservative theorem package \emph{within the declared closure hierarchy}, not an assumption-free theorem of the fully solved moving-throat PDE. Within that hierarchy, however, the algebraic $3$PN bottleneck is gone.

\subsection{Exact constants and exact compiler data}
\label{app:theorem-ledger-constants}

The strict one-body Schwarzschild gate fixes the minimal $3$PN repair ledger to
\[
\boxed{
\mu_{\rho,3}=\frac14,
\qquad
 d_3=-\frac{45}{4},
\qquad
 s_{24}=-\frac1{16}.
}
\]
These are the exact carried constants that replace the earlier possibility of a one-parameter continuation of the $2$PN self sector.

Once the lower-order ledger is frozen and the fixed ADM Hamiltonian chart is chosen, the full generic-frame $3$PN interaction compiler becomes
\[
\boxed{
\Delta H_3=-\Delta L_3,
\qquad
C_{\rm gen}=-I_{50}.
}
\]
So the earlier COM-blind and contact/gauge bookkeeping freedoms do not survive as true algebraic ambiguity in the chosen chart.

For the grouped real $P_2$ middle block, the richer local nine-family compiler is exactly invertible, with
\[
\boxed{
\det M_{\rm mid}=-\frac4{27}.
}
\]
This is the decisive exact statement behind the successful middle-block closure: the grouped conservative fit is neither underdetermined nor merely gauge-selected.

After the grouped closure is imposed, the remaining static scalar piece is uniquely geometry-side. In COM form the exact coefficient is
\[
\boxed{
\Delta l_{15}^{(g)}=
\frac{\nu\bigl(408\nu^2+1232\nu-2080+63\pi^2\bigr)}{96},
}
\]
and in generic-frame coefficient language
\[
\boxed{
 c_U^{(g)}=
 \frac{408\nu^2+1232\nu-2080+63\pi^2}{96}.
}
\]
Likewise, the apparent pure-kinetic remainder is
\[
\boxed{
\Delta l_1=
\frac{3\nu(3\nu-1)(4\nu-1)}{16}.
}
\]
The pure-kinetic collapse appendix shows that this is not a new throat datum, but the ordinary-Lagrangian counterimage of the universal free relativistic two-body COM Hamiltonian.

\subsection{Exact split formulas}
\label{app:theorem-ledger-split}

The exact isotropic one-body $3$PN target remains
\[
\boxed{
\frac{L_{3\mathrm{PN},1\mathrm{body}}}{m}
=
\frac5{128}v^8
+\frac{11}{16}Uv^6
+\frac{47}{16}U^2v^4
+\frac{13}{8}U^3v^2
-\frac18U^4.
}
\]
This is the immutable test-mass gate beneath the full comparable-mass lift.

Beyond the frozen one-body/self-static seed, the exact comparable-mass COM residual is fully assigned by the three-lane split
\[
\boxed{
\begin{aligned}
\Delta \hat L_3^{\rm GR}
=\;&\underbrace{\Delta l_1\,v^8}_{\text{compiler image of the free COM Hamiltonian}}
+\underbrace{L_{P_2}^{\rm mid}}_{\text{exact grouped real }P_2\text{ middle-block closure}}
\\
&+\underbrace{\Delta l_{15}^{(g)}U^4}_{\text{unique geometry completion}}.
\end{aligned}
}
\]
So after the earlier appendices are assembled, the conservative $3$PN residual really does stop with exactly three lanes and no fourth algebraic family.

The middle block itself is carried by the richer local grouped family
\[
(T_{20},T_{21},T_{22},S_{20},S_{21},S_{22},V_{20},V_{21},V_{22}),
\qquad
A\in\{20,21,22\},
\]
with the useful identifications
\[
S_A=U^2C_A^2,
\qquad
V_A=\frac{v^2}{U}S_A,
\]
and the exact inverse map guaranteed by the richer grouped-$P_2$ compiler appendix. The exact grouped fit fixes the nine middle slots and leaves only the static scalar slot to be completed.

Writing
\[
p\equiv m_A,
\qquad
q\equiv m_B,
\]
the full static completion is uniquely
\[
\boxed{
\Delta L_{3,\rm ct}^{\rm scalar}
=
\frac{G^4pq}{96r^4}
\bigl(408\nu^2+1232\nu-2080+63\pi^2\bigr)
(p^2q+pq^2),
}
\]
with fixed-chart Hamiltonian image
\[
\boxed{
\Delta H_{3,\rm ct}^{\rm scalar}=-\Delta L_{3,\rm ct}^{\rm scalar}.
}
\]
This is the exact sigma-collapsed pair/geometry remainder from the static $P_0$/geometry counterterm and sigma-collapse appendix.

The pure-kinetic sector likewise needs no new interaction module. The exact free generic-frame block is
\[
\boxed{
L_{3,\rm free}^{\rm gen}=
\frac{5}{128}\bigl(m_Av_A^8+m_Bv_B^8\bigr),
}
\]
while the exact free COM Hamiltonian target has
\[
\boxed{
 h_1^{\rm free}=
 \frac{-5+35\nu-70\nu^2+35\nu^3}{128},
 \qquad
 h_2=h_3=h_4=h_5=0.
}
\]
The exact COM compiler then produces $\Delta l_1$ automatically, so no independent generic-frame pure-kinetic interaction family should be added in the fixed ADM chart.

\subsection{Fixed versus open data}
\label{app:theorem-ledger-fixed-open}

Within the declared conservative hierarchy, the following should now be treated as fixed.

\begin{enumerate}
    \item The one-body $3$PN repair ledger
    \[
    (\mu_{\rho,3},d_3,s_{24})=
    \left(\frac14,-\frac{45}{4},-\frac1{16}\right).
    \]
    \item The exact generic-frame uniqueness theorem in the fixed ADM chart,
    \[
    \Delta H_3=-\Delta L_3,
    \qquad
    C_{\rm gen}=-I_{50}.
    \]
    \item The exact grouped real $P_2$ middle-block closure, with invertible richer compiler and
    \[
    \det M_{\rm mid}=-\frac4{27}.
    \]
    \item The unique geometry-side completion of the remaining static slot.
    \item The pure-kinetic collapse to the universal free relativistic COM Hamiltonian.
\end{enumerate}

In that precise sense, the conservative $3$PN two-body ledger is now fully assigned in the fixed ADM chart. What remains open is no longer algebraic coefficient matching.

The open problems are now dynamical and constitutive. First, the richer grouped real $P_2$ constitutive map should be derived directly from the moving-throat problem rather than by inverse matching. Second, the unique geometry completion should be derived directly from the throat-side geometry/compressibility dynamics. Third, those two conservative lanes must be carried into the already narrowed $2.5$PN bridge, where the live universal route has already been reduced to the grouped real quadrupole bundle and its passive/outgoing normalization problem.

The final carry-forward statement is therefore
\[
\boxed{
\begin{gathered}
\text{the conservative $3$PN algebra is closed;}\\
\text{the remaining bottleneck is the moving-throat origin of the grouped-}P_2\text{ and geometry lanes.}
\end{gathered}
}
\]

\end{document}
